\chapter[The Deligne series III: fusion modules and Ravinder's kernel]{The Deligne series III: fusion modules, Demazure modules and Ravinder's kernel}\label{ch:L26}
\begingroup\setlength{\emergencystretch}{3em}\sloppy

\providecommand{\ltwentysixg}{\mathfrak g}
\providecommand{\ltwentysixgt}{\mathfrak g[t]}
\providecommand{\ltwentysixWloc}{W_{\mathrm{loc}}}
\providecommand{\ltwentysixD}{\mathcal D}
\providecommand{\ltwentysixF}{\mathcal F}
\providecommand{\ltwentysixU}{U}
\providecommand{\ltwentysixcyc}{\operatorname{cyc}}
\providecommand{\ltwentysixHilb}{\operatorname{Hilb}}
\providecommand{\ltwentysixqb}[2]{\genfrac{[}{]}{0pt}{}{#1}{#2}_q}

This chapter was written against an earlier state of the LeanVOA library (28 October 2026).
For this edition, \texttt{scripts/textbook/sync.py} checks that every Lean quotation below occurs verbatim (up to whitespace and comments) in the source files of that release. What was machine-checked for this chapter,
and what is taken from the project's build records, is stated once in Section~\ref{l26:sec:statusconv} and in detail in the
verification record (Section~\ref{l26:sec:verification}).

The chapter covers the Lean files that the book assigns to it (\texttt{Notes/Textbook/assignment.txt}: the rules for
\lean{DeligneFusion*}, \lean{LatticeRealization*}, \lean{*FusionKernel*}, \lean{D4FusionF}, \lean{F4LieKernel*} and
\lean{G2Fusion}). The files that carry the mathematics of the proved layer are presented declaration by declaration; certificate
data and certificate checks (62 row-check files for the lattice realizations alone) are described collectively in
Section~\ref{l26:sec:certs}. The Chari--Venkatesh files that prove the $A_2$ total bound (\lean{DeligneFusionCV*}) are
presented through their headline statements only (Section~\ref{l26:sec:cv}).

\paragraph{Prerequisites.} The reader should know the following, which earlier chapters develop.
\begin{itemize}
\item The Deligne series, its levels $k=-1-h^\vee/6$, and the chain of coefficientwise inequalities that proves
  classical freeness (Chapter~\ref{ch:L24}, Section~\ref{l24:sec:chain-paper}). In particular the inputs
  \textbf{U5} (the fusion module bound) and \textbf{U6} (the mixed fusion recursion), which Chapter~\ref{ch:L24}
  leaves as named hypotheses, are the subject of this chapter.
\item Current modules \lean{CurrentModule}, the cyclic degree pieces \lean{cycDeg}, the $\mathfrak{sl}_3$ local Weyl
  relations \lean{LocalWeylRel} and input U5 in its $A_2$ form \lean{SectorWeyl.FusionModuleBound}
  (Definition~\ref{l24:def:U5} and Theorem~\ref{l24:thm:seven}); the type-generic versions
  \lean{LocalWeylRelG} and \lean{FusionModuleBoundG}, and the sector data \lean{DeligneSectorData}
  (Definitions~\ref{l24:def:U5} and \ref{l24:def:sectordata}); the recursion \lean{mixed} and the structure
  \lean{MixedFusionRecursion} (Definition~\ref{l24:def:chaininputs}).
\item The level-one affine Demazure characters and their $q=1$ totals (Chapter~\ref{ch:L25},
  Section~\ref{l25:sec:demazure}, Definition~\ref{l25:def:demazure} and Theorem~\ref{l25:thm:total}); for $A_2$ the
  Gelfand--Tsetlin polynomial \lean{demazureHilb} (Chapter~\ref{ch:L25}).
\item Lattice vertex algebras $V_L$, their charge and conformal-weight gradings, and the vertex-operator vanishing
  $e^\alpha_{(n)}e^\beta=0$ for $n\ge-\langle\alpha,\beta\rangle$ (Chapter~\ref{ch:L03} and the lattice chapters).
\end{itemize}

\paragraph{What this chapter is about.} The \emph{seven types} of this chapter are the members
$A_2$, $D_4$, $E_6$, $E_7$, $E_8$, $G_2$, $F_4$ of the Deligne series. The eighth member, $A_1$, is treated separately
(Chapter~\ref{ch:L24}, Section~\ref{l24:sec:math}: its relation module has a different shape, and it is the $m=1$ case of
Chapter~\ref{ch:L20}). The chain of Chapter~\ref{ch:L24} bounds the operator count of $L_k(\ltwentysixg)$ by a sum over
$m$ of ``$m$-particle sectors''. Each sector contributes the graded character $\ltwentysixF_m(q)$ of a \emph{fusion module}
$F_m$, defined in Section~\ref{l26:sec:fusion} by generators and relations. The chain needs two facts about these modules.
\begin{itemize}
\item \textbf{U5}: every module satisfying the relations of $F_m$ has graded dimension at most $\ltwentysixF_m(q)$ and total
  dimension at most $(\dim\ltwentysixg)^m$.
\item \textbf{U6}: the $\ltwentysixF_m$ are tied to the level-one Demazure characters $\ltwentysixD_n(q)$ by Ravinder's mixed
  recursion \eqref{l26:eq:mixed}.
\end{itemize}
The library proves U6 by \emph{defining} $\ltwentysixF_m$ as the unique solution of the recursion. It reduces U5, by
Ravinder's squeeze, to three statements.
\begin{enumerate}
\item[(i)] \textbf{Ravinder's kernel lemma.} Its Lean analogue is \emph{proved} for all seven types. Ravinder's paper
  states it for every simple $\ltwentysixg$ (only his Theorem~2 and Corollary~3 assume simply laced); for $G_2$, $F_4$ the
  library proves the statement for its own family of relations, which imposes the same relations as Ravinder's
  (Remark~\ref{l26:rem:fusionset}).
\item[(ii)] \textbf{U7}, the existence of a model of the Demazure module with the right graded dimension: for each $n$, some
  current module with a vector satisfying the relations of $M_{n,0}$ (Section~\ref{l26:sec:ravinder}) whose degree pieces
  are at least as large as those of $D(1,n\theta)$ (Lean: \lean{DemazureRealization n}, Definition~\ref{l26:def:squeezeinputs}).
  U7 is \emph{proved} for $A_2$ and for $D_4$ (\lean{D4Char.d4_demazureRealization}, Chapter~\ref{ch:L27}). For $E_6$, $E_7$,
  $E_8$ it is reduced to a named dimension statement on lattice vertex algebras. For $G_2$, $F_4$ it is a named hypothesis
  (there is no lattice route). For every type U7 is a published theorem (Ravinder's Proposition~1 with the Demazure character
  formula); for the five remaining types what is missing is its formalization.
\item[(iii)] \textbf{The total-dimension bound} $\dim F_m\le(\dim\ltwentysixg)^m$. Its finiteness half is \emph{proved} for
  all seven types. Its dimension half is \emph{proved for $A_2$} (\lean{CVUniv.fusionTotalDim}, Section~\ref{l26:sec:cv}) and
  remains a named hypothesis (\lean{FusionTotalDimG}) for the other six types.
\end{enumerate}
Feigin's proof of the total bound (\emph{SIGMA} 4 (2008) 070, arXiv:0806.4851, Cor.~2) has a gap that we document; a repair
(Kus--Littelmann, Thm.~4.2, valid for every simple $\ltwentysixg$, with one proof step flagged by the project) and an independent
$\mathfrak{sl}_3$ route are published, and the library's $A_2$ proof follows the independent route (Setia--Rani--Khandai). So for $A_2$ every input of
the chain in this chapter is a theorem, and Chapter~\ref{ch:L27} states the resulting unconditional classical freeness of
$L_{-3/2}(\mathfrak{sl}_3)$.

\paragraph{The findings of this chapter, in brief.}
\begin{enumerate}
\item \textbf{U6 is a theorem by construction} (\lean{mixedFusionRecursion_fusionOf}): the mixed recursion is
  unitriangular, so for any Demazure series it has exactly one solution \lean{fusionOf}. The price is that the
  identification ``\lean{fusionOf} $=$ the true fusion character'' becomes a \emph{consequence} of U5, not an input.
\item \textbf{Ravinder's squeeze is formalized} for every type (\lean{squeeze}, \lean{squeezeG}): nonnegativity, the
  spanning half of Ravinder's exact sequences, a Demazure lower bound and a total upper bound force
  $\mathrm{ch}\,M_{0,m}=\ltwentysixF_m$ coefficientwise.
\item \textbf{The Lean analogue of Ravinder's kernel lemma is unconditional for all seven types.} For $A_2$ it is
  stated in the $A_2$-specific form \lean{RavinderKernelRel} (through \lean{MabRel}) and follows from Garland's
  identities (\lean{ravinderKernelRel}). For the other six types it is \lean{RavinderKernelRelG} for the library's family
  \lean{ravFam}, and follows from a decidable table of structure-constant facts and the
  integrability statement \lean{ThetaNilG}, which is itself proved per type by Heisenberg chains or by an
  $\operatorname{ad}$-nilpotency lemma (\lean{ravinderKernelRelG_of_table}, \lean{d4_thetaNil},
  \lean{g2_thetaNil}, \dots). For $G_2$, $F_4$ the Lean statement is about the library's own family, which imposes the same relations
  as Ravinder's (Remark~\ref{l26:rem:fusionset}); Ravinder's Theorem~1, Corollaries~1--2 and Proposition~1 are stated
  for every simple $\ltwentysixg$.
\item \textbf{The finiteness clause of the total bound is proved for all seven types} (\lean{cycDeg_finiteDimensional},
  \lean{fusionTotalBoundG_of_dim}), by Kac's lemma and a kernel-checked nilpotency certificate.
\item \textbf{For $A_2$, U7 is unconditional} (\lean{EvalModel.demazureRealization_all}), by a fusion product of
  $2N$ evaluation modules at distinct points, combined with the Chari--Loktev spanning theorem for $\mathfrak{sl}_3$
  (\lean{clSpanning}), which is formalized in full, and the Chari--Pressley $\mathfrak{sl}_2$ box theorem (\lean{box_span}).
\item \textbf{The $A_2$ total bound is a theorem}, \lean{CVUniv.fusionTotalDim} ($\dim F_m\le 8^m$ for every module with the
  relations; Theorem~\ref{l26:thm:cvtotal}), proved from two base families of Chari--Venkatesh modules that are themselves proved.
  The earlier reductions to \lean{RavinderKernelLower} and \lean{IdealLower} (each implies \lean{FusionTotalDim}) remain in the
  library as alternative routes; those two \lean{Prop}s are not proved and are no longer needed.
\item \textbf{Lattice realizations} (the relations half of U7) are proved for $D_4$, $E_6$, $E_7$, $E_8$
  (\lean{LatData.sat}, \lean{e8_lattice_sat}); the dimension half is the named \lean{D4LatticeDim}, \dots,
  \lean{E8LatticeDim}, numerically sharp where tested.
\item \textbf{The literature behind the $A_2$ bound.} Feigin's Cor.~2, which Ravinder cites, rests on a sketched Prop.~1
  whose natural reading fails; a repair is published (Kus--Littelmann 2015, Thm.~4.2), and an independent upper-bound route exists
  (Setia--Rani--Khandai, arXiv:2405.11944, \S8), which is the one formalized. Section~\ref{l26:sec:open} records the paper
  analysis, including a \emph{slice lemma} in mixed evaluation models that would give an alternative proof; it is not proved and
  not needed. For the other six types the total bound is open in Lean but published (Kus--Littelmann, with the flagged step,
  which for $D_4$ at $m=2$ is insufficient as written although the theorem is not known to fail; Chapter~\ref{ch:L27}).
\end{enumerate}

\paragraph{The logical arc and its ledger.} Table~\ref{l26:tab:ledger} is the map of the chapter: for each input of the
chain that concerns fusion modules, where it is stated, and whether it is proved, reduced, or open.
\begin{table}[ht]\centering\small
\caption{The fusion inputs of the Deligne chain and their status in Lean.}\label{l26:tab:ledger}
\begin{tabular}{p{0.20\textwidth}p{0.40\textwidth}p{0.30\textwidth}}
\toprule
input & Lean & status \\
\midrule
U6 & \lean{fusionOf}, \lean{mixedFusionRecursion_fusionOf} & proved (by construction) \\
squeeze & \lean{squeeze}, \lean{squeezeG}, \lean{fusionModuleBoundG_of_kernel} & proved \\
universal modules & \lean{univRep}, \lean{univRepG}, \lean{finrank_cycDeg_le_univG} & proved \\
exact-sequence step & \lean{finrank_step}, \lean{finrank_stepG} & proved \\
kernel lemma & \lean{ravinderKernelRel} ($A_2$, \lean{MabRel} form); \lean{ravinderKernelRelG_of_table} $+$ \lean{X_thetaNil} & proved, all types (Lean analogue; $G_2$, $F_4$: library family) \\
finiteness clause & \lean{cycDeg_finiteDimensional}; \lean{X_fusionTotalBoundG_of_dim} & proved, all types \\
total bound ($A_2$) & \lean{CVUniv.fusionTotalDim} (of \lean{FusionTotalDim}) & proved \\
total bound (others) & \lean{FusionTotalDimG} & \textbf{named} \\
U7 ($A_2$) & \lean{EvalModel.demazureRealization_all} & proved \\
U7 ($D_4$) & \lean{D4Char.d4_demazureRealization} & proved (Chapter~\ref{ch:L27}) \\
U7 (other types) & \lean{DemazureRealizationG}; lattice \lean{LatData.sat} $+$ \lean{XLatticeDim} & \textbf{named} (relations half proved for $E_{6,7,8}$; no lattice route for $G_2$, $F_4$) \\
$A_2$ alternative routes & \lean{RavinderKernelLower}, \lean{IdealLower} & named, unproved, not needed; each implies \lean{FusionTotalDim} \\
slice lemma & none (paper only) & open; not needed \\
\bottomrule
\end{tabular}
\end{table}

\paragraph{How the $A_2$ argument fits together.} Table~\ref{l26:tab:dag} shows which theorem feeds which, and what
hypothesis is passed on. Reading it top to bottom: the kernel lemma and finiteness make the universal characters satisfy the
spanning inequality (R); U7 gives the Demazure lower bound (D); the total bound gives (T); the squeeze turns (R), (D), (T)
into U5; and U5 with U6 gives classical freeness. \lean{FusionTotalDim} enters \emph{only} through (T), and it is proved
(Theorem~\ref{l26:thm:cvtotal}).
\begin{table}[ht]\centering\small
\caption{Dependencies of the $A_2$ argument.}\label{l26:tab:dag}
\begin{tabular}{p{0.30\textwidth}p{0.30\textwidth}p{0.30\textwidth}}
\toprule
result & feeds & open hypothesis it passes on \\
\midrule
Garland (Thm.~\ref{l26:thm:garland}) & kernel lemma; box spanning & none \\
kernel lemma (Thm.~\ref{l26:thm:kernelA2}) & exact-sequence step & none \\
finiteness (Thm.~\ref{l26:thm:finite}) & finiteness of universal pieces & none \\
exact-sequence step (Thm.~\ref{l26:thm:step}) & (R) in the assembly & none \\
box spanning (Thm.~\ref{l26:thm:box}) & Chari--Loktev & none \\
Chari--Loktev (Thm.~\ref{l26:thm:cl}) & U7 & none \\
U7 (Thm.~\ref{l26:thm:u7}) & (D) in the assembly & none \\
\lean{FusionTotalDim} (Thm.~\ref{l26:thm:cvtotal}) & (T) in the assembly & none (proved) \\
squeeze (Thm.~\ref{l26:thm:squeeze}) & U5 for $\mathrm{fusionF}$ & the squeeze inputs \\
assembly (Thm.~\ref{l26:thm:a2assembly}) & classical freeness & \lean{FusionTotalDim} (as stated) \\
final forms (Thm.~\ref{l26:thm:u7final}) & Chapter~\ref{ch:L27} & none, once Thm.~\ref{l26:thm:cvtotal} is applied \\
\bottomrule
\end{tabular}
\end{table}

\paragraph{The named statements and how they are related.} Table~\ref{l26:tab:named} collects every named \lean{Prop} of
this chapter that is \emph{not} proved. For each type other than $A_2$ there is one open core, the total bound
\lean{FusionTotalDimG}, together with U7; for $A_2$ the total bound is proved, and the two $A_2$ rows are alternative routes that
are no longer needed.
\begin{table}[ht]\centering\small
\caption{The named, unproved statements of this chapter. ``$\Rightarrow$'' means a Lean theorem.}\label{l26:tab:named}
\begin{tabular}{p{0.22\textwidth}p{0.36\textwidth}p{0.32\textwidth}}
\toprule
statement & relation proved in Lean & what it would take \\
\midrule
\lean{RavinderKernelLower} & $\Rightarrow$ \lean{FusionTotalDim} (Thm.~\ref{l26:thm:a2red}) & a large kernel in one model (slice lemma) \\
\lean{IdealLower} & $\Rightarrow$ \lean{FusionTotalDim} (Thm.~\ref{l26:thm:a2red}) & a large ideal image in one model \\
\lean{FusionTotalDimG} ($G_2,D_4,F_4,E_{6,7,8}$) & for $D_4$: implied by a rank bound on a presented module $\operatorname{PTE}_m$ for all $m$ (proved for $m\le2$), and by $\mathfrak{sl}_4$ core bounds with at least two adjoint sites (\lean{D4LevelReduction.fusionTotalDim_of_coreBounds}, with the one-adjoint-site bounds proved) (Chapter~\ref{ch:L27}, Section~\ref{l27:sec:feigin}); for $D_4$ this is the only remaining input & a formalization of Kus--Littelmann Thm.~4.2 (published; one step flagged, insufficient as written for $D_4$ at $m=2$) \\
\lean{DemazureRealizationG} (non-$A_2$) & implied by \lean{XLatticeDim} for $D_4$, $E_{6,7,8}$ (Thm.~\ref{l26:thm:latsat}); for $D_4$ proved for all $n$ (\lean{D4Char.d4_demazureRealization}; aggregate build green, 3 November 2026; Chapter~\ref{ch:L27}) & a formalization of the Demazure character formula (published: Ravinder Prop.~1 $+$ Kumar/Mathieu), or a $D/E$ Chari--Loktev \\
\lean{D4LatticeDim}, \dots, \lean{E8LatticeDim} & $\Rightarrow$ U7 for that type & a spanning bound in $V_L$ \\
\lean{CLIndependence} & none; not used & --- \\
\bottomrule
\end{tabular}
\end{table}

%======================================================================
\section{Conventions, notation, and reading the Lean}\label{l26:sec:conventions}
%======================================================================

\subsection{Machine status, stated once}\label{l26:sec:statusconv}

\textbf{Axioms.} Every status box reports the \texttt{\#print axioms} result for the theorems it discusses. Unless the box
says otherwise, the result is exactly Lean's three standard axioms \lean{propext}, \lean{Classical.choice}, \lean{Quot.sound}
(``standard''). \lean{VOA.pbw} is a theorem of the library (Chapter~\ref{ch:L02}); no project axiom is used in this chapter.

\textbf{Re-run or record.} ``(re-run)'' means that we ran \texttt{\#print axioms} ourselves for this chapter
(Section~\ref{l26:sec:verification} lists the commands). ``(record)'' means that the result is taken from the project's
build records (the project's internal build log and workstream logs) and was not re-run by us.

\textbf{The aggregate build.} The project counts a result only when the aggregate build is green. Every file whose
declarations this chapter presents is in the import closure of the umbrella module \lean{VOA.Deligne}; the umbrella build and
the aggregate build that the records report are described in Section~\ref{l26:sec:verification}.

\textbf{Conditional is not unconditional.} A clean \texttt{\#print axioms} never reports explicit hypotheses. A theorem that
takes \lean{(h : P)} for a named \lean{Prop} \lean{P} is \emph{conditional on} \lean{P}, and the status boxes say so.

\subsection{Notation used in this chapter}\label{l26:sec:notation}

Several letters are used for more than one object. Table~\ref{l26:tab:notation} resolves them; the right column says where
each object is defined.
\begin{center}\small
\begin{longtable}{p{0.20\textwidth}p{0.58\textwidth}p{0.14\textwidth}}
\caption{Notation.}\label{l26:tab:notation}\\
\toprule
symbol & meaning & where \\
\midrule
\endhead
$x(r)$ & $x\otimes t^r\in\ltwentysixgt$; $r$ is its degree & \S\ref{l26:sec:current}\\
$e_\alpha$, $f_\alpha$, $h_i$ & raising and lowering root vectors ($\alpha>0$), simple coroots. ``Lowering'' letters are the $f_\alpha$ & \S\ref{l26:sec:current}\\
$\theta$ & the highest root; $V(\theta)=\ltwentysixg$ is the adjoint representation & \S\ref{l26:sec:current}\\
$V(\lambda)$ & the irreducible $\ltwentysixg$-module of highest weight $\lambda$ & \S\ref{l26:sec:fusion}\\
$V(z)$ & the evaluation module at $z\in\C$ & Def.~\ref{l26:def:fusion}\\
$V_1*\dots*V_m$, $V^{*m}$ & fusion product & Def.~\ref{l26:def:fusion}\\
$\ltwentysixcyc(w)_d$ & span of the words of degree $d$ applied to $w$ (Lean \lean{cycDeg}) & \S\ref{l26:sec:current}\\
$q$, $[q^N]f$ & formal variable of characters; the coefficient of $q^N$ (Lean \lean{coeff N}) & \S\ref{l26:sec:current}\\
$\ltwentysixHilb_q$ & graded character $\sum_d\dim(\cdot)_d\,q^d$ & \S\ref{l26:sec:current}\\
$\ltwentysixWloc(\lambda)$ & local Weyl module & \eqref{l26:eq:weyl}\\
$D(1,\lambda)$; $\ltwentysixD_n$ & level-one Demazure module; its character for $\lambda=n\theta$ & \S\ref{l26:sec:current}\\
$D(N)$ & short for $D(1,N\theta)=M_{N,0}$ in Section~\ref{l26:sec:open} only & \S\ref{l26:sec:open}\\
$F_m$; $\ltwentysixF_m$; $\mathrm{fusionF}_m$ & fusion \emph{module} (generators and relations); a \emph{series} solving the mixed recursion; the explicit $A_2$ solution & \eqref{l26:eq:Fm}, Def.~\ref{l26:def:fusionOf}\\
$F_k$, $F(u)$ & operators in Garland's identity (Lean \lean{F}); not related to $F_m$ or to the Lie type $F_4$ & Thm.~\ref{l26:thm:garland}\\
$M_{a,b}$ & Ravinder's modules; $M_{a,b}(q)$ their characters & \eqref{l26:eq:Mab}\\
$M$ (Lean \lean{M}) & any family of power series to which the squeeze is applied & Def.~\ref{l26:def:squeezeinputs}\\
$M_m$ & the sector module of Chapter~\ref{ch:L24} (only in cross-references) & Ch.~\ref{ch:L24}\\
(TB), (KL) & the total bound \lean{FusionTotalBound}; the kernel relation \lean{RavinderKernelRel} & Def.~\ref{l26:def:mab}\\
(R), (D), (T) & the three squeeze inequalities & \S\ref{l26:sec:squeeze-paper}\\
$U_{a,b}$, $w_U$, $M^{\mathrm{univ}}_{a,b}$ & universal module of $M_{a,b}$, its generator, its character (Lean \lean{Mchar}) & Def.~\ref{l26:def:umod}\\
$v_K$ & the kernel vector $f_\theta(2a+b+1)w$ of $M_{a+1,b}$ & \S\ref{l26:sec:ravinder}\\
$\tau_s$ & shift of the grading by $s$ & \eqref{l26:eq:ses}\\
$T_N$ & the evaluation model ($A_2$) & Def.~\ref{l26:def:evalmodel}\\
$T$ & a current module without grading (Lean section variable \lean{T}) in Def.~\ref{l26:def:evalmodel}; Feigin's tensor model in \S\ref{l26:sec:open} & \\
$T_{a,b}$ & mixed evaluation model of the slice lemma & \S\ref{l26:sec:open}\\
$\mathrm{tot}(a,b)$ & total dimension of $U_{a,b}$ & \S\ref{l26:sec:a2red}\\
$J_m$; $J_{a,b}$; $J_N$ & the box ideal of Chari--Pressley; the relation subspace of $U_{a,b}$; the ideal image $(\ltwentysixg\otimes t^N\C[t])\,D(N)$ & Thm.~\ref{l26:thm:box}, Def.~\ref{l26:def:umod}, \S\ref{l26:sec:open}\\
$R$ & the polynomial ring $\C[z_0,z_1,\dots]$ in Thm.~\ref{l26:thm:box}; a relation specification in \S\ref{l26:sec:generic} & \\
$Q(r,s)$ & $[u^s]F(u)^r$, the Garland coefficient (also $Q(r,j)$ applied to $w$) & Thm.~\ref{l26:thm:garland}\\
$\Phi_{<d}$, $\Phi_{\le d}$ & the degree filtration of a cyclic module & Def.~\ref{l26:def:fusion}, \ref{l26:def:evalmodel}\\
$S$ & sector data (Chapter~\ref{ch:L24}); in \S\ref{l26:sec:open}, a kernel subspace & \\
\bottomrule
\end{longtable}
\end{center}

\subsection{Reading the Lean in this chapter}\label{l26:sec:readlean}

Chapter~\ref{ch:01}, Section~\ref{ch01:sec:lean}, is the general primer. The idioms that the printed code of this chapter
needs are these.
\begin{itemize}
\item \textbf{Truncated blocks.} Proof bodies are omitted. A printed block therefore ends at \lean{:=}, \lean{:= by} or
  \lean{where}, and the proof follows in the file. This is not a syntax error.
\item \textbf{Section variables.} A Lean \lean{section} may declare \lean{variable {V : Type} {ρ : ...}} once; every
  declaration of the section then takes these as arguments, although they are not printed in its statement. Hypotheses
  declared this way (for example \lean{(hXY : X * Y - Y * X = Z)}) become arguments of a theorem when the theorem uses them or
  when they are \lean{include}d. The dictionaries list the section variables of every quoted declaration.
\item \textbf{Binders.} \lean{(x : T)} is an explicit argument, \lean{{x : T}} an implicit one inferred by Lean, and
  \lean{[FiniteDimensional ℂ V]} an instance argument: a hypothesis that Lean finds automatically where it is known.
\item \textbf{Named statements are not axioms.} \lean{def P : Prop := ...} only gives a \emph{name} to a statement; it asserts
  nothing. A theorem with a hypothesis \lean{(h : P)} is proved \emph{assuming} \lean{P}. This chapter's honesty rests on this
  distinction: \lean{FusionTotalDim} and its relatives are named statements, and nothing in the library postulates them.
  A named statement becomes a theorem when the library constructs a proof of it, as \lean{CVUniv.fusionTotalDim} does for
  \lean{FusionTotalDim}.
  \lean{structure S : Prop where f1 : ... f2 : ...} is a bundle of named conditions; a proof \lean{h : S} gives \lean{h.f1}, \lean{h.f2}.
\item \textbf{\lean{abbrev} vs \lean{def}}: both define; \lean{abbrev} unfolds automatically.
\item \textbf{Lists and words.} A word of current letters is a \lean{List (Fin 8 × ℕ)}: a letter $(c,r)$ is $x_c(r)$.
  \lean{Fin 8} $=\{0,\dots,7\}$. \lean{List.replicate k x} is the list $x,\dots,x$ of length $k$.
\item \textbf{Finite sums.} \lean{∑ i ∈ s, f i} is a sum over the finite set \lean{s}; \lean{Finset.range n}
  $=\{0,\dots,n-1\}$; \lean{Finset.antidiagonal n} $=\{(i,j):i+j=n\}$.
\item \textbf{Small notation.} \lean{•} is scalar multiplication; \lean{f^[n]} is the $n$-fold iterate of \lean{f};
  \lean{(ρ.op 5 0) ^ n} is the $n$-th power in $\End(V)$; \lean{x.1}, \lean{x.2} are the components of a pair; \lean{ℂ ∙ w} is
  the line through $w$; natural-number subtraction is truncated at $0$.
\item \textbf{Linear algebra.} \lean{Module.finrank ℂ V} is $\dim V$, \emph{and is $0$ when $V$ is infinite-dimensional};
  \lean{FiniteDimensional ℂ V} says $\dim V<\infty$. \lean{Submodule.span ℂ s} is the span; \lean{V ⧸ U} the quotient and
  \lean{U.mkQ} the quotient map. \lean{PowerSeries ℚ} is $\Q[[q]]$, and \lean{X} is $q$.
\item \textbf{Kernel checks.} \lean{decide +kernel} proves a decidable statement by evaluating it in Lean's trusted kernel.
  \lean{native_decide}, which would trust compiled code, is never used in these files.
\end{itemize}

%======================================================================
\section{The mathematics}\label{l26:sec:math}
%======================================================================

\subsection{Current algebras and their cyclic modules}\label{l26:sec:current}

Let $\ltwentysixg$ be a finite-dimensional simple Lie algebra over $\C$ with a basis $(x_c)_{c\in\sigma}$ and
structure constants $f$: $[x_a,x_b]=\sum_e f(a,b,e)\,x_e$. The \emph{current algebra} is
$\ltwentysixgt=\ltwentysixg\otimes\C[t]$ with bracket $[x\otimes t^r,y\otimes t^s]=[x,y]\otimes t^{r+s}$. We write
$x(r)$ for $x\otimes t^r$. It is graded by the $t$-degree: $x(r)$ has degree $r$.

A $\ltwentysixgt$-module is a vector space $V$ with operators $x(r)$ obeying the bracket. If $w\in V$, the
\emph{cyclic submodule} generated by $w$ is spanned by the \emph{loop words} $x_{c_1}(r_1)\cdots x_{c_\ell}(r_\ell)w$;
the degree of such a word is $r_1+\dots+r_\ell$. We write $\ltwentysixcyc(w)_d$ for the span of the words of degree
$d$ applied to $w$ (Lean: \lean{cycDeg ρ w d}). When $V$ is graded and $w$ is homogeneous, $\ltwentysixcyc(w)_d$ is
exactly the degree-$d$ part of the cyclic module; in general it is just a subspace, and the library works with it
because \lean{CurrentModule} carries no grading operator. The \emph{graded character} of a cyclic module is
$\ltwentysixHilb_q=\sum_d\dim\ltwentysixcyc(w)_d\,q^d$, and $[q^N]f$ denotes the coefficient of $q^N$ in a series $f$.

Fix a Cartan subalgebra, positive roots, the highest root $\theta$, simple roots $\alpha_i$ with coroots $h_i$, and
Chevalley generators $e_\alpha$ (raising) and $f_\alpha$ (lowering), $\alpha>0$. For a dominant weight $\lambda$ the
\textbf{local Weyl module} $\ltwentysixWloc(\lambda)$ (Chari--Pressley) is the $\ltwentysixgt$-module generated by $w$ with
relations
\begin{equation}\label{l26:eq:weyl}
e_\alpha(r)w=0,\qquad h(r)w=\delta_{r,0}\,\lambda(h)\,w,\qquad f_i(0)^{\lambda(h_i)+1}w=0 .
\end{equation}
It is finite-dimensional and graded. The \textbf{level-one Demazure module} $D(1,\lambda)$ is the $\ltwentysixgt$-submodule
of the basic level-one representation of the affine algebra generated by an extremal weight vector; for $\lambda=n\theta$ and
$\ltwentysixg$ simply laced it is isomorphic to $\ltwentysixWloc(n\theta)$ (Chari--Loktev for $\mathfrak{sl}_{r+1}$,
Fourier--Littelmann in general). We write
\begin{equation*}
\ltwentysixD_n(q)=\ltwentysixHilb_q D(1,n\theta),\qquad \ltwentysixD_n(1)=\dim D(1,n\theta)=(\dim\ltwentysixg+1)^n .
\end{equation*}
For $\mathfrak{sl}_3$: $\ltwentysixD_1=8+q$, $\ltwentysixD_2=27+28q+17q^2+8q^3+q^4$, totals $9$ and $81$. As
$\ltwentysixg$-modules, $D(1,\theta)=\ltwentysixg\oplus\C$ (adjoint in degree $0$, trivial in degree $1$).

\begin{example}[$\mathfrak{sl}_2$: the local Weyl module $\ltwentysixWloc(m)$]\label{l26:ex:sl2}
For $\mathfrak{sl}_2$ with generators $e,h,f$ and $\lambda=m$, Chari--Pressley prove that the vectors
$f(s_1)\cdots f(s_\ell)w$ with $0\le s_1\le\dots\le s_\ell\le m-\ell$ form a basis of $\ltwentysixWloc(m)$. For $m=2$:
$w$; $f(0)w$, $f(1)w$; $f(0)^2w$: four vectors, and $\dim\ltwentysixWloc(m)=\sum_\ell\binom m\ell=2^m$. The constraint
$s_\ell\le m-\ell$ (a ``box'') is the content of the theorem; it comes from Garland's identities
(Section~\ref{l26:sec:garland}).
\end{example}

\subsection{Fusion products, concretely}\label{l26:sec:fusion}

\begin{definition}[Evaluation modules; fusion products]\label{l26:def:fusion}
(a) For a $\ltwentysixg$-module $V$ and $z\in\C$ the \textbf{evaluation module} $V(z)$ is $V$ with $x(r)$ acting by $z^rx$
(with $0^0=1$). At $z=0$ every $x(r)$, $r\ge1$, acts by $0$, and $V(0)$ is a graded module concentrated in degree $0$.
We write $V(\lambda)$ for the irreducible $\ltwentysixg$-module of highest weight $\lambda$; $V(\theta)=\ltwentysixg$.

(b) More generally, for a graded cyclic $\ltwentysixgt$-module $V$ and $z\in\C$, the \textbf{translate} $V^{z}$ is $V$ with
$x(r)$ acting as $x\otimes(t+z)^r=\sum_s\binom rs z^{r-s}x(s)$. For $V=V(\lambda)(0)$ this is $V(\lambda)(z)$.

(c) Let $V_1,\dots,V_m$ be graded cyclic $\ltwentysixgt$-modules with cyclic vectors $v_j$, and $z_1,\dots,z_m$ distinct.
Put $T=V_1^{z_1}\otimes\dots\otimes V_m^{z_m}$ (with $x(r)$ acting by the Leibniz rule) and $v=v_1\otimes\dots\otimes v_m$.
Feigin and Loktev show that $T$ is cyclic on $v$. Its \textbf{degree filtration} is
\begin{equation*}
\Phi_{\le d}=\operatorname{span}\{x_1(r_1)\cdots x_\ell(r_\ell)v:\ r_1+\dots+r_\ell\le d\},\qquad \Phi_{<d}=\Phi_{\le d-1}.
\end{equation*}
Then $\Phi_{\le d}\subseteq\Phi_{\le d+1}$, $x(r)\Phi_{\le d}\subseteq\Phi_{\le d+r}$, and $\bigcup_d\Phi_{\le d}=T$. The
\textbf{fusion product} is the associated graded space
\begin{equation*}
V_1*\dots*V_m=\bigoplus_{d\ge0}\Phi_{\le d}/\Phi_{<d},
\end{equation*}
on which $x(r)$ acts by the map $\Phi_{\le d}/\Phi_{<d}\to\Phi_{\le d+r}/\Phi_{<d+r}$ that it induces. We write $V^{*m}$ for
$m$ factors equal to $V$.
\end{definition}

The fusion product is a graded cyclic $\ltwentysixgt$-module (generated by the class of $v$), and its dimension is
$\prod_j\dim V_j$, because an associated graded space has the dimension of the filtered space. Its definition depends on the
points $z_j$. That it does not is a conjecture of Feigin and Loktev; it is a theorem whenever the fusion product is shown to
have a presentation by generators and relations that does not mention the points, since then all choices give the same
module. For $V(\theta)^{*m}$ that presentation is Ravinder's \eqref{l26:eq:Fm}, a literature result resting on the input
discussed in the Warning at the end of this section.

\textbf{Why the fusion product satisfies truncation relations.} Let all $V_j=V(\lambda_j)(0)$. For a polynomial $P(t)$,
$x\otimes P(t)$ acts on $T$ by $\sum_jP(z_j)\,x^{(j)}$, where $x^{(j)}$ acts on the $j$-th factor. Take $P(t)=\prod_j(t-z_j)$,
of degree $m$: then $x\otimes P(t)$ acts by $0$, and $P(t)=t^m+(\text{lower})$ gives $x(m)v\in\Phi_{<m}$. In the associated
graded space this becomes the exact relation $x(m)\bar v=0$. This is how relations of the shape $f_\alpha(m)w=0$ arise.

\begin{example}[$\mathfrak{sl}_3$: the fusion square of the adjoint]\label{l26:ex:fus88}
Take $\ltwentysixg=\mathfrak{sl}_3$, $V_1=V_2=V(\theta)$, the points $z_1=1$, $z_2=-1$, and $v=e_\theta\otimes e_\theta$. As
$\ltwentysixg$-modules,
\begin{equation*}
V(\theta)\otimes V(\theta)=8\otimes8=\underbrace{27\oplus8\oplus1}_{\operatorname{Sym}^2}\oplus\underbrace{10\oplus\overline{10}\oplus8}_{\Lambda^2},
\end{equation*}
of dimension $64$. Let $\sigma$ be the swap of the two factors. Since $z_2=-z_1$, $\sigma x(r)\sigma=(-1)^rx(r)$, and
$\sigma v=v$; so $\sigma$ preserves each $\Phi_{\le d}$ and acts on $\Phi_{\le d}/\Phi_{<d}$ by $(-1)^d$. Hence the even-degree
pieces of the fusion product are, as $\ltwentysixg$-modules, the pieces of $\operatorname{Sym}^2$, and the odd-degree pieces those
of $\Lambda^2$. Degree $0$ is $U(\ltwentysixg)v=V(2\theta)$, of dimension $27$. So
\begin{equation*}
\text{degree }0:\ 27,\qquad \text{degree }1:\ 10\oplus\overline{10}\oplus8\ (=28),\qquad \text{degree }2:\ 8\oplus1\ (=9),
\end{equation*}
and the graded character is $27+28q+9q^2$; the $8$ in degree $1$ is the antisymmetric one ($[x,y]$), the $8$ in degree $2$ the
symmetric one. An exact rational computation of the filtration, made for this chapter (not part of the repository), confirms
this decomposition degree by degree at the points $(1,-1)$ and at $(0,1)$. The same numbers come out of the mixed recursion
\eqref{l26:eq:mixed} below: $\ltwentysixD_2=\ltwentysixF_2+q^2(1+q)\ltwentysixF_1+q^4$ with $\ltwentysixD_2=27+28q+17q^2+8q^3+q^4$ and
$\ltwentysixF_1=8$ gives $\ltwentysixF_2=27+28q+9q^2$. So at $m=2$ representation theory, the recursion and the total $8^2$ agree.
(In Lean, $\dim F_2\le64$ is the case $m=2$ of \lean{CVUniv.fusionTotalDim}, Theorem~\ref{l26:thm:cvtotal}.)
\end{example}

\begin{example}[$\mathfrak{sl}_3$: the Demazure module $D(1,\theta)$ as a fusion product]\label{l26:ex:eval1}
Take the standard representation $\C^3$ at $z=0$ and its dual $(\C^3)^*$ at $z=1$, with $v=e_1\otimes e_3^*$ (weight $\theta$).
Then $T=\C^3\otimes(\C^3)^*\cong\operatorname{End}\C^3$ has dimension $9$. Degree $0$ is $U(\ltwentysixg)v$, the traceless
matrices, $V(\theta)=8$. In degree $1$, $x(1)$ acts by $0\cdot x^{(1)}+1\cdot x^{(2)}$, the action on the dual factor alone, which
reaches the identity matrix; so degree $1$ is the trivial module $\C$. The graded character is $8+q=\ltwentysixD_1$. This is the
case $N=1$ of the evaluation model of Theorem~\ref{l26:thm:u7}.
\end{example}

\subsection{The fusion modules $F_m$}

The relations satisfied by the fusion product $V(\theta)^{*m}$ suggest the module
\begin{equation}\label{l26:eq:Fm}
F_m=\ltwentysixWloc(m\theta)\big/\bigl\langle f_\alpha(m)w:\ \alpha>0\bigr\rangle .
\end{equation}
By the paragraph on truncation relations, $V(\theta)^{*m}$ is a quotient of $F_m$, so $\dim F_m\ge(\dim\ltwentysixg)^m$.
Ravinder's presentation theorem says that $F_m\cong V(\theta)^{*m}$ (stated for every simple $\ltwentysixg$). Its proof takes the
terminal case from Feigin's Cor.~2 (Ravinder's Prop.~5), which is exactly the upper bound $\dim F_m\le(\dim\ltwentysixg)^m$ and has
the gap recorded in the Warning below; Kus--Littelmann's Thm.~4.2 proves the same bound for every simple $\ltwentysixg$, with one
step flagged by the project (Chapter~\ref{ch:L27}, item~(i) of the literature paragraph); in Lean that upper bound is the statement \lean{FusionTotalDim}, proved for $\mathfrak{sl}_3$
by another route (Theorem~\ref{l26:thm:cvtotal}) and open for the other types (\lean{FusionTotalDimG}). So: \emph{if} Ravinder's identification holds, then $\ltwentysixHilb_qF_m(1)=(\dim\ltwentysixg)^m$ and
$[q^0]\ltwentysixHilb_qF_m=\dim V(m\theta)$.

We write $\ltwentysixF_m(q)$ for the series defined by inverting the mixed recursion \eqref{l26:eq:mixed} below (in Lean
\lean{fusionOf}, Definition~\ref{l26:def:fusionOf}). For $\mathfrak{sl}_3$ these are
\begin{equation*}
\ltwentysixF_1=8,\qquad \ltwentysixF_2=27+28q+9q^2,\qquad \ltwentysixF_3=64+97q+152q^2+111q^3+71q^4+16q^5+q^6 .
\end{equation*}
They are \emph{outputs of the recursion} applied to the Demazure characters, with totals $8$, $64$, $512$; they are not computed
characters of the modules $F_m$. That they \emph{are} the characters of $F_m$ is literature (Ravinder) plus numerics (exact through
$m=4$), and in Lean a consequence of U5 (Theorem~\ref{l26:thm:squeeze}). Example~\ref{l26:ex:fus88} checks $m=2$ directly.

The statement $\dim F_m\le(\dim\ltwentysixg)^m$ is \textbf{U5 in total form}. In the Lean statements below, the
relations \eqref{l26:eq:weyl} plus $f_\alpha(m)w=0$ are \lean{LocalWeylRel} ($A_2$) and \lean{LocalWeylRelG}
(general), Definition~\ref{l24:def:U5}.

\subsection{Ravinder's modules $M_{a,b}$ and the exact sequence}\label{l26:sec:ravinder}

Ravinder (arXiv:1409.0274) interpolates between the Demazure module and the fusion module with a two-parameter family.
For $a,b\ge0$ put $k=a+b$ and let
\begin{equation}\label{l26:eq:Mab}
M_{a,b}=\ltwentysixWloc(k\theta)\Big/\Bigl\langle f_\alpha(k)w\ (\alpha>0,\ \alpha\neq\theta),\ \ f_\theta(2a+b)w\Bigr\rangle .
\end{equation}
Then $M_{0,m}=F_m$, and $M_{n,0}=\ltwentysixWloc(n\theta)=D(1,n\theta)$ ($\ltwentysixg$ simply laced): at $(a,b)=(n,0)$
every truncation already holds in $\ltwentysixWloc(n\theta)$, because $w$ is a highest-weight vector of weight
$(n\theta)(h_\alpha)\in\{n,2n\}$ for the $\mathfrak{sl}_2$-triple of $\alpha$, and the $\mathfrak{sl}_2$ box of
Example~\ref{l26:ex:sl2} kills $f_\alpha(s)w$ for $s$ at least that weight. Ravinder identifies $M_{a,b}$ with the
fusion product $D(1,\theta)^{*a}*V(\theta)^{*b}$.

\begin{remark}[The Lean fusion sets]\label{l26:rem:fusionset}
The Lean sectors impose $f_\alpha(k)w=0$ for all positive $\alpha$ for $A_2$, $D_4$, $G_2$, $F_4$, but only for
$\{\alpha:\langle\alpha,\theta^\vee\rangle=1\}\cup\{\theta\}$ for $E_6$, $E_7$, $E_8$ (Ravinder's choice;
$\langle\alpha,\theta^\vee\rangle=2(\alpha,\theta)/(\theta,\theta)$, which for a short root of $G_2$ or $F_4$ differs from
$(\theta,\alpha^\vee)$). For simply laced types the two
readings give the same module: for $\alpha\perp\theta$ the $\mathfrak{sl}_2[t]$-submodule of the $\alpha$-triple generated by $w$
is a local Weyl module of weight $0$, hence trivial, so $f_\alpha(r)w=0$ for all $r$ already (Chapter~\ref{ch:L24}). For
$\mathfrak{sl}_3$ the positive roots $\alpha_1,\alpha_2,\theta$ all have $(\theta,\alpha^\vee)\ge1$, so nothing changes. The same
holds for $G_2$ and $F_4$: since $\theta$ is long, every positive $\alpha\neq\theta$ has $\langle\alpha,\theta^\vee\rangle\in\{0,1\}$; those
with value $1$, short roots included, are in Ravinder's set, and for those with value $0$ one has $f_\alpha(0)w=0$ and then
$f_\alpha(r)w=0$ for all $r$ by $[h(1),f_\alpha(r)]=-\alpha(h)f_\alpha(r+1)$ and $h(1)w=0$ (Chapter~\ref{ch:L27},
Lemma~\ref{l27:lem:truncweyl}). So the two readings give the same module for every type.
\end{remark}

For a graded module $M$ and $s\in\N$ let $\tau_sM$ be $M$ with every degree raised by $s$. Ravinder proves the short exact
sequence
\begin{equation}\label{l26:eq:ses}
0\longrightarrow \tau_{2a+b+1}\,M_{a,b}\longrightarrow M_{a+1,b}\longrightarrow M_{a,b+1}\longrightarrow 0 .
\end{equation}
The surjection exists because the relations of $M_{a,b+1}$ are those of $M_{a+1,b}$ (same $k=a+b+1$) plus
$f_\theta(2a+b+1)w=0$; its kernel is generated by the \textbf{kernel vector} $v_K=f_\theta(2a+b+1)w$. The left map sends the
generator of $M_{a,b}$ to $v_K$; that it is well defined is \textbf{Ravinder's kernel lemma}: \emph{$v_K$ satisfies the relations
of $M_{a,b}$}. Injectivity of the left map is the deep half.

On characters, write $M_{a,b}(q)$. The sequence gives $M_{a+1,b}=M_{a,b+1}+q^{2a+b+1}M_{a,b}$. Iterating from $(n,0)$
down to the row $a=0$ expresses the Demazure character through the fusion characters:
\begin{equation}\label{l26:eq:mixed}
\ltwentysixD_n(q)=\sum_{d=0}^{n}q^{dn}\ltwentysixqb nd\,\ltwentysixF_{n-d}(q),
\end{equation}
with $\ltwentysixqb nd$ the Gaussian binomial. This is \textbf{U6} (Ravinder Cor.~1). For example
$\ltwentysixD_1=\ltwentysixF_1+q$ and $\ltwentysixD_2=\ltwentysixF_2+q^2(1+q)\ltwentysixF_1+q^4$.

\begin{physics}
\textbf{Dictionary of the mnemonics (heuristic throughout this chapter).} In conformal field theory a current algebra
$\ltwentysixgt$ is the algebra of non-negative modes of the chiral currents $J^a(z)$ of a WZW-type model. A cyclic module
generated by $w$ is the space of states created from a primary state by the currents. In the physics boxes of this chapter:
a \emph{particle} is a tensor factor $V_j(z_j)$, i.e.\ a point $z_j$ where a primary field (in the adjoint, or in $D(1,\theta)$) is
inserted; the \emph{energy} of a state is its $t$-degree (a loop grading; it is the eigenvalue of $L_0$ minus a constant only in
the lattice picture of Section~\ref{l26:sec:lattice}); a \emph{bound} particle is a $D(1,\theta)=\ltwentysixg\oplus\C$ factor in its
trivial component $\C$, a \emph{free} one an adjoint factor; the \emph{binding energy} $2a+b+1$ is the grading shift of the left map
in \eqref{l26:eq:ses}. There is no Hamiltonian in this chapter; the words are aids to memory, never part of a proof.

With this dictionary, the fusion product is the operator product expansion made algebraic: $m$ primaries sit at points
$z_1,\dots,z_m$ and are brought together; the associated graded space records, energy level by energy level, which
descendant states survive the collision. The Demazure module $D(1,n\theta)$ is the space of states of a level-one chiral boson
(a lattice vertex algebra) generated by the currents from the vertex operator $e^{n\theta}$. Ravinder's sequence splits these
states according to whether one of the particles is bound or not.
\end{physics}

\subsection{The squeeze}\label{l26:sec:squeeze-paper}

Suppose we have power series $M_{a,b}(q)$ with nonnegative coefficients satisfying
\begin{enumerate}
\item[(R)] $M_{a+1,b}\le M_{a,b+1}+q^{2a+b+1}M_{a,b}$ coefficientwise (the \emph{spanning half} of \eqref{l26:eq:ses});
\item[(D)] $\ltwentysixD_n\le M_{n,0}$ coefficientwise (a Demazure \emph{lower} bound);
\item[(T)] $\sum_d[q^d]M_{0,m}\le(\dim\ltwentysixg)^m$ (the \emph{total} upper bound).
\end{enumerate}
Then $M_{0,m}=\ltwentysixF_m$ for every $m$, where $\ltwentysixF$ is the solution of \eqref{l26:eq:mixed}. Indeed,
iterating (R) gives $M_{n,0}\le\sum_dq^{dn}\ltwentysixqb nd M_{0,n-d}$; by induction on $n$ all terms with $d\ge1$ equal
the corresponding terms of \eqref{l26:eq:mixed}; with (D) this leaves $\ltwentysixF_n\le M_{0,n}$ coefficientwise.
Both sides have total $(\dim\ltwentysixg)^n$ at $q=1$ (the left side by the count, the right side by (T)), and an
inequality between series with nonnegative differences and equal sums is an equality. This argument is Ravinder's
Cor.~1; the library formalizes it for every type. Applied to the characters of \emph{universal} modules
(Section~\ref{l26:sec:univ}), it yields U5 for $\ltwentysixF$.

\subsection{Where this chapter enters the Deligne chain}

In the chain of Chapter~\ref{ch:L24}, Theorem~\ref{l24:thm:B6} (freeness of $M_m$) and Theorem~\ref{l24:thm:sfb} (the
sector fusion bound) consume U5 as \lean{FusionModuleBoundG S F}, and the resummation D3
(Theorem~\ref{l24:thm:D3}) consumes U6 as \lean{MixedFusionRecursion F Dem}. The present chapter produces, for each
type, a family \lean{F} and proofs of these two structures from smaller inputs:
\begin{equation*}
\underbrace{\text{U6}}_{\text{by construction}},\qquad
\underbrace{\text{U5}}_{\text{squeeze}}\ \Longleftarrow\ \underbrace{\text{kernel lemma}}_{\text{proved}}
+\underbrace{\text{U7}}_{A_2:\ \text{proved}}+\underbrace{\text{finiteness}}_{\text{proved}}
+\underbrace{\dim\le(\dim\ltwentysixg)^m}_{A_2:\ \text{proved; others named}}+\underbrace{\text{count}}_{\text{Chapter~\ref{ch:L25}}} .
\end{equation*}
The final per-type theorems that assemble these with the carrier and character sides are in Chapter~\ref{ch:L27}.

\begin{warning}
\textbf{The literature behind the total bound, and its gap.} The dimension half of the total bound,
$\dim F_m\le(\dim\ltwentysixg)^m$, is the statement \lean{FusionTotalDim} for $A_2$ and \lean{FusionTotalDimG} for the other
types. The
standard citation is Feigin, \emph{The PBW filtration}, SIGMA 4 (2008) 070, arXiv:0806.4851, Cor.~2 (Prop.~2.1 and Cor.~2.2 in arXiv v1),
which Ravinder cites for the terminal case of his sequence. For $A_2$, Feigin's Prop.~1, on which Cor.~2 rests, is a sketch whose natural
reading is false (Section~\ref{l26:sec:open}). Kus and Littelmann (Pacific J.\ Math.\ 278 (2015), Thm.~3.5, Cor.~3.6(3), Thm.~4.2)
acknowledge Feigin's gap and repair it, for every simple $\ltwentysixg$ (with $\lambda=\theta$, $\lambda_0=0$ their Thm.~4.2 has no type
restriction), and Setia--Rani--Khandai (arXiv:2405.11944, \S8) give an independent upper-bound route through
Chari--Venkatesh modules. One passage of the Kus--Littelmann proof (transfer of the annihilator ideal to the associated graded, p.~437) needs
further justification before it can be formalized; this is not a counterexample. For $D_4$ at $m=2$ a computation (not Lean) shows
that the passage as written is insufficient, with an explicit witness, while the statement it is meant to prove holds there
(Chapter~\ref{ch:L27}, Section~\ref{l27:sec:feigin}); the theorem is not known to fail. In Lean the statement is \emph{proved for
$A_2$} (\lean{CVUniv.fusionTotalDim}, Theorem~\ref{l26:thm:cvtotal}), along the Setia--Rani--Khandai route and without Feigin's
Prop.~1; it is a named, unproved hypothesis for $G_2$, $D_4$, $F_4$, $E_6$, $E_7$, $E_8$. For these six types the published
proof is Kus--Littelmann's Thm.~4.2, and the flagged p.~437 step applies to them equally, since the specialization is the same
(Chapter~\ref{ch:L27}, Table~\ref{l27:tab:precise}).
\end{warning}

%======================================================================
\section{The $A_2$ fusion modules in Lean}\label{l26:sec:a2rel}
%======================================================================

Throughout the $A_2$ sections the Lie algebra is $\mathfrak{sl}_3$ in the basis of \lean{AffineSl3.bracketTable}:
indices $0,1,2$ are $e_1,e_2,e_\theta$, indices $3,4,5$ are $f_1,f_2,f_\theta$, and $6,7$ are $h_1,h_2$; $\theta=\alpha_1+\alpha_2$
and $\theta(h_1)=\theta(h_2)=1$. A \lean{CurrentModule (bracketTable) V} is a family of linear maps
\lean{ρ.op c r}$=x_c(r)$ on $V$ satisfying the current-algebra bracket, \lean{wordOp ρ l w} applies a loop word
(last letter first), and \lean{cycDeg ρ w d} is $\ltwentysixcyc(w)_d$ (Chapter~\ref{ch:L24},
Definition~\ref{l24:def:U5}, where these are quoted).

\begin{definition}[The $A_2$ relations and two named statements]\label{l26:def:mab}
Let $\rho$ be a current module and $w\in V$.
\begin{enumerate}
\item[(a)] $w$ satisfies the \emph{relations of $M_{a,b}$} if, with $k=a+b$:
$e_c(r)w=0$ for $c=e_1,e_2,e_\theta$ and all $r$; $h(r)w=0$ for $r\ge1$ and $h(0)w=k\,w$ for $h=h_1,h_2$;
$f_1(0)^{k+1}w=f_2(0)^{k+1}w=0$; $f_1(k)w=f_2(k)w=0$; and $f_\theta(2a+b)w=0$.
\item[(b)] \textbf{(TB)}, the \emph{total bound} \lean{FusionTotalBound}: for every $m$ and every module with the local Weyl and
fusion relations for $m\theta$, every $\ltwentysixcyc(w)_d$ is finite-dimensional and $\sum_{d<B}\dim\ltwentysixcyc(w)_d\le8^m$
for every $B$.
\item[(c)] \textbf{(KL)}, the \emph{kernel relation} \lean{RavinderKernelRel}: for $a\ge1$, if $w$ satisfies the
relations of $M_{a,b}$, then $f_\theta(2a+b-1)w$ satisfies those of $M_{a-1,b}$.
\end{enumerate}
\end{definition}
\noindent File \leanpath{VOA/\allowbreak DeligneAudit/\allowbreak DeligneFusionRel.lean}.
\begin{leancode}
structure MabRel {V : Type} [AddCommGroup V] [Module ℂ V]
    (ρ : CurrentModule (bracketTable) V) (w : V) (a b : ℕ) : Prop where
  raising : ∀ (e : Fin 8), e.val < 3 → ∀ r, ρ.op e r w = 0
  cartan_pos : ∀ (h : Fin 8), 6 ≤ h.val → ∀ r, 1 ≤ r → ρ.op h r w = 0
  cartan_zero : ∀ (h : Fin 8), 6 ≤ h.val → ρ.op h 0 w = ((a + b : ℕ) : ℂ) • w
  lowering_f1 : (ρ.op 3 0)^[a + b + 1] w = 0
  lowering_f2 : (ρ.op 4 0)^[a + b + 1] w = 0
  trunc_f1 : ρ.op 3 (a + b) w = 0
  trunc_f2 : ρ.op 4 (a + b) w = 0
  trunc_theta : ρ.op 5 (2 * a + b) w = 0

def FusionTotalBound : Prop :=
  ∀ (m : ℕ) (V : Type) [AddCommGroup V] [Module ℂ V] (ρ : CurrentModule (bracketTable) V) (w : V),
    LocalWeylRel ρ w m → (∀ b, FiniteDimensional ℂ (cycDeg ρ w b)) ∧
      ∀ B : ℕ, ∑ b ∈ Finset.range B, Module.finrank ℂ (cycDeg ρ w b) ≤ 8 ^ m

def RavinderKernelRel : Prop :=
  ∀ (a b : ℕ) (V : Type) [AddCommGroup V] [Module ℂ V] (ρ : CurrentModule (bracketTable) V) (w : V),
    1 ≤ a → MabRel ρ w a b → MabRel ρ (ρ.op 5 (2 * a + b - 1) w) (a - 1) b
\end{leancode}

\begin{dictionary}
\textbf{Decoding.} \lean{(ρ.op 3 0)^[k] w} is the $k$-fold iterate of $f_1(0)$ applied to $w$. The Cartan fields say
$h_i(0)w=(a+b)w$, i.e.\ $w$ has weight $(a+b)\theta$, since $\theta(h_i)=1$. The truncations \lean{trunc_f1},
\lean{trunc_f2}, \lean{trunc_theta} are $f_1(k)w=f_2(k)w=0$ and $f_\theta(2a+b)w=0$. These are exactly the relations
\eqref{l26:eq:Mab}: for $\mathfrak{sl}_3$ the positive roots are $\alpha_1,\alpha_2,\theta$, all ``fusion roots''.
\lean{mabRel_zero_iff} (same file) proves \lean{MabRel ρ w 0 m ↔ LocalWeylRel ρ w m}, so $M_{0,m}$ is the fusion
module \eqref{l26:eq:Fm}.

\textbf{Index conventions.} In \lean{RavinderKernelRel} the indices are shifted by one relative to
\eqref{l26:eq:ses}: with $a\ge1$ it maps $M_{a,b}$ to $M_{a-1,b}$ via $f_\theta(2(a-1)+b+1)=f_\theta(2a+b-1)$. The natural
subtraction \lean{2 * a + b - 1} and \lean{a - 1} are honest because of the hypothesis \lean{1 ≤ a}.

\textbf{The finiteness clause in \lean{FusionTotalBound}.} \lean{Module.finrank} of an infinite-dimensional space is $0$
in Mathlib, so a bare dimension inequality would say nothing about a module with an infinite-dimensional degree piece.
The clause \lean{FiniteDimensional ℂ (cycDeg ρ w b)} closes that loophole. It is proved below
(Theorem~\ref{l26:thm:finite}), so the only content left is the dimension sum.

\textbf{How it could be vacuous, and why it is not.} \lean{FusionTotalBound} quantifies over \emph{all} modules
satisfying the relations, including the universal one (Section~\ref{l26:sec:univ}), so it is a genuine upper bound on
$\dim F_m$; it is not satisfied trivially. \lean{RavinderKernelRel} could be vacuous only if no module satisfied
\lean{MabRel} with a nonzero vector. The zero action does \emph{not} provide one: \lean{cartan_zero} then reads
$0=(a+b)w$, which forces $w=0$ because $a\ge1$. The lattice module (\lean{latRep_mabRel}, Theorem~\ref{l26:thm:latrep})
covers only $(a,b)=(N,0)$. The witness for every $(a,b)$ is the universal module: \lean{univ_mabRel}
(Definition~\ref{l26:def:umod}) proves \lean{MabRel (univRep a b) (wU a b) a b}. That $w_U\neq0$ is not stated in Lean; on
paper it holds because $U_{a,b}$ maps onto the irreducible $\mathfrak{sl}_3$-module $V((a+b)\theta)$ with $x(r)$ acting by $0$ for
$r\ge1$, whose highest-weight vector satisfies every relation of \lean{MabRel} (all truncations have positive mode index when
$a+b\ge1$).
\end{dictionary}

\begin{physics}
(Heuristic, in the dictionary of Section~\ref{l26:sec:ravinder}.) $w$ is a highest-weight state of weight $(a+b)\theta$, the state
of $a+b$ particles fused together; the truncations say that the modes $f_\alpha(k)$ above a threshold annihilate it.
The threshold for $f_\theta$ distinguishes the $a$ factors that carry the Demazure module $D(1,\theta)=\ltwentysixg\oplus\C$
(two units of degree for $f_\theta$) from the $b$ factors in the adjoint representation (one unit).
\end{physics}

\begin{status}
Definitions. (TB) is reduced to \lean{FusionTotalDim} (Theorem~\ref{l26:thm:finite}); its literature sources are Feigin's Cor.~2
(with a gap) and Kus--Littelmann's Thm.~4.2 (with one step flagged by the project); Ravinder's Cor.~2 at $a=0$ quotes Feigin. (KL) is a theorem (Theorem~\ref{l26:thm:kernelA2}).
\end{status}

\begin{theorem}[Finiteness clause for $\mathfrak{sl}_3$]\label{l26:thm:finite}
If $w$ satisfies the local Weyl and fusion relations for $m\theta$, then every $\ltwentysixcyc(w)_b$ is
finite-dimensional. Consequently \lean{FusionTotalBound} is equivalent to its dimension part \lean{FusionTotalDim}: for
every module with these relations and every $B$, $\sum_{b<B}\dim\ltwentysixcyc(w)_b\le8^m$.
\end{theorem}
\noindent File \leanpath{VOA/\allowbreak DeligneAudit/\allowbreak DeligneFusionFinite.lean}.
\begin{leancode}
theorem cycDeg_finiteDimensional {w : V} {m : ℕ} (hW : LocalWeylRel ρ w m) (b : ℕ) :
    FiniteDimensional ℂ (cycDeg ρ w b) := by

def FusionTotalDim : Prop :=
  ∀ (m : ℕ) (V : Type) [AddCommGroup V] [Module ℂ V] (ρ : CurrentModule (bracketTable) V) (w : V),
    LocalWeylRel ρ w m → ∀ B : ℕ, ∑ b ∈ Finset.range B, Module.finrank ℂ (cycDeg ρ w b) ≤ 8 ^ m

theorem fusionTotalBound_of_dim (h : FusionTotalDim) : FusionTotalBound :=
\end{leancode}

\begin{dictionary}
In \lean{cycDeg_finiteDimensional} the module $V$ and the action \lean{ρ : CurrentModule (bracketTable) V} are
section variables (\lean{variable {V : Type} ... {ρ : ...}} at the head of the section \lean{Fin}); they are
universally quantified, with no further hypothesis. \lean{FusionTotalDim} is the bare dimension statement; the theorem
\lean{fusionTotalBound_of_dim} supplies the finiteness half. The converse is trivial, so the two named statements are
equivalent. The bound is total (summed over all $b<B$, for every $B$), not graded.
\end{dictionary}

\textbf{Proof on paper.} (1) \emph{PBW sorting.} Call $e_1,e_2,e_\theta,h_1,h_2$ (at any grade) the \emph{non-lowering} letters.
Take a word applied to $w$ and look at its rightmost non-lowering letter $y(r)$, the one that acts first among them. Using
$y(r)x(s)=x(s)y(r)+[y,x](r+s)$, move it to the right past the lowering letters that act before it, until it reaches $w$; there it
kills $w$ or multiplies it by a scalar. Each commutator term is a word of the same degree with one fewer letter. Induction on
length (and on the number of non-lowering letters) shows: \emph{every $\ltwentysixcyc(w)_b$ is spanned by words in the lowering
letters $f_1(r),f_2(r),f_\theta(r)$ only}, and, sorting further, by such words in a fixed order with the grade-zero letters acting
first. (2) \emph{Local nilpotency.} Each $f_\alpha(0)$ satisfies $(\operatorname{ad}f_\alpha(0))^3=0$ on every letter
(\lean{ad_cube}: the root strings of $\mathfrak{sl}_3$ have length at most three), and $f_\alpha(0)^{2m+1}w=0$ for each
negative root (for $f_\theta$ via the Heisenberg lemma of Section~\ref{l26:sec:garland}). Kac's lemma (stated and proved as
Theorem~\ref{l26:thm:finiteG}(1) below: if $(\operatorname{ad}A)^KB=0$ and $A^Nu=0$ then $A^{N+K}Bu=0$) then shows that
$f_\alpha(0)^{2\ell+2m+1}$ kills any word of length $\ell$ applied to $w$. (3) \emph{Length bound.} A non-vanishing sorted
word of degree $b$ has at most $b$ letters of positive grade, and between them runs of grade-zero letters whose length is bounded
by (2); hence its length is bounded by an explicit number $G(m,b)$ (Lean: \lean{Gr m b 8}, a closed-form function of $m$, $b$ and
the alphabet size $8$). (4) There are finitely many words of bounded length with letters of grade $\le b$ in an alphabet of eight
letters, so $\ltwentysixcyc(w)_b$ lies in a finite span. \qed

\begin{physics}
(Heuristic.) In a fusion of $m$ adjoint particles there are only finitely many states at each energy: the zero modes act as a
finite-dimensional representation of $\ltwentysixg$ (integrability), and each excitation costs energy.
\end{physics}

\begin{status}
\lean{cycDeg_finiteDimensional} and \lean{fusionTotalBound_of_dim}: proved, \texttt{\#print axioms} $=$
\lean{[propext, Classical.choice, Quot.sound]} (re-run, Section~\ref{l26:sec:verification}).
\lean{FusionTotalDim}: a named \lean{Prop} here; it is proved as \lean{CVUniv.fusionTotalDim} (Theorem~\ref{l26:thm:cvtotal}).
\end{status}

%======================================================================
\section{U6 by construction: the fusion characters}\label{l26:sec:u6}
%======================================================================

\begin{definition}[The fusion characters as the inverse of the mixed recursion]\label{l26:def:fusionOf}
For any sequence of power series $\mathrm{Dem}_n$ define recursively
\begin{equation*}
\mathrm{fusionOf}(\mathrm{Dem})_n=\mathrm{Dem}_n-\sum_{d=1}^{n}q^{dn}\ltwentysixqb nd\,\mathrm{fusionOf}(\mathrm{Dem})_{n-d}.
\end{equation*}
For $A_2$, $\mathrm{fusionF}=\mathrm{fusionOf}(\ltwentysixD)$ with $\ltwentysixD_n$ the Gelfand--Tsetlin polynomial
\lean{demazureHilbSeries n}.
\end{definition}
\noindent File \leanpath{VOA/\allowbreak DeligneAudit/\allowbreak D4FusionF.lean}.
\begin{leancode}
def fusionOf (Dem : ℕ → PowerSeries ℚ) : ℕ → PowerSeries ℚ
  | n => Dem n -
      ∑ d : Fin n, X ^ ((d.1 + 1) * n) * qbinom (X : PowerSeries ℚ) n (d.1 + 1) *
        fusionOf Dem (n - (d.1 + 1))
decreasing_by omega

theorem mixedFusionRecursion_fusionOf (Dem : ℕ → PowerSeries ℚ) :
    MixedFusionRecursion (fusionOf Dem) Dem := by

theorem fusionOf_unique {F Dem : ℕ → PowerSeries ℚ} (h : MixedFusionRecursion F Dem) (n : ℕ) :
    F n = fusionOf Dem n := by
\end{leancode}
\noindent File \leanpath{VOA/\allowbreak DeligneAudit/\allowbreak DeligneFusionInverse.lean}.
\begin{leancode}
def fusionF : ℕ → PowerSeries ℚ
  | n => demazureHilbSeries n -
      ∑ d : Fin n, X ^ ((d.1 + 1) * n) * qbinom (X : PowerSeries ℚ) n (d.1 + 1) * fusionF (n - (d.1 + 1))
decreasing_by omega

theorem mixedFusionRecursion_fusionF : MixedFusionRecursion fusionF demazureHilbSeries := by
\end{leancode}

\begin{dictionary}
\lean{X} is the formal variable $q$; \lean{qbinom (X : PowerSeries ℚ) n k} is the Gaussian binomial
$\ltwentysixqb nk$; the sum over \lean{d : Fin n} with \lean{d.1 + 1} runs over $d=1,\dots,n$, so the $d=0$ term (which is
$\mathrm{fusionOf}_n$ itself) is moved to the left. \lean{decreasing_by omega} is Lean's termination proof.\lean{MixedFusionRecursion F Dem} (Definition~\ref{l24:def:chaininputs}) says $\mathrm{Dem}_n=\mathrm{mixed}(F)_{n,0}$,
where $\mathrm{mixed}(F)_{0,b}=F_b$ and $\mathrm{mixed}(F)_{a+1,b}=\mathrm{mixed}(F)_{a,b+1}+q^{2a+b+1}\mathrm{mixed}(F)_{a,b}$;
the library's lemma \lean{mixed_eq} (\leanpath{GaussianTransform.lean}, Chapter~\ref{ch:L24}) evaluates this as
$\sum_dq^{dn}\ltwentysixqb nd F_{n-d}$, which is \eqref{l26:eq:mixed}.

\textbf{What ``U6 by construction'' does and does not say.} \lean{mixedFusionRecursion_fusionOf} is true for any
\lean{Dem} whatsoever; it carries no information about fusion products. The mathematical content of Ravinder's Cor.~1,
``the true graded character of $F_m$ satisfies \eqref{l26:eq:mixed}'', is recovered only when U5 is proved for
\lean{fusionOf Dem} (Theorem~\ref{l26:thm:squeeze}): then $\mathrm{ch}\,F_m\le\mathrm{fusionOf}(\mathrm{Dem})_m$, and the squeeze
gives equality. By \lean{fusionOf_unique}, nothing is lost by choosing this $F$: any $F$ that satisfies U6 is
\lean{fusionOf Dem}.
\end{dictionary}

\begin{warning}
\textbf{Reading the fusion characters.} \lean{fusionF} and \lean{fusionOf Dem} are defined by inverting the mixed recursion. The
docstring statement that \lean{fusionF m} \emph{is} Ravinder's $f_m$ (the graded character of $V(\theta)^{*m}$) is a literature
claim plus numerics ($m\le5$), not a Lean fact. In Lean the identification follows only once U5 is proved, by the squeeze. Do not
take \lean{mixedFusionRecursion_fusionOf} as evidence about fusion products.
\end{warning}

\begin{example}[$\mathfrak{sl}_3$, $n\le2$]\label{l26:ex:fusionF}
$\mathrm{fusionF}_0=\ltwentysixD_0=1$; $\mathrm{fusionF}_1=\ltwentysixD_1-q\ltwentysixqb11\mathrm{fusionF}_0=8+q-q=8$;
$\mathrm{fusionF}_2=\ltwentysixD_2-q^2(1+q)\cdot8-q^4=27+28q+9q^2$. At $q=1$ these are $1,8,64$.
\end{example}

\textbf{Proof on paper.} Expanding $\mathrm{mixed}(F)_{n,0}$ by \lean{mixed_eq} and splitting off the $d=0$ term gives
$F_n+\sum_{d\ge1}q^{dn}\ltwentysixqb nd F_{n-d}$; with $F=\mathrm{fusionOf}(\mathrm{Dem})$ this is
$\mathrm{Dem}_n$ by definition. Uniqueness: if two sequences satisfy the recursion, their difference
$\delta$ satisfies $\delta_n=-\sum_{d\ge1}(\dots)\delta_{n-d}$ and $\delta_0=0$, so $\delta=0$ by strong induction. \qed

\begin{physics}
(Heuristic.) A ``bootstrap'' of the partition functions: once the degree-graded counts of the Demazure states are known,
the recursion determines the counts of the fusion states uniquely. It is the generating-function form of the statement
that states with $n$ particles decompose into $d$ bound particles and $n-d$ free ones.
\end{physics}

\begin{status}
Proved; standard axioms for \lean{mixedFusionRecursion_fusionOf} (re-run) and \lean{mixedFusionRecursion_fusionF} (record). The per-type series are \lean{fusionF} ($A_2$), \lean{d4FusionF},
\lean{g2FusionF}, \lean{fusionOf (hilbSeries e8DemazureHilb)}, etc.
\end{status}

\begin{theorem}[$A_2$ classical freeness from U5 alone]\label{l26:thm:a2u5}
Let $W=L_{-3/2}(\mathfrak{sl}_3)$ with its explicit presentation (Chapter~\ref{ch:L24}). If
\lean{SectorWeyl.FusionModuleBound fusionF} holds, then $W$ is classically free for the presentation and the kernel of
the current-coordinate map equals $(I_2^{\mathrm{ns}})$.
\end{theorem}
\noindent File \leanpath{VOA/\allowbreak DeligneAudit/\allowbreak DeligneFusionInverse.lean}.
\begin{leancode}
theorem classicallyFree_of_fusionBound (hU5 : SectorWeyl.FusionModuleBound fusionF) :
    CaseClassicallyFree .A2 W explicitPresentation ∧
      RingHom.ker (MvPolynomial.aeval fun i =>
        liDegreeZeroClass W 1 ⟨explicitPresentation.gen i,
          explicitPresentation.gen_weight_one i⟩).toRingHom = I2ns :=
\end{leancode}

\begin{dictionary}
\lean{CaseClassicallyFree .A2 W explicitPresentation} and the kernel identity are the conclusion of the capstone of
Chapter~\ref{ch:L24} (Definition~\ref{l24:def:lpcf} and Theorem~\ref{l24:thm:capstone}); \lean{I2ns} is the ideal of the
eight quadrics. The proof is \lean{classicallyFree_of_fusion fusionF hU5 mixedFusionRecursion_fusionF}, i.e.\ the
$A_2$ chain of Chapter~\ref{ch:L24} with U6 discharged. The only hypothesis is U5 for the explicit series \lean{fusionF}.
\end{dictionary}

\textbf{Proof on paper.} The chain of Chapter~\ref{ch:L24} (Theorem~\ref{l24:thm:seven} and the $A_2$ capstone) takes U5 and
U6 for a family $F$; U6 for $F=\mathrm{fusionF}$ is Definition~\ref{l26:def:fusionOf}. \qed

\begin{physics}
$L_{-3/2}(\mathfrak{sl}_3)$ is, by the four-dimensional/two-dimensional correspondence of Beem et al.\ (a physics input,
arXiv:1312.5344), the chiral algebra of the Argyres--Douglas theory $H_2$. Its classical freeness follows once the multi-particle
fusion spaces are known to have the expected size.
\end{physics}

\begin{status}
Standard axioms (record), \textbf{conditional on} \lean{SectorWeyl.FusionModuleBound fusionF}. The rest of
the chapter reduces that hypothesis.
\end{status}

%======================================================================
\section{The squeeze}\label{l26:sec:squeeze}
%======================================================================

\begin{definition}[The squeeze inputs and U7]\label{l26:def:squeezeinputs}
\begin{enumerate}
\item[(a)] For $f,g\in\Q[[q]]$ write $f\le g$ (\emph{coefficientwise order}) if $[q^N]f\le[q^N]g$ for all $N$.
\item[(b)] A family $M_{a,b}\in\Q[[q]]$ satisfies the \emph{squeeze inputs} if (nonneg) its coefficients are $\ge0$;
(R) $M_{a+1,b}\le M_{a,b+1}+q^{2a+b+1}M_{a,b}$; (D) $\ltwentysixD_n\le M_{n,0}$; (T) $\sum_{N<B}[q^N]M_{0,m}\le8^m$ for all
$m,B$. These are the three inequalities of Section~\ref{l26:sec:squeeze-paper}.
\item[(c)] It satisfies the \emph{universal bound} if every module whose generator satisfies the relations of $M_{a,b}$ has
$\dim\ltwentysixcyc(w)_N\le[q^N]M_{a,b}$.
\item[(d)] \textbf{U7 at $n$} (\lean{DemazureRealization n}): there is a current module and a vector $w$ satisfying the
relations of $M_{n,0}$ with $\dim\ltwentysixcyc(w)_N\ge[q^N]\ltwentysixD_n$ for every $N$.
\end{enumerate}
\end{definition}
\noindent File \leanpath{VOA/\allowbreak DeligneAudit/\allowbreak DeligneFusionSqueeze.lean}.
\begin{leancode}
def CoeffLe (f g : PowerSeries ℚ) : Prop := ∀ N, coeff N f ≤ coeff N g

structure FusionSqueezeInputs (M : ℕ → ℕ → PowerSeries ℚ) : Prop where
  /-- Graded characters have nonnegative coefficients. -/
  nonneg : ∀ a b N, 0 ≤ coeff N (M a b)
  /-- Spanning half of Ravinder's exact sequences. -/
  rec_le : ∀ a b, CoeffLe (M (a + 1) b) (M a (b + 1) + X ^ (2 * a + b + 1) * M a b)
  /-- `Hilb_q D(1,nθ) ≤ ch W_loc(nθ)`, with U7 (Demazure = Gelfand–Tsetlin). -/
  dem_le : ∀ n, CoeffLe (demazureHilbSeries n) (M n 0)
  /-- F4: `dim M_{0,m} ≤ 8^m`. -/
  total_le : ∀ m B, ∑ N ∈ Finset.range B, coeff N (M 0 m) ≤ 8 ^ m
\end{leancode}
\noindent File \leanpath{VOA/\allowbreak DeligneAudit/\allowbreak DeligneFusionDemazure.lean}.
\begin{leancode}
def MabUniversalBound (M : ℕ → ℕ → PowerSeries ℚ) : Prop :=
  ∀ (a b : ℕ) (V : Type) [AddCommGroup V] [Module ℂ V] (ρ : CurrentModule (bracketTable) V)
    (w : V), MabRel ρ w a b → ∀ N, (Module.finrank ℂ (cycDeg ρ w N) : ℚ) ≤ coeff N (M a b)

def DemazureRealization (n : ℕ) : Prop :=
  ∃ (V : Type) (_ : AddCommGroup V) (_ : Module ℂ V) (ρ : CurrentModule (bracketTable) V) (w : V),
    MabRel ρ w n 0 ∧ ∀ N, demazureHilb n N ≤ Module.finrank ℂ (cycDeg ρ w N)
\end{leancode}

\begin{dictionary}
\lean{coeff N} is the $q^N$-coefficient. In \lean{FusionSqueezeInputs} the field \lean{dem_le} uses
\lean{demazureHilbSeries n}, the Gelfand--Tsetlin polynomial of Chapter~\ref{ch:L25} viewed as a power series; that it
equals $\ltwentysixHilb_qD(1,n\theta)$ is a literature identification (Sanderson, Ion) checked numerically, not a Lean
fact, but nothing below needs it: what is used is only the inequality (D) for that polynomial. \lean{total_le} sums
over $N<B$ for every $B$, i.e.\ it bounds the total. Its docstring label ``F4'' is the project's name for the total bound,
(TB) in this chapter; it has nothing to do with the Lie type $F_4$.

\textbf{The \lean{finrank} convention in \lean{MabUniversalBound}.} As for \lean{FusionTotalBound}
(Definition~\ref{l26:def:mab}), \lean{Module.finrank} is $0$ on an infinite-dimensional space, so \lean{MabUniversalBound}
alone says nothing about a module with an infinite-dimensional degree piece. Where the squeeze uses it as a genuine bound,
finiteness of the universal pieces is supplied separately by \lean{univ_fd} (Theorem~\ref{l26:thm:a2assembly}); the
universal property then bounds every model by the universal one. In \lean{DemazureRealization}, \lean{demazureHilb n N : ℕ}
is the $q^N$-coefficient of the Gelfand--Tsetlin polynomial \lean{gl3Hilb (2 * n) n 0} of Chapter~\ref{ch:L25}, i.e.\ of
\lean{demazureHilbSeries n}; it is a lower bound on \lean{finrank}, so a model with an infinite-dimensional piece (where
\lean{finrank} is $0$) cannot satisfy it vacuously unless \lean{demazureHilb n N} $=0$.

\lean{DemazureRealization n} asks for \emph{some} model, with a lower bound in every degree. It is an existence
statement, so it is not vacuous in either direction: it would fail if no module with the $M_{n,0}$ relations were that
large. The universal bound turns it into (D) (Theorem~\ref{l26:thm:squeeze}).

\textbf{A stale docstring.} The docstring of \lean{DemazureRealization} says that neither the presentation of $D(1,n\theta)$
nor its character ``is formalized''. That is accurate, but it suggests the \lean{Prop} is open; in fact
\lean{EvalModel.demazureRealization_all} (Theorem~\ref{l26:thm:u7}) proves \lean{DemazureRealization N} for every $N$ without
either fact.
\end{dictionary}

\begin{physics}
(Heuristic.) (R) is the ``at most'' half of the particle decomposition; (D) says there are at least as many states as in the
level-one Demazure space of a chiral boson; (T) caps the number of states of $m$ fused adjoint particles at the number
of states of $m$ free ones. The universal bound expresses that the module defined by generators and relations is the
largest Hilbert space compatible with those null vectors.
\end{physics}

\begin{status}
Definitions. \lean{FusionSqueezeInputs} and \lean{MabUniversalBound} are structures of hypotheses, discharged below for
the universal characters \lean{Mchar}; \lean{DemazureRealization n} is proved for all $n$ (Theorem~\ref{l26:thm:u7}).
\end{status}

\begin{theorem}[Ravinder's squeeze; U5 from the squeeze]\label{l26:thm:squeeze}
If $M$ satisfies the squeeze inputs, then $M_{0,m}=\mathrm{fusionF}_m$ for every $m$, and $\mathrm{fusionF}$ has nonnegative
coefficients. If moreover every module with the local Weyl and fusion relations for $m\theta$ is bounded degreewise by
$M_{0,m}$, then U5 holds for $\mathrm{fusionF}$. Finally U5 for $\mathrm{fusionF}$ follows from the universal bound, U7 for
all $n$, (nonneg), (R) and (T).
\end{theorem}
\noindent File \leanpath{VOA/\allowbreak DeligneAudit/\allowbreak DeligneFusionSqueeze.lean}.
\begin{leancode}
theorem squeeze (h : FusionSqueezeInputs M) (m : ℕ) : M 0 m = fusionF m := by

theorem fusionF_coeff_nonneg (h : FusionSqueezeInputs M) (m N : ℕ) :
    0 ≤ coeff N (fusionF m) := by

theorem fusionModuleBound_of_squeeze (h : FusionSqueezeInputs M)
    (huniv : ∀ (m : ℕ) (V : Type) [AddCommGroup V] [Module ℂ V]
      (ρ : CurrentModule (bracketTable) V) (w : V), LocalWeylRel ρ w m →
        ∀ b, (Module.finrank ℂ (cycDeg ρ w b) : ℚ) ≤ coeff b (M 0 m)) :
    FusionModuleBound fusionF := by
\end{leancode}
\noindent File \leanpath{VOA/\allowbreak DeligneAudit/\allowbreak DeligneFusionDemazure.lean}.
\begin{leancode}
theorem dem_le_of_realization (hU : MabUniversalBound M) {n : ℕ} (hD : DemazureRealization n) :
    CoeffLe (demazureHilbSeries n) (M n 0) := by

theorem fusionModuleBound_of_realizations (hU : MabUniversalBound M)
    (hD : ∀ n, DemazureRealization n) (hnonneg : ∀ a b N, 0 ≤ coeff N (M a b))
    (hrec : ∀ a b, CoeffLe (M (a + 1) b) (M a (b + 1) + X ^ (2 * a + b + 1) * M a b))
    (htot : ∀ m B, ∑ N ∈ Finset.range B, coeff N (M 0 m) ≤ 8 ^ m) :
    FusionModuleBound fusionF :=
\end{leancode}

\begin{dictionary}
\lean{M : ℕ → ℕ → PowerSeries ℚ} is a section variable (\lean{variable {M : ...}}) in both files: it is universally
quantified, and the hypotheses on it are exactly those printed. The conclusion of \lean{squeeze} is an equality of
power series, i.e.\ equality in every degree. \lean{FusionModuleBound fusionF} is the $A_2$ form of U5
(Definition~\ref{l24:def:U5}): a graded bound \lean{le} and a total bound \lean{total}. The hypothesis \lean{huniv} of
\lean{fusionModuleBound_of_squeeze} says every module with \lean{LocalWeylRel ρ w m} is bounded by $M_{0,m}$ in every
degree; it is the universal bound at $a=0$ (\lean{huniv_of_universalBound}, same file).
\end{dictionary}

\textbf{Proof on paper.} Section~\ref{l26:sec:squeeze-paper} gives the argument; we add the two points the Lean proof
must handle. First, \lean{le_mixed}: iterating (R) gives $M_{a,b}\le\mathrm{mixed}(M_{0,\cdot})_{a,b}$, by induction on
$a$, using that $\le$ is preserved by sums and by multiplication with $q^k$ (both because the comparison is
coefficientwise and $q^k$ only shifts). Second, the passage from ``$\mathrm{fusionF}_n\le M_{0,n}$ coefficientwise and
equal totals'' to equality: choose $B$ beyond the (finite) support of $\mathrm{fusionF}_n$ and beyond the degree $N$ of
interest; then $\sum_{k<B}([q^k]M_{0,n}-[q^k]\mathrm{fusionF}_n)\le8^n-8^n=0$ with nonnegative summands, so every
summand vanishes. The totals of $\mathrm{fusionF}_n$ are $8^n$ by \lean{fusionF_sum} (from \lean{fusionPoly_eval_one}:
the polynomial version evaluated at $q=1$ inverts $9^n=\sum_d\binom nd8^{n-d}$). For U5: the graded half is
$\dim\ltwentysixcyc(w)_b\le[q^b]M_{0,m}=[q^b]\mathrm{fusionF}_m$; the total half sums this and uses nonnegativity of
$\mathrm{fusionF}$ to extend the range of summation. \qed

\begin{physics}
(Heuristic.) A pure counting argument about a partition function: a degree-by-degree upper bound on the space of
states, a degree-by-degree lower bound on the Demazure side, and a bound on the total number of states together determine the
partition function exactly. No dynamics enters.
\end{physics}

\begin{status}
Proved; standard axioms (\lean{squeeze} re-run; the others record). No hypothesis beyond those printed.
\end{status}

%======================================================================
\section{Universal modules and the exact-sequence step}\label{l26:sec:univ}
%======================================================================

\begin{definition}[The universal modules $U_{a,b}$ and their characters]\label{l26:def:umod}
Let $\mathrm{FreeW}$ be the vector space with basis the loop words (lists of letters $(c,r)\in\{0,\dots,7\}\times\N$).
Let $J_{a,b}\subseteq\mathrm{FreeW}$ be spanned by: all \emph{commutation instances}
$p\,xy\,s-p\,yx\,s-\sum_ef(x,y,e)\,p\,(e,r_x+r_y)\,s$ (for prefixes $p$, suffixes $s$); all words $p\,t$ with $t$ a
relation tail of $M_{a,b}$ (a single raising letter, a positive-grade Cartan letter, $f_i(0)^{a+b+1}$, $f_1(a+b)$,
$f_2(a+b)$, $f_\theta(2a+b)$); and all $p\,h(0)-(a+b)\,p$. Then $U_{a,b}=\mathrm{FreeW}/J_{a,b}$, with $x_c(r)$
acting by prepending the letter, and generator $w_U=[\,\text{empty word}\,]$. Its character is
$M_{a,b}^{\mathrm{univ}}(q)=\sum_N\dim\ltwentysixcyc(w_U)_N\,q^N$ (\lean{Mchar a b}).
\end{definition}
\noindent File \leanpath{VOA/\allowbreak DeligneAudit/\allowbreak DeligneFusionUniv.lean}.
\begin{leancode}
abbrev Ltr := Fin 8 × ℕ

abbrev FreeW := List Ltr →₀ ℂ

def commInst (p s : List Ltr) (x y : Ltr) : FreeW :=
  sw (p ++ x :: y :: s) - sw (p ++ y :: x :: s) -
    ∑ e, bracketTable x.1 y.1 e • sw (p ++ (e, x.2 + y.2) :: s)

def relTail (a b : ℕ) : Set (List Ltr) :=
  {t | (∃ e r, e.val < 3 ∧ t = [(e, r)]) ∨ (∃ h r, 6 ≤ h.val ∧ 1 ≤ r ∧ t = [(h, r)]) ∨
    t = List.replicate (a + b + 1) (3, 0) ∨ t = List.replicate (a + b + 1) (4, 0) ∨
    t = [(3, a + b)] ∨ t = [(4, a + b)] ∨ t = [(5, 2 * a + b)]}

def relGen (a b : ℕ) : Set FreeW :=
  {g | ∃ p : List Ltr, (∃ s x y, g = commInst p s x y) ∨ (∃ t ∈ relTail a b, g = sw (p ++ t)) ∨
    (∃ h : Fin 8, 6 ≤ h.val ∧ g = sw (p ++ [(h, 0)]) - ((a + b : ℕ) : ℂ) • sw p)}

def Jrel (a b : ℕ) : Submodule ℂ FreeW := Submodule.span ℂ (relGen a b)

abbrev UMod (a b : ℕ) := FreeW ⧸ Jrel a b

def univRep (a b : ℕ) : CurrentModule (bracketTable) (UMod a b) where
  op c r := opU a b (c, r)
  comm c d r s v := by

def wU (a b : ℕ) : UMod a b := Submodule.Quotient.mk (sw [])

def Mchar (a b : ℕ) : PowerSeries ℚ :=
  PowerSeries.mk fun N => (Module.finrank ℂ (cycDeg (univRep a b) (wU a b) N) : ℚ)
\end{leancode}
\begin{leancode}
theorem univ_mabRel (a b : ℕ) : MabRel (univRep a b) (wU a b) a b where
\end{leancode}

\begin{dictionary}
\lean{List Ltr →₀ ℂ} is the space of finitely supported functions on words, i.e.\ formal linear combinations.
\lean{sw l} is the basis vector of the word \lean{l}; words are read with the head acting last, so prepending a letter is
applying an operator. \lean{relTail} lists exactly the tails of \lean{MabRel} (the indices $e<3$ are the raising
letters, $h\ge6$ the Cartan letters). \lean{univRep} is a \lean{CurrentModule}; its \lean{comm} field is proved from the
commutation instances with empty prefix (the proof body, elided here, is the only place where the quotient is used).
\lean{univ_mabRel} (statement line quoted; the proof follows in \lean{where} form) says the generator satisfies
\lean{MabRel}. The universal property (\lean{univMap}, \lean{finrank_cycDeg_le_univ}, same file) is: every module with
\lean{MabRel ρ w a b} receives a linear map from $U_{a,b}$ sending $[l]$ to \lean{wordOp ρ l w}, so
$\ltwentysixcyc(w)_N$ is the image of $\ltwentysixcyc(w_U)_N$.

\textbf{Why this construction and not an enveloping algebra.} The library has no $U(\ltwentysixgt)$; the free space on
words modulo prefixed relations is the cyclic module presented by generators and relations, built directly.
\end{dictionary}

\begin{physics}
(Heuristic.) $U_{a,b}$ is the ``Verma-like'' space: all states that can be created by current modes from a state obeying only the
stated null-vector conditions. Every physical realization is a quotient of it.
\end{physics}

\begin{status}
Definitions; \lean{univRep}, \lean{univ_mabRel}, \lean{finrank_cycDeg_le_univ}: \texttt{\#print axioms} standard (record).
\end{status}

\begin{theorem}[The exact-sequence step: quotient and rank--nullity]\label{l26:thm:step}
(i) If $w$ satisfies the relations of $M_{a+1,b}$ and $S$ is an invariant subspace containing $f_\theta(2a+b+1)w$, then
the image of $w$ in $V/S$ satisfies the relations of $M_{a,b+1}$. (ii) Assume the kernel relation (KL). Then for every
degree $N$, with $s=2a+b+1$,
\begin{equation*}
\dim\ltwentysixcyc_{U_{a+1,b}}(w_U)_N\ \le\ \dim\ltwentysixcyc_{U_{a,b+1}}(w_U)_N+[s\le N]\,\dim\ltwentysixcyc_{U_{a,b}}(w_U)_{N-s},
\end{equation*}
provided the two right-hand pieces are finite-dimensional, and then the left piece is finite-dimensional too.
\end{theorem}
\noindent File \leanpath{VOA/\allowbreak DeligneAudit/\allowbreak DeligneFusionGrade.lean}.
\begin{leancode}
theorem mabRel_quot {w : V} {a b : ℕ} (h : MabRel ρ w (a + 1) b)
    (hv : ρ.op 5 (2 * a + b + 1) w ∈ S) : MabRel (quotRep ρ S hS) (S.mkQ w) a (b + 1) where
\end{leancode}
\begin{leancode}
def vK : UMod (a + 1) b := (univRep (a + 1) b).op 5 (2 * a + b + 1) (wU (a + 1) b)

theorem fd_step (hR : RavinderKernelRel) (N : ℕ)
    [FiniteDimensional ℂ (cycDeg (univRep a (b + 1)) (wU a (b + 1)) N)]
    [FiniteDimensional ℂ (cycDeg (univRep a b) (wU a b) (N - (2 * a + b + 1)))] :
    FiniteDimensional ℂ (cycDeg (univRep (a + 1) b) (wU (a + 1) b) N) := by

theorem finrank_step (hR : RavinderKernelRel) (N : ℕ)
    [FiniteDimensional ℂ (cycDeg (univRep a (b + 1)) (wU a (b + 1)) N)]
    [FiniteDimensional ℂ (cycDeg (univRep a b) (wU a b) (N - (2 * a + b + 1)))] :
    Module.finrank ℂ (cycDeg (univRep (a + 1) b) (wU (a + 1) b) N) ≤
      Module.finrank ℂ (cycDeg (univRep a (b + 1)) (wU a (b + 1)) N) +
        if 2 * a + b + 1 ≤ N then
          Module.finrank ℂ (cycDeg (univRep a b) (wU a b) (N - (2 * a + b + 1))) else 0 := by
\end{leancode}

\begin{dictionary}
In \lean{mabRel_quot} the section variables are \lean{ρ}, a submodule \lean{S} and \lean{hS : ∀ c r, S ≤ S.comap (ρ.op c r)}
(invariance under every current operator); \lean{quotRep ρ S hS} is the induced current module on $V/S$ and \lean{S.mkQ}
the quotient map. In \lean{fd_step} and \lean{finrank_step}, \lean{a b} are section variables (\lean{variable (a b : ℕ)}),
and the bracketed \lean{[FiniteDimensional ...]} are instance hypotheses: finite-dimensionality of the two right-hand
pieces. The ``if'' encodes $[s\le N]$; \lean{N - (2 * a + b + 1)} is natural subtraction, harmless because it is used
only under $s\le N$.

\textbf{Which half of Ravinder's sequence.} This is the \emph{spanning} (upper bound) half: the kernel of
$U_{a+1,b}\to U_{a,b+1}$ is the cyclic module of $v_K$, and that module is a quotient of $U_{a,b}$ shifted by $s$. The
injectivity of $\tau_sM_{a,b}\to M_{a+1,b}$ is not claimed and not needed.
\end{dictionary}

\textbf{Proof on paper.} (i) All relations of $M_{a,b+1}$ except $f_\theta(2a+b+1)w=0$ are among those of $M_{a+1,b}$
(same $k=a+b+1$, same raising, Cartan, lowering and $f_1,f_2$ truncations); the remaining one holds in $V/S$ because
$f_\theta(2a+b+1)w\in S$. (ii) Let $S_K$ be the cyclic submodule of $U_{a+1,b}$ generated by $v_K$. By (i) the quotient
$U_{a+1,b}/S_K$ is a module with the $M_{a,b+1}$ relations, hence (universal property) a quotient of $U_{a,b+1}$, degree
by degree. By (KL), $v_K$ satisfies the $M_{a,b}$ relations, so $S_K$ is the image of $U_{a,b}$ under $[l]\mapsto lv_K$,
which raises degrees by $s$; hence the degree-$N$ piece of the cyclic module of $w_U$ intersected with $S_K$ is the image of
$\ltwentysixcyc_{U_{a,b}}(w_U)_{N-s}$ (using a degree projection on $U_{a+1,b}$, \lean{projU}, which exists because all
relations are homogeneous). Rank--nullity for the quotient map restricted to $\ltwentysixcyc(w_U)_N$ gives the inequality. \qed

\begin{physics}
(Heuristic mnemonic.) States of $a+1$ Demazure particles and $b$ adjoint ones are at most those in which one Demazure particle has
been replaced by an adjoint one, plus a copy of the smaller system shifted up by the ``binding energy'' $2a+b+1$, which is the
grading shift of \eqref{l26:eq:ses}.
\end{physics}

\begin{status}
Proved, standard axioms (record). Hypothesis of (ii): \lean{RavinderKernelRel}, which is a theorem
(Theorem~\ref{l26:thm:kernelA2}).
\end{status}

\begin{theorem}[The $A_2$ assembly]\label{l26:thm:a2assembly}
(1) \emph{Structural statements.} Assume (TB) and (KL). Then all degree pieces of all $U_{a,b}$ are finite-dimensional; if
moreover U7 holds for every $n$, the universal characters \lean{Mchar} satisfy the squeeze inputs.
(2) \emph{Corollary.} If \lean{FusionTotalDim} holds and U7 holds for every $n$, then $L_{-3/2}(\mathfrak{sl}_3)$ is classically
free with kernel $(I_2^{\mathrm{ns}})$. ((KL) is a theorem, Section~\ref{l26:sec:garland}, and (TB) follows from
\lean{FusionTotalDim} by Theorem~\ref{l26:thm:finite}, so neither appears as a hypothesis of (2).)
\end{theorem}
\noindent File \leanpath{VOA/\allowbreak DeligneAudit/\allowbreak DeligneFusionAssembly.lean}.
\begin{leancode}
theorem univ_fd (hF : FusionTotalBound) (hR : RavinderKernelRel) :
    ∀ a b N, FiniteDimensional ℂ (cycDeg (univRep a b) (SectorWeyl.wU a b) N)

theorem fusionSqueezeInputs_Mchar (hF : FusionTotalBound) (hR : RavinderKernelRel)
    (hD : ∀ n, DemazureRealization n) : FusionSqueezeInputs Mchar where

theorem classicallyFree_of_totalBound (hF : FusionTotalBound) (hD : ∀ n, DemazureRealization n) :
    CaseClassicallyFree .A2 W explicitPresentation ∧
      RingHom.ker (MvPolynomial.aeval fun i =>
        liDegreeZeroClass W 1 ⟨explicitPresentation.gen i,
          explicitPresentation.gen_weight_one i⟩).toRingHom = I2ns :=

theorem classicallyFree_of_totalDim (hF : FusionTotalDim) (hD : ∀ n, DemazureRealization n) :
    CaseClassicallyFree .A2 W explicitPresentation ∧
      RingHom.ker (MvPolynomial.aeval fun i =>
        liDegreeZeroClass W 1 ⟨explicitPresentation.gen i,
          explicitPresentation.gen_weight_one i⟩).toRingHom = I2ns :=
\end{leancode}

\begin{dictionary}
\lean{univ_fd} is by induction on $a$: at $a=0$ finiteness is the first half of \lean{FusionTotalBound}, applied to the
universal module (whose generator satisfies \lean{LocalWeylRel} by \lean{univ_mabRel} and \lean{mabRel_zero_iff}); the
step is \lean{fd_step}. In \lean{fusionSqueezeInputs_Mchar}: \lean{nonneg} is free (dimensions); \lean{rec_le} is
\lean{finrank_step}; \lean{dem_le} is \lean{dem_le_of_realization}; \lean{total_le} is the second half of
\lean{FusionTotalBound} for the universal module. In the two \lean{classicallyFree} theorems, the kernel lemma has
disappeared from the hypotheses because \lean{ravinderKernelTrunc} is a theorem (next section), and the finiteness clause
because of Theorem~\ref{l26:thm:finite}. The file's docstring still lists \lean{RavinderKernelTrunc} as a named input; that
is out of date (see Section~\ref{l26:sec:defects}).
\end{dictionary}

\textbf{Proof on paper.} Combine Theorems~\ref{l26:thm:step}, \ref{l26:thm:squeeze} and \ref{l26:thm:a2u5}. \qed

\begin{physics}
(Heuristic.) The $A_2$ chiral algebra's classical freeness rests on one counting statement about fused adjoint particles and
one statement about the Demazure space, the latter proved in Section~\ref{l26:sec:u7}.
\end{physics}

\begin{status}
Standard axioms (record), \textbf{conditional on} \lean{FusionTotalDim} and \lean{∀ n, DemazureRealization n} as stated.
Both are now theorems: the second is discharged in \lean{classicallyFree_of_totalDim'} (Theorem~\ref{l26:thm:u7final}) and the
first by \lean{CVUniv.fusionTotalDim} (Theorem~\ref{l26:thm:cvtotal}); the composite is the unconditional theorem of
Chapter~\ref{ch:L27}.
\end{status}

%======================================================================
\section{Garland's identities and Ravinder's kernel lemma for $A_2$}\label{l26:sec:garland}
%======================================================================

The kernel lemma needs three truncations of the kernel vector $v_K=f_\theta(2a+b-1)w$. Their source is a classical
identity of Garland, in the form used by Chari--Pressley: in an $\mathfrak{sl}_2[t]$-module, if $w$ is killed by $e(1)$
and by $h(r)$, $r\ge1$, and $f(0)^Nw=0$, then $[u^s]f(u)^r\,w=0$ whenever $r\ge1$ and $r+s\ge N$, where
$f(u)=\sum_kf(k)u^k$. The library proves this special case, in which every $h(r)$, $r\ge1$, kills $w$; it is all that is
needed.

\begin{theorem}[Heisenberg nilpotency; Garland's identity with $P(u)=1$]\label{l26:thm:garland}
(1) Let $X,Y,Z$ be operators with $[X,Y]=Z$ and $Z$ commuting with $X$ and $Y$. If $X^{p+1}w=Y^{p+1}w=0$ then
$Z^{2p+1}w=0$. (2) Let $E_1$, $H_k$ ($k\ge1$), $F_k$ ($k\ge0$) be operators with $[E_1,F_k]=H_{k+1}$,
$[H_{m+1},F_k]=-2F_{m+1+k}$ and $[F_j,F_k]=0$, and let $E_1w=0$, $H_{m+1}w=0$ for all $m$. If $F_0^Nw=0$, then for
$r\ge1$ and $r+s\ge N$,
\begin{equation*}
[u^s]\,F(u)^r\,w=\sum_{k_1+\dots+k_r=s}F_{k_1}\cdots F_{k_r}\,w=0 .
\end{equation*}
\end{theorem}
\noindent File \leanpath{VOA/\allowbreak DeligneAudit/\allowbreak DeligneFusionGarland.lean}.
\begin{leancode}
variable {X Y Z : Module.End ℂ V} (hXY : X * Y - Y * X = Z) (hXZ : X * Z = Z * X)
  (hYZ : Y * Z = Z * Y)

theorem heis_nilpotent (hXZ : X * Z = Z * X) {p : ℕ} {w : V} (hX : (X ^ (p + 1)) w = 0)
    (hY : (Y ^ (p + 1)) w = 0) : (Z ^ (2 * p + 1)) w = 0 := by

def Fs : PowerSeries (Module.End ℂ V) := PowerSeries.mk F


variable {E1 H F} (hEF : ∀ k, E1 * F k - F k * E1 = H (k + 1))
  (hHF : ∀ m k, H (m + 1) * F k - F k * H (m + 1) = (-2 : ℂ) • F (m + 1 + k))

theorem garland (r : ℕ) (hr : 1 ≤ r) {N : ℕ} (hN : (F 0 ^ N) w = 0) (s : ℕ) (hs : N ≤ r + s) :
    (coeff s (Fs F ^ r) : Module.End ℂ V) w = 0 := by
\end{leancode}

\begin{dictionary}
\textbf{Hidden section hypotheses.} In \lean{heis_nilpotent}, the section variables \lean{hXY : X * Y - Y * X = Z} and
\lean{hYZ : Y * Z = Z * Y} are \lean{include}d, so the full statement is
\lean{heis_nilpotent hXY hYZ hXZ hX hY}; the printed \lean{hXZ} is the third hypothesis. In \lean{garland} the five
hypotheses of the section \lean{Kill} (quoted above) are \lean{include}d: \lean{hEF} ($[E_1,F_k]=H_{k+1}$), \lean{hHF}
($[H_{m+1},F_k]=-2F_{m+1+k}$), \lean{hFF} ($F$'s commute), \lean{hE} ($E_1w=0$), \lean{hH} ($H_{m+1}w=0$); the call
pattern is \lean{garland hEF hHF hFF hE hH r hr hN s hs}. \lean{Fs F} is $F(u)=\sum_kF_ku^k$ as a power series with
coefficients in \lean{Module.End ℂ V} (noncommutative coefficients are fine because the $F_k$ commute among themselves).
\lean{coeff s (Fs F ^ r)} is $[u^s]F(u)^r$.

\textbf{Which identity this is.} Garland's general identity allows the $h(k)$, $k\ge1$, to act on $w$ by arbitrary scalars,
which enter through a generating polynomial (``$P(u)$''). We use only the case in which all $h(k)$, $k\ge1$, kill $w$ ($P=1$),
and only the operator $E_1=e\otimes t$, not all $e\otimes t^k$. This is weaker than Garland's theorem but sufficient for every
module in this chapter, because in all $M_{a,b}$ the generator is killed by $h(k)$, $k\ge1$.
\end{dictionary}

\textbf{Proof on paper.} (1) One shows by induction $\operatorname{ad}_X^{\,i}(Y^{i+r})=\frac{(i+r)!}{r!}Z^iY^r$, so
$\operatorname{ad}_X^{\,p'}(Y^{p'})=p'!\,Z^{p'}$ with $p'=2p+1$. Expanding $\operatorname{ad}_X^{\,p'}(Y^{p'})w$ as
$\sum_j\pm\binom{p'}jX^{p'-j}Y^{p'}X^jw$, it suffices that $Y^{p'}X^jw=0$ for every $j$. If $j\ge p+1$ then $X^jw=0$. If
$j\le p$, commute $Y^{p'}$ past $X^j$: since $Z$ is central, $Y^{p'}X^j=\sum_{i\le j}c_i\,X^{j-i}Z^iY^{p'-i}$ for constants
$c_i$, and every term kills $w$ because $p'-i\ge 2p+1-p=p+1$ and $Y^{p+1}w=0$. Hence $p'!\,Z^{p'}w=0$, i.e.\
$Z^{2p+1}w=0$ (Lean: the auxiliary statement \lean{h0} inside \lean{heis_nilpotent}).

(2) Work with power series in $u$ whose coefficients are operators on $V$, and write $F(u)=\sum_kF_ku^k$,
$H(u)=\sum_kH_{k+1}u^k$ and $F'(u)=\sum_k(k+1)F_{k+1}u^k$. The relations say
\begin{equation*}
[E_1,F(u)]=H(u),\qquad [H(u),F(u)]=-2\sum_{m,k}F_{m+k+1}u^{m+k}=-2F'(u),\qquad H(u)w=0,\qquad E_1w=0,
\end{equation*}
and $F(u)$, $F'(u)$ commute. \emph{One step.} $E_1F(u)^2w=[E_1,F(u)^2]w=\bigl(H(u)F(u)+F(u)H(u)\bigr)w=[H(u),F(u)]w=-2F'(u)w$.
\emph{General step.} In the same way $E_1F(u)^nw=\sum_{i=0}^{n-1}F(u)^iH(u)F(u)^{n-1-i}w$, and $H(u)F(u)^jw=[H(u),F(u)^j]w=-2jF(u)^{j-1}F'(u)w$;
summing, $E_1F(u)^nw=-n(n-1)F(u)^{n-2}F'(u)w$. (Every equation here is an equality of the vectors obtained by applying a series to
$w$, coefficient by coefficient; that is all that ``modulo the series killing $w$'' means.) With $n=q+2$ and
$F^qF'=(F^{q+1})'/(q+1)$, comparing the coefficients of $u^j$ gives, for $Q(r,j)=[u^j]F(u)^rw$,
\begin{equation*}
E_1\,Q(q+2,j)=-(q+2)(j+1)\,Q(q+1,j+1).
\end{equation*}
Iterating $s$ times, $E_1^s\,Q(q+1+s,0)$ is a nonzero multiple of $Q(q+1,s)$. But $Q(q+1+s,0)=F_0^{q+1+s}w=0$ when $q+1+s\ge N$.
So $Q(r,s)=0$ for $r=q+1\ge1$ and $r+s\ge N$. \qed

\begin{physics}
$F(u)$ is the generating function of the modes of the current $f(z)$. For the level-one boson, the normally ordered square
$f(z)^2$ vanishes: $f$ is the vertex operator $e^{-\alpha}(z)$, $\langle\alpha,\alpha\rangle=2$, and the operator product
$e^{-\alpha}(z)e^{-\alpha}(w)$ has a double zero $(z-w)^2$ at coincident points. The current modes commute (they are bosonic); the
vanishing is a property of this operator product, and calling it an ``exclusion'' is only a mnemonic. Garland's identities are the
algebraic shadow of such nilpotency of currents on a highest-weight state.
\end{physics}

\begin{status}
Proved; standard axioms (\lean{garland} re-run).
\end{status}

\begin{theorem}[Ravinder's kernel lemma for $\mathfrak{sl}_3$]\label{l26:thm:kernelA2}
Let $w$ satisfy the relations of $M_{a,b}$. Then $f_\theta(0)^{2(a+b)+1}w=0$ and
$\sum_{i+j=s}f_\theta(i)f_\theta(j)w=0$ for $s+1\ge2(a+b)$. If $a\ge1$, the kernel vector $v=f_\theta(2a+b-1)w$ satisfies
$f_1(a-1+b)v=f_2(a-1+b)v=0$ and $f_\theta(2(a-1)+b)v=0$, and therefore all relations of $M_{a-1,b}$.
\end{theorem}
\noindent File \leanpath{VOA/\allowbreak DeligneAudit/\allowbreak DeligneFusionKernel.lean}.
\begin{leancode}
def RavinderKernelTrunc : Prop :=
  ∀ (a b : ℕ) (V : Type) [AddCommGroup V] [Module ℂ V] (ρ : CurrentModule (bracketTable) V) (w : V),
    1 ≤ a → MabRel ρ w a b →
      ρ.op 3 (a - 1 + b) (ρ.op 5 (2 * a + b - 1) w) = 0 ∧
      ρ.op 4 (a - 1 + b) (ρ.op 5 (2 * a + b - 1) w) = 0 ∧
      ρ.op 5 (2 * (a - 1) + b) (ρ.op 5 (2 * a + b - 1) w) = 0

theorem ravinderKernelRel_of_trunc (hT : RavinderKernelTrunc) : RavinderKernelRel := by
\end{leancode}
\noindent File \leanpath{VOA/\allowbreak DeligneAudit/\allowbreak DeligneFusionGarlandSl3.lean}.
\begin{leancode}
theorem ftheta_pow (h : MabRel ρ w a b) : ((ρ.op 5 0) ^ (2 * (a + b) + 1)) w = 0 := by

theorem theta_garland (h : MabRel ρ w a b) (s : ℕ) (hs : 2 * (a + b) ≤ s + 1) :
    ∑ x ∈ Finset.antidiagonal s, ρ.op 5 x.1 (ρ.op 5 x.2 w) = 0 := by

theorem ravinderKernelTrunc : RavinderKernelTrunc := by

theorem ravinderKernelRel : RavinderKernelRel := ravinderKernelRel_of_trunc ravinderKernelTrunc
\end{leancode}

\begin{dictionary}
In \lean{ftheta_pow} and \lean{theta_garland}, \lean{ρ}, \lean{w}, \lean{a}, \lean{b} are section variables; there are no
included hypotheses besides the printed \lean{h : MabRel ρ w a b}. \lean{(ρ.op 5 0) ^ n} is the $n$-th power in
$\End(V)$; \lean{Finset.antidiagonal s} is $\{(i,j):i+j=s\}$. \lean{RavinderKernelTrunc} is stated as a named \lean{Prop}
(its docstring still calls it ``named literature input''); \lean{ravinderKernelTrunc} proves it, so
\lean{RavinderKernelRel} (Definition~\ref{l26:def:mab}) is a theorem.
\end{dictionary}

\textbf{Proof on paper.} \emph{Integrability of $f_\theta$.} In $\mathfrak{sl}_3$, $[f_1,f_2]=-f_\theta$ and $f_\theta$
commutes with $f_1,f_2$: the three lowering vectors span a Heisenberg algebra. Since $f_1(0)^{k+1}w=f_2(0)^{k+1}w=0$,
Theorem~\ref{l26:thm:garland}(1) gives $f_\theta(0)^{2k+1}w=0$ with $k=a+b$. \emph{The $\theta$-Garland identity.}
Apply Theorem~\ref{l26:thm:garland}(2), $r=2$, to the $\mathfrak{sl}_2$-triple $(e_\theta(1),\,h_1+h_2,\,f_\theta)$ with
$N=2k+1$: $\sum_{i+j=s}f_\theta(i)f_\theta(j)w=0$ for $s+2\ge2k+1$. \emph{Extracting single terms.} At $s=2(2a+b)-3$, all
terms with an index $\ge2a+b$ vanish by the truncation $f_\theta(2a+b)w=0$ (and its consequences $f_\theta(j)w=0$ for
$j\ge2a+b$, which follow from $[h_1(1),f_\theta(r)]=-f_\theta(r+1)$); the only surviving pair is $\{2a+b-2,2a+b-1\}$, so
$2f_\theta(2a+b-2)f_\theta(2a+b-1)w=0$, which is the $f_\theta$-truncation of $v$. For $f_1$: apply $e_2(0)$ to the
$\theta$-Garland identity at $s=(a+b-1)+(2a+b-1)$; since $[e_2,f_\theta]$ is a multiple of $f_1$, each term becomes
$f_1(i)f_\theta(j)w+f_\theta(i)f_1(j)w$, and the truncations kill all but $f_1(a+b-1)f_\theta(2a+b-1)w$. Symmetrically for $f_2$
with $e_1(0)$. The other relations of $M_{a-1,b}$ for $v$ (raising, Cartan, lowering) follow from commutators with
$f_\theta$ and $\mathfrak{sl}_2$ representation theory (\lean{kernel_raising}, \dots, \lean{kernel_lowering_f2}). \qed

\begin{physics}
(Heuristic.) The kernel vector is the state in which one Demazure particle has emitted its top $f_\theta$ excitation. The lemma says
this state behaves exactly like the ground state of the smaller system: it is annihilated by the same set of modes.
\end{physics}

\begin{status}
Proved, unconditional; \texttt{\#print axioms} standard (\lean{ravinderKernelRel} re-run).
\end{status}

%======================================================================
\section[The Demazure side for $A_2$]{The Demazure side for $A_2$: Chari--Pressley, Chari--Loktev, Chari--Venkatesh, and U7}\label{l26:sec:u7}
%======================================================================

The Demazure lower bound (D) needs a module with the $M_{N,0}$ relations whose degree pieces are at least
$\ltwentysixD_N$. The library proves this through three classical results about current algebras, formalized here for
$\mathfrak{sl}_2$ and $\mathfrak{sl}_3$, and one construction.

\subsection{Chari--Pressley: box spanning for $\mathfrak{sl}_2$}

\begin{theorem}[Chari--Pressley, Prop.~6.1]\label{l26:thm:box}
In $R=\C[z_0,z_1,\dots]$ let $Z(u)=\sum_iz_iu^i$, $x(r,s)=[u^s]Z(u)^r$ and $J_m=(x(r,s):r\ge1,\ r+s\ge m+1)$. Then every
$p\in R$ is a combination of \emph{box monomials} $z_{i_1}\cdots z_{i_\ell}$ with $i_j+\ell\le m$, plus an element of $J_m$.
Consequently, in every $\mathfrak{sl}_2[t]$-module, if $w$ is killed by $e\otimes t$ and by $h(k)$, $k\ge1$, and
$f(0)^{m+1}w=0$, then every polynomial in the $f(k)$ applied to $w$ is a combination of box monomials applied to $w$.
\end{theorem}
\noindent File \leanpath{VOA/\allowbreak DeligneAudit/\allowbreak DeligneFusionSl2.lean}.
\begin{leancode}
def Zs : PowerSeries R := PowerSeries.mk fun i => z i

def xx (r s : ℕ) : R := PowerSeries.coeff s (Zs ^ r)

def J (m : ℕ) : Ideal R := Ideal.span {p | ∃ r s, 1 ≤ r ∧ m + 1 ≤ r + s ∧ p = xx r s}

def BoxP (m : ℕ) (f : ℕ →₀ ℕ) : Prop := ∀ i ∈ f.support, i + Finsupp.degree f ≤ m

def Bk (m k : ℕ) : Submodule ℂ R :=
  Submodule.span ℂ {p | ∃ f, BoxP m f ∧ f 0 ≤ k ∧ p = MvPolynomial.monomial f 1}

def T (m k : ℕ) : Submodule ℂ R := Bk m k ⊔ (J m).restrictScalars ℂ

theorem box_span : ∀ m (p : R), p ∈ T m m := by
\end{leancode}
\noindent File \leanpath{VOA/\allowbreak DeligneAudit/\allowbreak DeligneFusionSl2Module.lean}.
\begin{leancode}
theorem weyl_span_box {E1 : Module.End ℂ V} {H : ℕ → Module.End ℂ V}
    (hEF : ∀ k, E1 * F k - F k * E1 = H (k + 1))
    (hHF : ∀ m k, H (m + 1) * F k - F k * H (m + 1) = (-2 : ℂ) • F (m + 1 + k))
    {y : V} (hE : E1 y = 0) (hH : ∀ m, H (m + 1) y = 0) (m : ℕ) (hN : (F 0 ^ (m + 1)) y = 0)
    (p : R) :
    evAlg F hFF p y ∈
      Submodule.span ℂ {v | ∃ f, BoxP m f ∧ v = evAlg F hFF (MvPolynomial.monomial f 1) y} :=
\end{leancode}

\begin{dictionary}
\lean{R} is \lean{MvPolynomial ℕ ℂ} and \lean{z i} the variable $z_i$. A monomial is a finitely supported exponent
function \lean{f : ℕ →₀ ℕ}; \lean{Finsupp.degree f} is its total degree $\ell$, and \lean{BoxP m f} says every index $i$
that occurs satisfies $i+\ell\le m$. \lean{Bk m k} restricts the number of $z_0$'s to $\le k$ (used in the induction);
\lean{T m m} is box monomials plus $J_m$, so \lean{box_span} is exactly the first claim. In \lean{weyl_span_box} the
section variables are \lean{F : ℕ → Module.End ℂ V} and \lean{hFF} (the $F$'s commute); \lean{evAlg F hFF} is the
evaluation $z_i\mapsto F_i$ through the commutative subalgebra they generate. The hypotheses \lean{hEF}, \lean{hHF},
\lean{hE}, \lean{hH} are those of Garland (Theorem~\ref{l26:thm:garland}).
\end{dictionary}

\textbf{The cases $m\le2$ by hand.} Recall $x(1,s)=z_s$ and $x(2,s)=\sum_{i+j=s}z_iz_j$.
\begin{itemize}
\item $m=0$: $J_0$ contains every $x(1,s)=z_s$, so $R=\C\cdot1+J_0$; the only box monomial is $1$.
\item $m=1$: $J_1$ contains $z_s$ for $s\ge1$ and $x(2,0)=z_0^2$. Modulo $J_1$ every monomial is $0$ except $1$ and $z_0$; these are
  the box monomials ($\ell=1$ requires $i\le0$).
\item $m=2$: $J_2$ contains $z_s$ for $s\ge2$, $x(2,1)=2z_0z_1$, $x(2,2)=2z_0z_2+z_1^2$ and $x(3,0)=z_0^3$, but \emph{not}
  $x(2,0)=z_0^2$ ($r+s=2<3$). Modulo $J_2$: $z_0z_1\equiv0$, then $z_1^2\equiv-2z_0z_2\equiv0$, and $z_0^3\equiv0$. What survives is
  $1$; $z_0,z_1$; $z_0^2$: exactly the box monomials ($\ell=1$: $i\le1$; $\ell=2$: $i\le0$). So $\ltwentysixWloc(2)$ is spanned by
  four vectors, matching Example~\ref{l26:ex:sl2}.
\end{itemize}

\textbf{Proof on paper.} Induction on $m$, and for fixed $m$ on the number $k$ of factors $z_0$ in a monomial (the Lean
filtration \lean{Bk m k}). \emph{Monomials without $z_0$.} Such a monomial is $\sigma(p)$ for the shift
$\sigma:z_i\mapsto z_{i+1}$. Writing $Z(u)=z_0+u\,\sigma Z(u)$ gives $\sigma Z(u)^r=u^{-r}(Z(u)-z_0)^r$, hence
\begin{equation*}
\sigma x(r,s)=\sum_{j=1}^{r}\binom rj(-z_0)^{r-j}\,x(j,s+r)
\end{equation*}
(Chari--Pressley Lemma~6.5; the $j=0$ term vanishes because $s+r\ge1$). If $x(r,s)\in J_{m-1}$, i.e.\ $r+s\ge m$, every term has
$j+(s+r)\ge m+1$, so $\sigma(J_{m-1})\subseteq J_m$. By induction $p$ is a combination of box monomials for $m-1$ plus an element of
$J_{m-1}$; applying $\sigma$ raises every index by one, so box monomials for $m-1$ ($i+\ell\le m-1$) become box monomials for
$m$, and the error lands in $J_m$. \emph{Monomials with $z_0$.} The generator $x(r,s)$ with $s\ge1$ has the form
$r\,z_0^{r-1}z_s+(\text{terms with fewer }z_0)$. When a monomial $z_0^kq$ violates the box condition, a suitable $x(r,s)\in J_m$ with
this leading term lets one trade $z_0^{k}$ for terms with fewer factors $z_0$, which are box monomials modulo $J_m$ by the induction
on $k$ (Chari--Pressley Lemma~6.7). The resulting system of binomial identities is triangular; where Chari--Pressley invert a
binomial determinant (Lemma~6.6), the library instead shows that a polynomial with too many prescribed roots vanishes. The online
text of Chari--Pressley breaks off inside its inductive Claim, so the library's proof is a re-derivation. \emph{The module
statement.} Evaluating $z_i\mapsto f(i)$, the generators $x(r,s)$ of $J_m$ become Garland's elements $[u^s]F(u)^r$, which kill $w$
by Theorem~\ref{l26:thm:garland} (with $N=m+1$); so every polynomial in the $f(k)$ applied to $w$ is a combination of box monomials
applied to $w$. \qed

\begin{physics}
(Heuristic.) The $\mathfrak{sl}_2$ local Weyl module is the space of states of $m$ fused spin-$\frac12$ particles; the box rule says
the $\ell$-th excitation can carry at most $m-\ell$ units of degree, an exclusion rule reminiscent of fermionic quasi-particle bases.
\end{physics}

\begin{status}
Proved; standard axioms (\lean{box_span} re-run).
\end{status}

\subsection{Chari--Loktev: spanning for $\mathfrak{sl}_3$}

\begin{definition}[The Chari--Loktev family $\mathfrak B_N$]\label{l26:def:cl}
An index is a triple $(a,b,c)$ of weakly increasing lists of naturals; it is \emph{admissible} for $N$ if
$\max b\le N-|b|$, $\max c\le N-|c|$, $\max a\le N+|c|-|b|-|a|$ ($|\cdot|$ the length). Its word is
$f_1(a)\,f_\theta(b)\,f_2(c)$ ($f_2$ acting first), of degree $\sum a+\sum b+\sum c$. \textbf{CL spanning} at $N$: in every
module with the $M_{N,0}$ relations, $\ltwentysixcyc(w)_d$ is spanned by the words of admissible indices of degree $d$.
\textbf{CL count}: the number of admissible indices of degree $d$ equals $[q^d]\ltwentysixD_N$. \textbf{CL independence}:
the admissible words applied to $e^{N\theta}$ in the lattice vertex algebra $V_{A_2}$ are linearly independent.
\end{definition}
\noindent File \leanpath{VOA/\allowbreak DeligneAudit/\allowbreak DeligneFusionCL.lean}.
\begin{leancode}
def BoxT (B : ℤ) (s : List ℕ) : Prop := s.Pairwise (· ≤ ·) ∧ ∀ x ∈ s, (x : ℤ) ≤ B

def Adm (N : ℕ) (i : List ℕ × List ℕ × List ℕ) : Prop :=
  BoxT ((N : ℤ) - i.2.1.length) i.2.1 ∧ BoxT ((N : ℤ) - i.2.2.length) i.2.2 ∧
    BoxT ((N : ℤ) + i.2.2.length - i.2.1.length - i.1.length) i.1

def clWord (i : List ℕ × List ℕ × List ℕ) : List (Fin 8 × ℕ) :=
  i.1.map (fun s => ((3 : Fin 8), s)) ++ i.2.1.map (fun s => ((5 : Fin 8), s)) ++
    i.2.2.map (fun s => ((4 : Fin 8), s))

def tdeg (i : List ℕ × List ℕ × List ℕ) : ℕ := i.1.sum + i.2.1.sum + i.2.2.sum

def CLIdx (N d : ℕ) : Set (List ℕ × List ℕ × List ℕ) := {i | Adm N i ∧ tdeg i = d}

def CLCount (N : ℕ) : Prop := ∀ d, (CLIdx N d).ncard = A2.demazureHilb N d

def CLIndependence (N : ℕ) : Prop :=
  ∀ d, LinearIndependent ℂ (fun i : CLIdx N d => wordOp latRep (clWord i.1) (ch N N))

def CLSpanning (N : ℕ) : Prop :=
  ∀ (V : Type) [AddCommGroup V] [Module ℂ V] (ρ : CurrentModule (bracketTable) V) (w : V),
    MabRel ρ w N 0 → ∀ d, cycDeg ρ w d ≤
      Submodule.span ℂ (Set.range fun i : CLIdx N d => wordOp ρ (clWord i.1) w)
\end{leancode}

\begin{dictionary}
\lean{BoxT B s}: \lean{s} weakly increasing with entries $\le B$; $B\in\Z$, so a negative bound forces \lean{s} empty.
In \lean{Adm N (a, b, c)} the components are \lean{i.1}$=a$ ($f_1$-part), \lean{i.2.1}$=b$ ($f_\theta$-part),
\lean{i.2.2}$=c$ ($f_2$-part). \lean{clWord} maps index $3,5,4$ to $f_1,f_\theta,f_2$; \lean{wordOp} applies the last letter
first, so $f_2$ acts first. \lean{Set.ncard} is the cardinality of a finite set, and $0$ on an infinite one; \lean{CLCount} is meaningful
because \lean{CLIdx N d} is finite: \lean{BoxT} with a negative bound forces the list to be empty, so an admissible index has
$|b|,|c|\le N$ and $|a|\le N+|c|$, and every entry is bounded by $N+|c|\le2N$; finitely many weakly increasing lists remain.
\lean{latRep} and \lean{ch N N}
$=e^{N\theta}$ are from Theorem~\ref{l26:thm:latrep}. This is Chari--Loktev's set (2.6) at rank $r=2$,
$\lambda=N(\omega_1+\omega_2)$. Its total size is $\sum_{\ell,\ell'}\binom N\ell\binom N{\ell'}2^{N+\ell'-\ell}=9^N$.
\end{dictionary}

\begin{physics}
(Heuristic.) A Gelfand--Tsetlin ``quasi-particle'' basis: the $f_2$, $f_\theta$, $f_1$ excitations are added in that order, each with
an exclusion rule on its energy depending on how many of the others are present.
\end{physics}

\begin{status}
Definitions. \lean{CLCount} and \lean{CLSpanning} are proved for all $N$ (next theorem). \lean{CLIndependence} is a
named \lean{Prop} (Chari--Loktev's linear independence in the lattice model), checked numerically for $N\le4$ and not used:
U7 is proved without it.
\end{status}

\begin{theorem}[Chari--Loktev spanning and count for $\mathfrak{sl}_3$]\label{l26:thm:cl}
For every $N$, \lean{CLSpanning N} and \lean{CLCount N} hold. Consequently, in every module with the $M_{N,0}$ relations,
$\dim\ltwentysixcyc(w)_d\le[q^d]\ltwentysixD_N$ for every $d$.
\end{theorem}
\noindent File \leanpath{VOA/\allowbreak DeligneAudit/\allowbreak DeligneFusionCLAssembly.lean}.
\begin{leancode}
theorem clSpanning (N : ℕ) : CLSpanning N := by

theorem finrank_cycDeg_le {V : Type} [AddCommGroup V] [Module ℂ V] {ρ : CurrentModule (bracketTable) V}
    {w : V} {N : ℕ} (hR : MabRel ρ w N 0) (d : ℕ) :
    Module.finrank ℂ (cycDeg ρ w d) ≤ A2.demazureHilb N d := by
\end{leancode}
\noindent File \leanpath{VOA/\allowbreak DeligneAudit/\allowbreak DeligneFusionCLCount.lean}.
\begin{leancode}
theorem clCount (N : ℕ) : CLCount N := by
\end{leancode}

\begin{dictionary}
\lean{A2.demazureHilb N d} is $[q^d]$ of the Gelfand--Tsetlin polynomial \lean{gl3Hilb (2N) N 0}
(Chapter~\ref{ch:L25}); \lean{finrank_cycDeg_le} combines spanning with the count. This is an \emph{upper} bound on every
model, the universal one included: $M^{\mathrm{univ}}_{N,0}\le\ltwentysixD_N$.
\end{dictionary}

\textbf{The case $N=1$ by hand.} The admissible indices are those with $|b|,|c|\le1$ and entries bounded as in
Definition~\ref{l26:def:cl}. There are exactly nine:
\begin{equation*}
\begin{array}{lll}
w, & f_1(0)w, & f_2(0)w,\\
f_1(0)f_2(0)w, & f_1(0)^2f_2(0)w, & f_\theta(0)w,\\
f_\theta(0)f_2(0)w, & f_1(0)f_\theta(0)f_2(0)w, & f_1(1)f_2(0)w .
\end{array}
\end{equation*}
The first eight have degree $0$ and the last has degree $1$, matching $\ltwentysixD_1=8+q$ and the total $9^1$. In
$D(1,\theta)=\ltwentysixg\oplus\C$ the eight degree-$0$ words span the adjoint representation and $f_1(1)f_2(0)w$ spans the trivial
summand (compare Example~\ref{l26:ex:eval1}). For instance $f_1(1)$ is allowed after $f_2(0)$ because $f_2(0)$ raises the $h_1$-weight
from $1$ to $2$, and the $\alpha_1$-box then allows one letter of index up to $2-1=1$.

\textbf{Why the box conditions have this shape.} Write a word as $f_1(a)\,u$ with $u=f_\theta(b)f_2(c)$. The letters $f_\theta$ and
$f_2$ span an abelian ideal of the lowering subalgebra $\langle f_1,f_2,f_\theta\rangle$ (since $[f_1,f_2]=\pm f_\theta$ and $f_\theta$
commutes with both), so by PBW sorting every word in lowering letters is a combination of words of the form $f_1(\cdot)\cdots
f_\theta(\cdot)\cdots f_2(\cdot)\cdots w$; the $f_1$ part is analysed last, by restricting to the $\alpha_1$-copy of $\mathfrak{sl}_2[t]$.
This is the Gelfand--Tsetlin branching $\mathfrak{sl}_3\supset\mathfrak{sl}_2$. The three bounds:
\begin{itemize}
\item \emph{The $\alpha_2$-box}, $\max c\le N-|c|$: $w$ is a highest-weight vector of $h_2$-weight $N$ for the $\alpha_2$-copy of
  $\mathfrak{sl}_2[t]$, killed by $e_2(1)$ and the $h_2(r)$, $r\ge1$, with $f_2(0)^{N+1}w=0$; Theorem~\ref{l26:thm:box} with $m=N$ gives
  exactly this box for the $f_2$ letters.
\item \emph{The $\theta$-box}, $\max b\le N-|b|$: for the $\theta$-triple the weight of $w$ is $2N$, so the naive box would be $2N$.
  Chari--Loktev show (their Prop.~3.3.3) that in degree zero the vectors $f_1^{p-C}f_\theta^Cf_2^{q-C}w$ with $C>N$ are combinations of
  those with $C\le N$: excess $f_\theta$'s can be traded for pairs $f_1f_2$, because $f_\theta=\pm[f_1,f_2]$. With their Prop.~3.6.3
  (the $\theta$-box with bound $N$ on $f_2(c)w$ modulo a subspace stable under the relevant operators) this gives the bound $N$.
\item \emph{The $\alpha_1$-box}, $\max a\le(N+|c|-|b|)-|a|$: $f_2$ raises the $h_1$-weight by $1$ (as $\alpha_2(h_1)=-1$) and $f_\theta$
  lowers it by $1$ (as $\theta(h_1)=1$), so $u\,w$ has $h_1$-weight $N+|c|-|b|$. Modulo the span of words of higher index, $u\,w$ is
  a highest-weight vector for the $\alpha_1$-copy of $\mathfrak{sl}_2[t]$, and Theorem~\ref{l26:thm:box} with $m=N+|c|-|b|$ gives the box.
\end{itemize}

\textbf{Proof on paper.} One follows Chari--Loktev's proof for $\mathfrak{sl}_3$, restated on vectors of an arbitrary module
(there is no enveloping algebra in the library). Order the words $u=f_\theta(b)f_2(c)$ by an index (lengths, then sums). By
induction along this order, using the three bullets above (each applied modulo the span of the words of higher index), every word
applied to $w$ is a combination of admissible words of the same degree. Two technical points differ from the paper: Chari--Loktev's
auxiliary ideal $J'_N$ is replaced by a numerical statistic $\sigma(a,c)=|a|+\sum a+\sum c$ that does not increase under the
operators used, and error terms whose degree is not known in advance are controlled by lifting vectors to functions of degree and
weight. The count is a reindexing of the Gelfand--Tsetlin formula for $\ltwentysixD_N$, using that the Gaussian binomial
$\ltwentysixqb{n}{k}$ counts partitions in a $k\times(n-k)$ box. (In Lean: \lean{part_i}, \lean{alpha2_box}, \lean{prop333},
\lean{theta_box}, \lean{CLLift}, \lean{ncard_dec}.) \qed

\begin{physics}
(Heuristic.) The Demazure space is spanned by quasi-particle states obeying the exclusion rules; the count says these rules reproduce
the energy-graded partition function of the level-one boson's Demazure subspace.
\end{physics}

\begin{status}
Proved, unconditional; \texttt{\#print axioms} standard (\lean{clSpanning}, \lean{finrank_cycDeg_le} re-run). The
supporting files \lean{DeligneFusionECalc}, \lean{CLFilt}, \lean{CLGT}, \lean{CLNormal}, \lean{CLTheta}, \lean{CLLift} are
described in Section~\ref{l26:sec:certs}.
\end{status}

\subsection{Chari--Venkatesh: the $\mathfrak{sl}_2$ bound for partitions}

For $\mathfrak{sl}_2$ and a partition $\xi=(\xi_1\ge\dots\ge\xi_\ell>0)$ of $\lambda$, the \emph{Chari--Venkatesh module} $V(\xi)$ is
the $\mathfrak{sl}_2[t]$-module generated by a vector $y$ of weight $\lambda$, killed by all $e(r)$ and by $h(r)$, $r\ge1$, with
$f(0)^{\lambda+1}y=0$ and with the Garland coefficients $Q(r,s)y=[u^s]f(u)^ry$ set to zero in the region of
Theorem~\ref{l26:thm:cv}(1) below. Chari and Venkatesh prove $\dim V(\xi)=\prod_j(\xi_j+1)$ and identify $V(\xi)$ with the fusion
product $V(\xi_1)*\dots*V(\xi_\ell)$ of irreducible $\mathfrak{sl}_2$-modules. Two extreme cases: $\xi=(1,\dots,1)$ gives the
local Weyl module $\ltwentysixWloc(\lambda)$ of dimension $2^\lambda$, and $\xi=(\lambda)$ gives the evaluation module $V(\lambda)(0)$
of dimension $\lambda+1$ (the region then contains $Q(1,s)$, i.e.\ $f(s)y=0$, for all $s\ge1$). The library proves the upper bound.

\begin{theorem}[Chari--Venkatesh bound; truncated Weyl modules]\label{l26:thm:cv}
Let $E_s,H_p,F_k$ satisfy the $\mathfrak{sl}_2[t]$ relations, and $y$ be a vector with $E_sy=0$, $H_py=0$ ($p\ge1$),
$H_0y\in\C y$. (1) If $\xi$ is a partition (weakly decreasing, positive parts) and $y$ satisfies
$Q(r,s)y=0$ for all $r\ge1$, $s\ge0$ and $k\ge0$ with $1+rk+\sum_{j>k}\xi_j\le s+r$ (with $Q(r,s)=[u^s]F(u)^r$), then the $F$-span of $y$ is
finite-dimensional of dimension $\le\prod_j(\xi_j+1)$. (2) If $F_0^{\lambda+1}y=0$ and $F_jy=0$ for $j\ge r_0\ge1$, then
the same holds with $\xi=\mathrm{bal}(\lambda,r_0)$, the partition of $\lambda$ into $\min(\lambda,r_0)$ parts differing by at
most one.
\end{theorem}
\noindent File \leanpath{VOA/\allowbreak DeligneAudit/\allowbreak DeligneFusionCVBound.lean}.
\begin{leancode}
structure SL2 {V : Type u} [AddCommGroup V] [Module ℂ V] (E H F : ℕ → Module.End ℂ V) : Prop where
  hEF : ∀ s k, E s * F k - F k * E s = H (s + k)
  hHF : ∀ p k, H p * F k - F k * H p = (-2 : ℂ) • F (p + k)
  hHE : ∀ p q, H p * E q - E q * H p = (2 : ℂ) • E (p + q)
  hFF : ∀ i j, F i * F j = F j * F i
  hEE : ∀ i j, E i * E j = E j * E i

def CVRel {V : Type u} [AddCommGroup V] [Module ℂ V] (F : ℕ → Module.End ℂ V) (ξ : List ℕ) (y : V) :
    Prop :=
  ∀ r s k, 1 ≤ r → 1 + r * k + T k ξ ≤ s + r → Q F r s y = 0

def CVGoal (ξ : List ℕ) : Prop :=
  ∀ (V : Type u) [AddCommGroup V] [Module ℂ V] (E H F : ℕ → Module.End ℂ V), SL2 E H F →
    ∀ y : V, (∀ s, E s y = 0) → (∀ p, 1 ≤ p → H p y = 0) → (∃ c : ℂ, H 0 y = c • y) →
      CVRel F ξ y → FiniteDimensional ℂ (Fspan F y) ∧
        Module.finrank ℂ (Fspan F y) ≤ (ξ.map (· + 1)).prod

theorem cv_bound (ξ : List ℕ) (hs : ξ.Pairwise (· ≥ ·)) (hp : ∀ x ∈ ξ, 1 ≤ x) : CVGoal.{u} ξ :=
\end{leancode}
\noindent File \leanpath{VOA/\allowbreak DeligneAudit/\allowbreak DeligneFusionCVTrunc.lean}.
\begin{leancode}
def bal (lam r0 : ℕ) : List ℕ :=
  if lam < r0 then List.replicate lam 1
  else List.replicate (lam % r0) (lam / r0 + 1) ++ List.replicate (r0 - lam % r0) (lam / r0)

theorem tw_bound {V : Type u} [AddCommGroup V] [Module ℂ V] {E H F : ℕ → Module.End ℂ V}
    (hS : SL2 E H F) {y : V} (hE : ∀ s, E s y = 0) (hH : ∀ p, 1 ≤ p → H p y = 0)
    (hH0 : ∃ c : ℂ, H 0 y = c • y) {lam r0 : ℕ} (hr0 : 1 ≤ r0) (hF0 : (F 0 ^ (lam + 1)) y = 0)
    (htr : ∀ j, r0 ≤ j → F j y = 0) :
    FiniteDimensional ℂ (Fspan F y) ∧
      Module.finrank ℂ (Fspan F y) ≤ ((bal lam r0).map (· + 1)).prod :=
\end{leancode}

\begin{dictionary}
\lean{T k ξ} is $\sum_{j>k}\xi_j$ (\lean{(ξ.drop k).sum}, \leanpath{DeligneFusionCVPart.lean}); \lean{Q F r s} is
$[u^s]F(u)^r$ (\leanpath{DeligneFusionCVsl2.lean}); \lean{Fspan F y} is the span of all $F_{k_1}\cdots F_{k_n}y$. The
universe parameter \lean{.{u}} is bookkeeping. \lean{(ξ.map (· + 1)).prod} is $\prod(\xi_j+1)$. \lean{ξ.Pairwise (· ≥ ·)}
is ``weakly decreasing''. In \lean{CVRel} the quantifiers \lean{∀ r s k} range over all naturals, with \lean{1 ≤ r} the only
restriction; $k=0$ is included and gives the condition $1+|\xi|\le s+r$ with $|\xi|=\sum_j\xi_j$ (the parts are indexed from
$1$ in the text and from $0$ in \lean{List.drop}).
\end{dictionary}

\textbf{Proof on paper.} (1) Chari--Venkatesh \S6: induction on $(|\xi|,\ell(\xi),\xi_\ell)$. With
$y'=F_{\ell-1}^{\xi_\ell}y$ and $K$ the $F$-span of $y'$, $y$ modulo $K$ satisfies the relations of
$\xi^+=(\dots,\xi_{\ell-1}+1,\xi_\ell-1)$, and $y'$ satisfies those of $\xi^-=(\dots,\xi_{\ell-1}-\xi_\ell)$, while
$\prod(\xi_j+1)=\prod(\xi^+_j+1)+\prod(\xi^-_j+1)$. The second step uses CV Lemma~6.6, whose ``obvious induction'' is not
uniform in $j$; the library proves the precise shift (\lean{shift_aux}: the relation region shifts by $2$ for
$F_jv$ exactly when $j\ge k+1$, with exact shifts $\max(-2,k-1-j)$). (2) Every CV region of $\mathrm{bal}(\lambda,r_0)$ lies in
the Garland region $s+r\ge\lambda+1$ (Theorem~\ref{l26:thm:garland}) or in the pigeonhole region $s>(r_0-1)r$, where some
$F_j$ with $j\ge r_0$ occurs in every monomial (\lean{bal_region}). \qed

\begin{example}
$\lambda=3$, $r_0=2$: $\mathrm{bal}(3,2)=(2,1)$ and the bound is $3\cdot2=6=\dim V(2)*V(1)$. The naive ``box''
count with $F_ky=0$ for $k\ge2$ counts the seven monomials $1$; $z_0,z_1$; $z_0^2,z_0z_1,z_1^2$; $z_0^3$ (indices $i<2$,
$i+\ell\le3$) and gives $7$; see the dead ends of Section~\ref{l26:sec:open}.
\end{example}

\begin{physics}
(Heuristic.) Truncating the energy of each excitation of a fused spin chain distributes the $\lambda$ quanta as evenly as possible among
$r_0$ sites; the dimension is that of the fusion of the corresponding spins.
\end{physics}

\begin{status}
Proved; \texttt{\#print axioms} standard (\lean{cv_bound}, \lean{tw_bound} re-run). These results are \emph{not} used by
the headline chain. The relation \lean{CVRel} itself is reused: it is the building block of the three-root family through
which $\dim M_{0,m}\le8^m$ is proved (Section~\ref{l26:sec:cv}).
\end{status}

\subsection{U7 for $A_2$: the evaluation model}

\begin{definition}[Filtered relations; the evaluation model]\label{l26:def:evalmodel}
Let $T$ be a current module (no grading) and $w\in T$; let $\Phi_{<d}$ be the span of words of degree $<d$ applied to $w$
and $\Phi_{\le d}=\Phi_{<d+1}$. Say $w$ satisfies the \emph{filtered $M_{N,0}$ relations} if the raising and lowering
relations hold exactly, $h(r)w\in\C w$, $h(0)w=Nw$, and the truncations hold modulo lower filtration:
$f_i(N)w\in\Phi_{<N}$, $f_\theta(2N)w\in\Phi_{<2N}$. The \emph{evaluation model} is $T_N=\C^{\{1,2,3\}^{2N}}$, the tensor
product of $N$ copies of the standard representation $\C^3$ and $N$ copies of its dual, at the distinct points
$z=0,1,\dots,2N-1$, with $x(k)$ acting by $\sum_iz_i^k\,x^{(i)}$ ($x^{(i)}$ the action on factor $i$).
\end{definition}
\noindent File \leanpath{VOA/\allowbreak DeligneAudit/\allowbreak DeligneFusionGrFilt.lean}.
\begin{leancode}
def Fbelow (d : ℕ) : Submodule ℂ T :=
  Submodule.span ℂ {v | ∃ l : List (Fin 8 × ℕ), SectorGraded.wordDeg l < d ∧ v = wordOp ρ l w}

abbrev Fle (d : ℕ) : Submodule ℂ T := Fbelow ρ w (d + 1)

structure FiltRel (N : ℕ) : Prop where
  raising : ∀ (e : Fin 8), e.val < 3 → ∀ r, ρ.op e r w = 0
  cartan : ∀ (h : Fin 8), 6 ≤ h.val → ∀ r, ρ.op h r w ∈ ℂ ∙ w
  cartan_zero : ∀ (h : Fin 8), 6 ≤ h.val → ρ.op h 0 w = (N : ℂ) • w
  lowering_f1 : (ρ.op 3 0)^[N + 1] w = 0
  lowering_f2 : (ρ.op 4 0)^[N + 1] w = 0
  trunc_f1 : ρ.op 3 N w ∈ Fbelow ρ w N
  trunc_f2 : ρ.op 4 N w ∈ Fbelow ρ w N
  trunc_theta : ρ.op 5 (2 * N) w ∈ Fbelow ρ w (2 * N)
\end{leancode}
\noindent File \leanpath{VOA/\allowbreak DeligneAudit/\allowbreak DeligneFusionEvalModel.lean}.
\begin{leancode}
def matT : Fin 8 → Matrix (Fin 3) (Fin 3) ℂ :=
  ![Matrix.single 0 1 1, Matrix.single 1 2 1, Matrix.single 0 2 1, Matrix.single 1 0 1,
    Matrix.single 2 1 1, Matrix.single 2 0 1, Matrix.single 0 0 1 - Matrix.single 1 1 1,
    Matrix.single 1 1 1 - Matrix.single 2 2 1]

abbrev Site := Fin N ⊕ Fin N
abbrev Cfg := Site N → Fin 3
abbrev TT := Cfg N → ℂ

def R : Site N → Fin 8 → Matrix (Fin 3) (Fin 3) ℂ
  | .inl _, c => matT c
  | .inr _, c => -(Matrix.transpose (matT c))

def z : Site N → ℂ
  | .inl i => ((i : ℕ) : ℂ)
  | .inr i => ((N + i : ℕ) : ℂ)

def loc (i : Site N) (X : Matrix (Fin 3) (Fin 3) ℂ) : Module.End ℂ (TT N) where
  toFun f γ := ∑ k, X (γ i) k * f (Function.update γ i k)
  map_add' f g := by ext γ; simp [mul_add, Finset.sum_add_distrib]
  map_smul' c f := by ext γ; simp [Finset.mul_sum, mul_left_comm]

def op (c : Fin 8) (k : ℕ) : Module.End ℂ (TT N) := ∑ i, (z N i) ^ k • loc i (R N i c)
\end{leancode}

\begin{dictionary}
In \leanpath{DeligneFusionGrFilt.lean}, \lean{ρ} and \lean{w} are section variables, over a section variable
\lean{T : Type} with \lean{[AddCommGroup T] [Module ℂ T]}; no finiteness of $T$ is assumed in these definitions. \lean{ℂ ∙ w} is
the line through $w$. In \leanpath{DeligneFusionEvalModel.lean}, \lean{N} is a section variable: declared explicit
(\lean{variable (N : ℕ)}) before \lean{Site}, \lean{Cfg}, \lean{TT}, \lean{R}, \lean{z}, and made implicit (\lean{variable {N}})
for \lean{loc}, then explicit again for \lean{op}; this is why the printed text writes \lean{Site N} and \lean{z N i} but
\lean{loc i}. The \lean{op} printed here is the module-level \lean{EvalModel.op}, the operator family of the model, not the field
\lean{CurrentModule.op} (the two are linked by \lean{evalRep N}, whose \lean{op} field is this \lean{op N}).
\lean{Matrix.single i j 1} is the matrix unit $E_{ij}$ (so $e_1=E_{01}$, $f_\theta=E_{20}$, $h_1=E_{00}-E_{11}$);
\lean{mat_eq} and \lean{matT_bracket} check these against \lean{bracketTable}. \lean{Site N = Fin N ⊕ Fin N}: left sites
carry $\C^3$ (\lean{matT c}), right sites its dual (\lean{-(matT c)ᵀ}); a vector of $T_N$ is a function on
configurations \lean{Cfg N}. \lean{loc i X} applies the matrix $X$ at site $i$. The points are $z(\mathrm{inl}\ i)=i$,
$z(\mathrm{inr}\ i)=N+i$, all distinct. \lean{op c k} is $x_c(k)=\sum_iz_i^k\,x_c^{(i)}$, which assembles into the
current module \lean{evalRep N}; the generator \lean{w0 N} is the tensor product of highest-weight vectors.
\end{dictionary}

\begin{physics}
(Heuristic.) $2N$ quarks and antiquarks of $\mathfrak{su}(3)$ placed at distinct points of the line; the current modes act by
$\sum_iz_i^k\,(\cdot)^{(i)}$, the multipole moments of the charge distribution. Their fusion, as the points collide, is the
Demazure module.
\end{physics}

\begin{status}
Definitions.
\end{status}

\begin{theorem}[U7 for $A_2$]\label{l26:thm:u7}
(1) Let $T$ be a \emph{finite-dimensional} current module. If $w\in T$ satisfies the filtered $M_{N,0}$ relations and
$\dim\Phi_{\le D}\ge9^N$ for some $D$, then \lean{DemazureRealization N} holds. (2) The evaluation model satisfies the hypotheses of (1). Hence
\lean{DemazureRealization N} holds for every $N$.
\end{theorem}
\noindent File \leanpath{VOA/\allowbreak DeligneAudit/\allowbreak DeligneFusionGrFilt.lean}.
\begin{leancode}
theorem demazureRealization_of_filt {N : ℕ} (h : FiltRel ρ w N) (D : ℕ)
    (hD : 9 ^ N ≤ Module.finrank ℂ (Fle ρ w D)) : A2.DemazureRealization N := by
\end{leancode}
\noindent File \leanpath{VOA/\allowbreak DeligneAudit/\allowbreak DeligneFusionEvalModel.lean}.
\begin{leancode}
theorem demazureRealization_all (N : ℕ) : A2.DemazureRealization N := by
\end{leancode}

\begin{dictionary}
\textbf{A hidden instance hypothesis.} Besides the section variables \lean{T}, \lean{ρ}, \lean{w}, the theorem
\lean{demazureRealization_of_filt} carries the section variable \lean{[FiniteDimensional ℂ T]}, declared at line~182 of
\leanpath{DeligneFusionGrFilt.lean} (\lean{variable [FiniteDimensional ℂ T]}) and therefore not visible in the printed statement.
The full signature is
\lean{demazureRealization_of_filt ρ w [FiniteDimensional ℂ T] {N} (h : FiltRel ρ w N) (D) (hD : ...)}. It is used in
\lean{finrank_gr_ge} (the first theorem after that \lean{variable} line), which compares $\dim\Phi_{\le d}$ with
$\dim\Phi_{<d}+\dim G_d$ by rank--nullity and needs the pieces of $T$ to be finite-dimensional. The hypothesis is harmless in
the application: the evaluation model $T_N$ is a space of functions on the finite set \lean{Cfg N}, of dimension $9^N$, and
\lean{demazureRealization_all} supplies the instance. As stated in the source, (1) is therefore weaker than a version for
arbitrary $T$, which is why the theorem above says ``finite-dimensional''.

\lean{Fle ρ w D} is $\Phi_{\le D}$ and \lean{9 ^ N} its required total dimension. The conclusion of (1) is about a
\emph{different} module: the associated graded \lean{grRep ρ w} on $(\N\to_0T)/\mathrm{span}\{\text{single}(d,v):v\in\Phi_{<d}\}$,
with generator \lean{wg ρ w}. That module satisfies the exact relations \lean{MabRel N 0} (\lean{mabRel_gr}).
\end{dictionary}

\textbf{Proof on paper.} (1) The associated graded $G=\bigoplus_d\Phi_{\le d}/\Phi_{<d}$, with $x(k)$ raising degree by $k$, is a
graded module (as in Definition~\ref{l26:def:fusion}) whose generator $\bar w$, the class of $w$ in degree $0$, satisfies the exact
$M_{N,0}$ relations. The truncations become exact because of how $G$ is graded: $f_i(N)w$ is a word of degree $N$, so its class in
$G$ lies in $\Phi_{\le N}/\Phi_{<N}$; the hypothesis $f_i(N)w\in\Phi_{<N}$ says exactly that this class is $0$, i.e.\
$f_i(N)\bar w=0$ in $G$ (similarly for $f_\theta(2N)$). And $\sum_d\dim G_d=\dim\Phi_{\le D}\ge9^N$. By Theorem~\ref{l26:thm:cl}, $\dim G_d\le[q^d]\ltwentysixD_N$
for each $d$, and $\sum_d[q^d]\ltwentysixD_N=9^N$ (\lean{demazureHilb_sum}). A coefficientwise upper bound with equal totals
forces equality, so $\dim G_d=[q^d]\ltwentysixD_N$. (2) \emph{Relations.} Raising operators kill each highest-weight factor;
$h(r)w_0=(\sum_iz_i^r\lambda_i(h))w_0\in\C w_0$; $f_i(0)^{N+1}w_0=0$ by $\mathfrak{sl}_2$-theory on each factor and
nilpotency; the truncations: the polynomial $\prod_i(t-z_i)$ of degree $2N$ vanishes at every point, so
$f_\alpha(t^{2N})w_0$ is a combination of $f_\alpha(t^k)w_0$, $k<2N$ (Vandermonde, \lean{vand}); for $f_1$, only the $N$
sites where $f_1$ acts nontrivially on the highest weight vector matter, giving the truncation at $N$. \emph{Cyclicity.}
Lagrange interpolation gives $x^{(j)}=\sum_{k<2N}c_k\,x(k)$, with $c_k$ the coefficients of the polynomial that is $1$ at $z_j$ and $0$
at the other points; this isolates the action at site $j$ (\lean{loc_eq_sum}). Each factor is generated by its highest-weight vector
under $\ltwentysixg$, so every vector of $T_N$ is a combination of words in the site actions $x^{(j)}$ applied to $w_0$, and each
$x^{(j)}$ is a combination of $x(k)$ with $k<2N$. A word of length at most $L$ in the $x^{(j)}$ thus lies in $\Phi_{\le D}$ with
$D=(2N-1)L$, and since $T_N$ is finite-dimensional, finitely many words suffice: $\Phi_{\le D}=T_N$ for some $D$. The theorem needs
exactly this, $\dim\Phi_{\le D}\ge9^N$ for \emph{some} $D$, and $\dim T_N=3^{2N}=9^N$. \qed

\begin{physics}
(Heuristic.) The total number of states of $2N$ fundamental particles is $9^N$; fusing them cannot lose states, only redistribute them
in energy; and the Chari--Loktev exclusion rules say there is no room for more at any energy. Hence the energy
distribution is exactly the Demazure one.
\end{physics}

\begin{status}
Proved, unconditional; \texttt{\#print axioms} standard (\lean{demazureRealization_all},
\lean{demazureRealization_of_filt} re-run). The $A_2$ chain does not use \lean{CLIndependence}.
\end{status}

\subsection{The lattice model for $A_2$}

\begin{theorem}[Frenkel--Kac for $\mathfrak{sl}_3$; $e^{N\theta}$ satisfies the $M_{N,0}$ relations]\label{l26:thm:latrep}
Let $V_{A_2}$ be the lattice vertex algebra of the $A_2$ root lattice. The eight weight-one states
$e^{\pm\alpha}$ (rescaled by cocycle signs) and $\alpha_i(-1)\mathbf 1$ have $0$-products equal to \lean{bracketTable}; their
non-negative modes form a current module \lean{latRep} on $V_{A_2}$, and for every $N$ the vector $e^{N\theta}$ satisfies the
relations of $M_{N,0}$.
\end{theorem}
\noindent File \leanpath{VOA/\allowbreak DeligneAudit/\allowbreak DeligneFusionLattice.lean}.
\begin{leancode}
abbrev WA : VertexOperatorAlgebra (LatModuleGen a2) := latticeVOAGen a2 a2_even a2_posDef

abbrev ch (a b : ℤ) : LatModuleGen a2 := latChargeGen a2 (lv a b)

def Xs : Fin 8 → LatModuleGen a2 :=
  ![ch 1 0, ch 0 1, s12 • ch 1 1, s11 • ch (-1) 0, s22 • ch 0 (-1),
    (s11 * s21 * s22) • ch (-1) (-1), h1, h2]

theorem Xs_bracket (c d : Fin 8) :
    WA.smul (Xs c) 0 (Xs d) = ∑ e, AffineSl3.bracketTable c d e • Xs e := by

def latRep : SectorWeyl.CurrentModule (AffineSl3.bracketTable) (LatModuleGen a2) where
  op c n := mode (WA.Y (Xs c)) (n : ℤ)
  comm c d r s v := by
\end{leancode}
\begin{leancode}
theorem latRep_mabRel (N : ℕ) : MabRel latRep (ch N N) N 0 where
\end{leancode}

\begin{dictionary}
\lean{latChargeGen a2 β} is $e^\beta$; \lean{ch a b}$=e^{a\alpha_1+b\alpha_2}$, so \lean{ch N N}$=e^{N\theta}$. \lean{s11}, \lean{s12},
\lean{s21}, \lean{s22} are values of the lattice cocycle $\varepsilon(\alpha_i,\alpha_j)\in\{\pm1\}$. \lean{WA.smul x n y} is the
$n$-th product $x_{(n)}y$; \lean{mode (WA.Y x) n} is the mode $x_{(n)}$. The current-module property uses the commutator
formula for weight-one fields, whose central term vanishes because $\langle\cdot\rangle$-products of weight-one states land in
$V_0=\C\mathbf 1$ and $r+s>0$ (\lean{Deform.current_mode_commutator}).
\end{dictionary}

\textbf{Proof on paper.} The $0$-products are the Frenkel--Kac formulas: $e^\alpha_{(0)}e^\beta=\varepsilon(\alpha,\beta)e^{\alpha+\beta}$ if
$\langle\alpha,\beta\rangle=-1$, $=\varepsilon\,\alpha(-1)\mathbf 1$ if $\beta=-\alpha$, and $0$ otherwise; checked in eight row lemmas.
For $e^{N\theta}$: $e^\alpha_{(r)}e^{N\theta}=0$ when $r\ge-\langle\alpha,N\theta\rangle$, which gives the raising relations
($\langle\alpha,\theta\rangle\ge0$), the truncations $f_i(N)$ ($\langle-\alpha_i,\theta\rangle=-1$) and $f_\theta(2N)$; Cartan modes
act by $\langle\alpha_i,N\theta\rangle=N$ at $r=0$ and by $0$ for $r\ge1$; and $f_i(0)^{N+1}e^{N\theta}$ lies in charge
$N\theta-(N+1)\alpha_i$ of half-norm $N^2+N+1>N^2$ at conformal weight $N^2$, where the space is zero (the weight shell). \qed

\begin{physics}
The level-one $\widehat{\mathfrak{sl}}_3$ currents are realized by two free bosons through vertex operators (Frenkel--Kac).
The state $e^{N\theta}$, of conformal weight $N^2$, is an extremal state; the currents generate from it the Demazure subspace.
\end{physics}

\begin{status}
Proved, unconditional; \texttt{\#print axioms} standard (\lean{latRep_mabRel} re-run). The \emph{dimension} half
(the cyclic module of $e^{N\theta}$ has graded dimension $\ltwentysixD_N$) is not proved in Lean; it is not needed for $A_2$
(Theorem~\ref{l26:thm:u7}). Numerically it holds for $N\le3$.
\end{status}

%======================================================================
\section{Reductions of the $A_2$ total bound}\label{l26:sec:a2red}
%======================================================================

After Theorems~\ref{l26:thm:a2assembly} and \ref{l26:thm:u7}, the $A_2$ case rests on \lean{FusionTotalDim} alone. That
statement is proved in Section~\ref{l26:sec:cv}. This section records two sharper statements that also imply it. They are not
proved, and since \lean{FusionTotalDim} is now a theorem they are not needed; they are alternative routes, kept in the library.

\begin{definition}[Kernel lower bound; ideal lower bound]\label{l26:def:s2}
\leavevmode
\begin{enumerate}
\item[(a)] \textbf{\lean{RavinderKernelLower}} (the \emph{kernel lower bound}): for all $a,b$ there is a module with the
relations of $M_{a+1,b}$ in which the cyclic module of the kernel vector $f_\theta(2a+b+1)w$ has total dimension at least
$9^a8^b$.
\item[(b)] \textbf{\lean{IdealLower}} (the \emph{ideal lower bound}): for every $N$ there is a module with the relations of
$M_{N,0}$ in which the subspace $I\cdot\langle w\rangle$, spanned by $x(k)\cdot(\text{word})\cdot w$ with $k\ge N$, has total
dimension at least $9^N-8^N$.
\end{enumerate}
\end{definition}
\noindent File \leanpath{VOA/\allowbreak DeligneAudit/\allowbreak DeligneFusionKernelLower.lean}.
\begin{leancode}
def RavinderKernelLower : Prop :=
  ∀ a b : ℕ, ∃ (V : Type) (_ : AddCommGroup V) (_ : Module ℂ V) (ρ : CurrentModule (bracketTable) V)
    (w : V), MabRel ρ w (a + 1) b ∧
    ∃ B, 9 ^ a * 8 ^ b ≤ ∑ n ∈ Finset.range B,
      Module.finrank ℂ (cycDeg ρ (ρ.op 5 (2 * a + b + 1) w) n)

abbrev uTot (a b B : ℕ) : ℕ :=
  ∑ M ∈ Finset.range B, Module.finrank ℂ (cycDeg (univRep a b) (SectorWeyl.wU a b) M)
\end{leancode}
\noindent File \leanpath{VOA/\allowbreak DeligneAudit/\allowbreak DeligneFusionIdeal.lean}.
\begin{leancode}
def IVd {V : Type} [AddCommGroup V] [Module ℂ V] (ρ : CurrentModule (bracketTable) V) (w : V) (N d : ℕ) :
    Submodule ℂ V :=
  Submodule.span ℂ {v | ∃ (c : Fin 8) (k : ℕ) (l : List (Fin 8 × ℕ)), N ≤ k ∧
    k + SectorGraded.wordDeg l = d ∧ v = ρ.op c k (wordOp ρ l w)}

def IdealLower : Prop :=
  ∀ N : ℕ, ∃ (V : Type) (_ : AddCommGroup V) (_ : Module ℂ V) (ρ : CurrentModule (bracketTable) V)
    (w : V), MabRel ρ w N 0 ∧ ∃ B, 9 ^ N ≤ 8 ^ N + ∑ d ∈ Finset.range B, Module.finrank ℂ (IVd ρ w N d)
\end{leancode}

\begin{dictionary}
Both are existence statements of \emph{lower} bounds, stated as \lean{Prop}s. \lean{IVd ρ w N d} is the degree-$d$ part
of $(\ltwentysixg\otimes t^N\C[t])\cdot\langle w\rangle$ in the sense of words: letter $(c,k)$ with $k\ge N$ applied to a word of
degree $d-k$. The inequality \lean{9 ^ N ≤ 8 ^ N + ∑ ...} is ``$\sum_d\dim\ge9^N-8^N$'' without natural subtraction.

\textbf{What is and is not established about them.} Neither is trivially true: in any model of $M_{a+1,b}$ the cyclic
module of the kernel vector is a quotient of $U_{a,b}$ (kernel lemma), so \lean{RavinderKernelLower} forces
$\dim U_{a,b}\ge9^a8^b$, a genuine statement about the universal modules; and \lean{IdealLower} asks for a large ideal image
in a model of $M_{N,0}$. The \emph{only} relations to the total bound that the library proves are the two implications
\lean{RavinderKernelLower} $\Rightarrow$ \lean{FusionTotalDim} and \lean{IdealLower} $\Rightarrow$ \lean{FusionTotalDim}
(\lean{fusionTotalDim_of_kernelLower}, \lean{fusionTotalDim_of_idealLower}, Theorem~\ref{l26:thm:a2red}). The converses are
\emph{not} formalized. They are not consequences of what is proved here either: the spanning inequality
$\mathrm{tot}(a+1,b)\le\mathrm{tot}(a,b+1)+\mathrm{tot}(a,b)$ of Theorem~\ref{l26:thm:step} (with $\mathrm{tot}(a,b)$ the total dimension
of $U_{a,b}$), with $\mathrm{tot}(a,0)\le9^a$, does not bound $\mathrm{tot}(a,b)$ by $9^a8^b$. Mathematically
the converses would follow, conditionally on \lean{FusionTotalDim}, from Ravinder's exact sequence, which itself rests on the
literature input whose gap is discussed in Section~\ref{l26:sec:open}; they are literature-level claims, not
results of this chapter. The two statements were introduced because they are \emph{lower} bounds, which can in principle be
witnessed in a single explicit model, whereas \lean{FusionTotalDim} is an upper bound over all modules. With
\lean{FusionTotalDim} proved (Theorem~\ref{l26:thm:cvtotal}), whether they hold is no longer needed for any headline result; it is
not proved.
\end{dictionary}

\begin{physics}
(Heuristic.) \lean{RavinderKernelLower}: in the system of $a+1$ Demazure and $b$ adjoint particles, the states reachable after one particle
emits its top excitation are as many as the states of the smaller system. \lean{IdealLower}: the states killed when every
particle is forbidden to carry more than $N-1$ units of energy are at least $9^N-8^N$.
\end{physics}

\begin{status}
Named \lean{Prop}s, never axioms; neither is proved. Numerically, in $V_{A_2}$, \lean{IdealLower} holds with equality for
$N\le3$: the ideal image $J_N=(\ltwentysixg\otimes t^N\C[t])\cdot D(N)$ in the cyclic module $D(N)$ of $e^{N\theta}$ has
$\dim J_N=1,17,217=9^N-8^N$ and the kernel pieces of \lean{RavinderKernelLower} are exactly $9^{N-1-i}8^i$ for $N\le4$
(Section~\ref{l26:sec:open}).
\end{status}

\begin{theorem}[Each refinement implies the total bound]\label{l26:thm:a2red}
(1) \lean{RavinderKernelLower} implies $\sum_{N<B}\dim\ltwentysixcyc_{U_{a,b}}(w_U)_N\le9^a8^b$ for all $a,b,B$, hence
\lean{FusionTotalDim}. (2) In a module with the relations of $M_{0,N}$, every mode $x(k)$ with $k\ge N$ kills the whole
cyclic module. (3) \lean{IdealLower} implies \lean{FusionTotalDim}.
\end{theorem}
\noindent File \leanpath{VOA/\allowbreak DeligneAudit/\allowbreak DeligneFusionKernelLower.lean}.
\begin{leancode}
theorem uTot_le (hK : RavinderKernelLower) : ∀ b a B, uTot a b B ≤ 9 ^ a * 8 ^ b

theorem fusionTotalDim_of_kernelLower (hK : RavinderKernelLower) : FusionTotalDim := by
\end{leancode}
\noindent File \leanpath{VOA/\allowbreak DeligneAudit/\allowbreak DeligneFusionIdeal.lean}.
\begin{leancode}
theorem trunc_all {V : Type} [AddCommGroup V] [Module ℂ V] {ρ : CurrentModule (bracketTable) V}
    {w : V} {N : ℕ} (h : MabRel ρ w 0 N) :
    ∀ (l : List (Fin 8 × ℕ)) (c : Fin 8) (k : ℕ), N ≤ k → ρ.op c k (wordOp ρ l w) = 0

theorem fusionTotalDim_of_idealLower (hI : IdealLower) : SectorWeyl.FusionTotalDim := by
\end{leancode}

\begin{dictionary}
\lean{uTot a b B} is $\sum_{N<B}\dim\ltwentysixcyc_{U_{a,b}}(w_U)_N$. The induction in \lean{uTot_le} is on $b$ (first
argument), starting from the Chari--Loktev bound at $b=0$ (where $\sum_N[q^N]\ltwentysixD_a=9^a$). \lean{SectorWeyl.FusionTotalDim}
is Theorem~\ref{l26:thm:finite}'s statement.
\end{dictionary}

\textbf{Proof on paper.} (1) The relations of $M_{a,b+1}$ imply those of $M_{a+1,b}$ (\lean{mabRel_succ}: the only
difference is $f_\theta(2a+b+2)w=0$ versus $f_\theta(2a+b+1)w=0$, and the former follows from the latter via
$[h_\theta(1),f_\theta(r)]=-2f_\theta(r+1)$). So there is a surjection $\pi:U_{a+1,b}\to U_{a,b+1}$ killing the cyclic module of
$v_K$. Rank--nullity in each degree gives $u_{a,b+1}(n)+\dim(U\cdot v_K)_n\le u_{a+1,b}(n)$ (\lean{step_kernel}), and in any
model the cyclic module of $v_K$ is a quotient of the universal one, so the model's count is a lower bound
(\lean{kU_ge_model}). Summing and using \lean{RavinderKernelLower}: $\mathrm{uTot}(a,b+1)\le\mathrm{uTot}(a+1,b)-9^a8^b$. Induction on $b$ from
$\mathrm{uTot}(a,0)\le9^a$ gives $\mathrm{uTot}(a,b)\le9^{a+b}-\sum(\dots)=9^a8^b$, because
$\sum_{i<m}9^{m-1-i}8^i=9^m-8^m$. At $a=0$ this is $\le8^m$ for the universal module, which bounds every module.
(2) The subspace $\ltwentysixg\otimes t^N\C[t]$ is a Lie ideal of $\ltwentysixgt$ that kills $w$ in $M_{0,N}$ (raising and
Cartan relations, and the truncations propagated upward); an ideal that kills the generator kills the cyclic module.
(3) By (2), $U_{N,0}\to U_{0,N}$ kills $I\cdot U_{N,0}$, so $u_{0,N}(d)+\dim(I\cdot U_{N,0})_d\le u_{N,0}(d)$
(\lean{step_ideal}); the model's \lean{IVd} is a lower bound for the universal one (\lean{IVd_ge_model}); and $\sum_du_{N,0}(d)\le9^N$
by Chari--Loktev. \qed

\begin{physics}
(Heuristic.) Counting by subtraction: the $m$-adjoint system is what remains of the $m$-Demazure system after removing all states with a
particle in a bound configuration; a lower bound on what is removed is an upper bound on what remains.
\end{physics}

\begin{status}
Proved, \texttt{\#print axioms} standard (\lean{fusionTotalDim_of_kernelLower}, \lean{fusionTotalDim_of_idealLower},
\lean{trunc_all} re-run). Conditional on the respective named \lean{Prop}.
\end{status}

\begin{theorem}[The $A_2$ case: final forms in this chapter]\label{l26:thm:u7final}
$L_{-3/2}(\mathfrak{sl}_3)$ is classically free for its explicit presentation, with kernel $(I_2^{\mathrm{ns}})$, assuming
either \lean{RavinderKernelLower} or \lean{FusionTotalDim}.
\end{theorem}
\noindent File \leanpath{VOA/\allowbreak DeligneAudit/\allowbreak DeligneFusionU7.lean}.
\begin{leancode}
theorem classicallyFree_of_kernelLower' (hK : RavinderKernelLower) :
    CaseClassicallyFree .A2 W explicitPresentation ∧
      RingHom.ker (MvPolynomial.aeval fun i =>
        liDegreeZeroClass W 1 ⟨explicitPresentation.gen i,
          explicitPresentation.gen_weight_one i⟩).toRingHom = I2ns :=

theorem classicallyFree_of_totalDim' (hF : SectorWeyl.FusionTotalDim) :
    CaseClassicallyFree .A2 W explicitPresentation ∧
      RingHom.ker (MvPolynomial.aeval fun i =>
        liDegreeZeroClass W 1 ⟨explicitPresentation.gen i,
          explicitPresentation.gen_weight_one i⟩).toRingHom = I2ns :=
\end{leancode}

\begin{dictionary}
The second, applied to \lean{CVUniv.fusionTotalDim} (Theorem~\ref{l26:thm:cvtotal}), is the unconditional
\lean{A2.classicallyFree} of Chapter~\ref{ch:L27}, which \lean{DeligneSeries.A2_classicallyFree_final} re-exports; the first is
re-exported as \lean{A2_classicallyFree_final_S2}. U7 is supplied by \lean{EvalModel.demazureRealization_all}; the kernel lemma by \lean{ravinderKernelTrunc}; the
finiteness by Theorem~\ref{l26:thm:finite}. The only hypothesis left is the printed one. The conclusion is the full
target of Chapter~\ref{ch:L24}: classical freeness and exactness of the presentation.
\end{dictionary}

\textbf{Proof on paper.} Theorem~\ref{l26:thm:a2assembly} with Theorems~\ref{l26:thm:u7} and \ref{l26:thm:a2red}. \qed

\begin{physics}
Classical freeness (Chapter~\ref{ch:L24}) identifies the vacuum character of $L_{-3/2}(\mathfrak{sl}_3)$ with the Hilbert series
of the arc space of its associated scheme $\operatorname{Spec}\operatorname{Sym}\mathfrak{sl}_3/(I_2^{\mathrm{ns}})$, which is
non-reduced (its Casimir is nilpotent but not zero); that much is mathematics. The further identification of this vacuum character with the Schur index of the Argyres--Douglas theory $H_2$ is a
\emph{physics input}: the four-dimensional/two-dimensional correspondence of Beem et al.\ (arXiv:1312.5344; for the comparison with
arc spaces see Beem--Rastelli, arXiv:1707.07679). Granting it, the physical statement ``Schur index $=$ arc-space Hilbert series''
for $H_2$ follows from the unconditional Lean theorem of Chapter~\ref{ch:L27}. The Lean theorem says nothing about four
dimensions.
\end{physics}

\begin{status}
Standard axioms (re-run at an earlier commit, not re-run for this edition), \textbf{conditional} on one named \lean{Prop} as stated. The second form
becomes unconditional when \lean{CVUniv.fusionTotalDim} is supplied (Chapter~\ref{ch:L27}); the first stays conditional on the
unproved \lean{RavinderKernelLower}.
\end{status}

%======================================================================
\section[The $A_2$ total bound in Lean]{The $A_2$ total bound in Lean: the Chari--Venkatesh family}\label{l26:sec:cv}
%======================================================================

The total bound \lean{FusionTotalDim} ($\dim F_m\le8^m$ for every module with the relations of $M_{0,m}$) is a theorem of the
library, \lean{CVUniv.fusionTotalDim}. Its proof does not go through Feigin's Prop.~1, through \lean{RavinderKernelLower} or
\lean{IdealLower}, or through the slice lemma of Section~\ref{l26:sec:open}. It follows the enlarged-family argument of
Setia, Rani and Khandai (arXiv:2405.11944, \S8): one embeds $M_{0,m}$ into a larger family of $\mathfrak{sl}_3[t]$-modules defined by
Chari--Venkatesh relations along the three positive roots, proves a dimension bound for the whole family by induction, and reads
off the case of $m$ adjoint factors. We present the family, the two base statements, and the resulting theorem. The internal
construction (quotient maps between universal modules, the associated graded of the PBW filtration, and the finite-site
counts) is spread over the 139 files whose names begin \texttt{DeligneFusionCV}; we describe it only on paper.

\begin{definition}[The three-root Chari--Venkatesh family]\label{l26:def:cvfamily}
Write a dominant $\mathfrak{sl}_3$-weight as a pair $(\lambda_1,\lambda_2)$ of coordinates in the fundamental weights
$\omega_1,\omega_2$, so $\theta=(1,1)$. For pairs $\varepsilon,\nu,\rho\in\N^2$ and $j\in\N$ put
$\lambda=\varepsilon+2\nu+3\rho+j\theta$, and for each positive root a partition with parts $3,2,1$:
\begin{equation*}
\xi^{(1)}=(3^{\rho_1},2^{\nu_1},1^{\varepsilon_1+j}),\qquad \xi^{(2)}=(3^{\rho_2},2^{\nu_2},1^{\varepsilon_2+j}),\qquad
\xi^{(\theta)}=(3^{\rho_1+\rho_2},2^{\nu_1+\nu_2+j},1^{\varepsilon_1+\varepsilon_2}) .
\end{equation*}
A vector $w$ of a current module satisfies the \emph{CV relations} $\mathrm{CV}(\varepsilon,\nu,\rho;j)$ if $e_\alpha(k)w=0$ for all
$k$ and $\alpha>0$, $h_i(k)w=0$ for $k\ge1$, $h_i(0)w=\lambda_iw$, $f_i(0)^{\lambda_i+1}w=0$, and, along each positive root
$\alpha\in\{\alpha_1,\alpha_2,\theta\}$, the $\mathfrak{sl}_2$ Chari--Venkatesh relations of Theorem~\ref{l26:thm:cv} for the
partition $\xi^{(\alpha)}$ (\lean{CVRel}). $C_j(\varepsilon,\nu,\rho)$ is the universal cyclic module with these relations
(Definition~\ref{l26:def:relspec}), and
\begin{equation*}
\mathrm{cvCount}(\varepsilon,\nu,\rho;j)=3^{|\varepsilon|}\,6^{|\nu|}\,10^{|\rho|}\,8^{j},\qquad |\mu|=\mu_1+\mu_2 .
\end{equation*}
The \emph{rectangular base family} is the statement: for $\varepsilon\in\{(0,0),(1,0),(0,1)\}$ and all $\nu,\rho$,
$C_0(\varepsilon,\nu,\rho)$ is finite-dimensional of dimension $\le\mathrm{cvCount}(\varepsilon,\nu,\rho;0)$. The \emph{residual
base family} is: for all $\rho$, $C_1(0,0,\rho)$ is finite-dimensional of dimension $\le8\cdot10^{|\rho|}$.
\end{definition}
\noindent Files \leanpath{VOA/\allowbreak DeligneAudit/\allowbreak DeligneFusionCVFamily.lean} (\lean{parts}, \lean{CV3Rel}),
\leanpath{VOA/\allowbreak DeligneAudit/\allowbreak DeligneFusionCVUniv.lean} (\lean{CVMod}, \lean{sat_cvSpec_iff}) and
\leanpath{VOA/\allowbreak DeligneAudit/\allowbreak DeligneFusionCVTotal.lean} (the counts and the two base statements).
\begin{leancode}
abbrev Weight := ℕ × ℕ

def parts (r n e : ℕ) : List ℕ :=
  List.replicate r 3 ++ List.replicate n 2 ++ List.replicate e 1

structure CV3Rel {V : Type} [AddCommGroup V] [Module ℂ V]
    (R : CurrentModule bracketTable V) (w : V) (ε ν ϱ : Weight) (j : ℕ) : Prop where
  raising : ∀ (e : Fin 8), e.val < 3 → ∀ k, R.op e k w = 0
  cartan_pos : ∀ (h : Fin 8), 6 ≤ h.val → ∀ k, 1 ≤ k → R.op h k w = 0
  cartan_first : R.op 6 0 w = ((ε.1 + 2 * ν.1 + 3 * ϱ.1 + j : ℕ) : ℂ) • w
  cartan_second : R.op 7 0 w = ((ε.2 + 2 * ν.2 + 3 * ϱ.2 + j : ℕ) : ℂ) • w
  lowering_first : (R.op 3 0)^[ε.1 + 2 * ν.1 + 3 * ϱ.1 + j + 1] w = 0
  lowering_second : (R.op 4 0)^[ε.2 + 2 * ν.2 + 3 * ϱ.2 + j + 1] w = 0
  first : CVRel (R.op 3) (parts ϱ.1 ν.1 (ε.1 + j)) w
  second : CVRel (R.op 4) (parts ϱ.2 ν.2 (ε.2 + j)) w
  highest : CVRel (R.op 5) (parts (ϱ.1 + ϱ.2) (ν.1 + ν.2 + j) (ε.1 + ε.2)) w

theorem mabRel_zero_cv3 {N : ℕ} (h : MabRel R w 0 N) (hN : 0 < N) :
    CV3Rel R w (0, 0) (0, 0) (0, 0) N := by
\end{leancode}
\begin{leancode}
theorem sat_cvSpec_iff :
    UnivG.Sat bracketTable R w (cvSpec ε ν ϱ j) ↔ CV3Rel R w ε ν ϱ j := by

abbrev CVMod (ε ν ϱ : Weight) (j : ℕ) := UnivG.UModG bracketTable (cvSpec ε ν ϱ j)
\end{leancode}
\begin{leancode}
def SmallEpsilon (ε : Weight) : Prop := ε = (0, 0) ∨ ε = (1, 0) ∨ ε = (0, 1)

def cvCount (ε ν ϱ : Weight) (j : ℕ) : ℕ :=
  3 ^ (ε.1 + ε.2) * 6 ^ (ν.1 + ν.2) * 10 ^ (ϱ.1 + ϱ.2) * 8 ^ j

def CVBound (ε ν ϱ : Weight) (j : ℕ) : Prop :=
  FiniteDimensional ℂ (CVMod ε ν ϱ j) ∧
    Module.finrank ℂ (CVMod ε ν ϱ j) ≤ cvCount ε ν ϱ j

def SharpRectangularBases : Prop :=
  ∀ (ε ν ϱ : Weight), SmallEpsilon ε → CVBound ε ν ϱ 0

def ResidualTheta3Bound : Prop :=
  ∀ ϱ : Weight, FiniteDimensional ℂ (CVMod (0, 0) (0, 0) ϱ 1) ∧
    Module.finrank ℂ (CVMod (0, 0) (0, 0) ϱ 1) ≤ 8 * 10 ^ (ϱ.1 + ϱ.2)
\end{leancode}

\begin{dictionary}
\textbf{Indices.} The \lean{bracketTable} basis is that of Section~\ref{l26:sec:a2rel}: \lean{0,1,2} $=e_1,e_2,e_\theta$,
\lean{3,4,5} $=f_1,f_2,f_\theta$, \lean{6,7} $=h_1,h_2$. So \lean{raising} says $e_\alpha(k)w=0$ for all $k$;
\lean{cartan_pos}, $h_i(k)w=0$ for $k\ge1$; \lean{cartan_first}/\lean{second}, $h_i(0)w=\lambda_iw$; \lean{lowering_*},
$f_i(0)^{\lambda_i+1}w=0$ (\lean{f^[n]} is the $n$-fold iterate). \lean{parts r n e} is the partition $(3^r,2^n,1^e)$ as a
list, and \lean{CVRel F ξ w} is the Chari--Venkatesh condition of Theorem~\ref{l26:thm:cv}: all Garland coefficients
$Q(r,s)w=[u^s]F(u)^rw$ vanish in the region $1+rk+T_k(\xi)\le s+r$, $r\ge1$, where $T_k(\xi)$ is the sum of the parts after the
first $k$. The three \lean{CVRel} fields use the partitions $\xi^{(1)},\xi^{(2)},\xi^{(\theta)}$ of the definition, for the
$\mathfrak{sl}_2$ triples of $\alpha_1$, $\alpha_2$, $\theta$ (the letters \lean{3}, \lean{4}, \lean{5}).

\textbf{The universal module.} \lean{cvSpec ε ν ϱ j} is a relation specification (Definition~\ref{l26:def:relspec}) whose tails
are the single words $e_\alpha(1)^sf_\alpha(0)^{r+s}$ that Garland's identity makes equivalent to the coefficient relations, and
\lean{sat_cvSpec_iff} proves that satisfying it is \emph{exactly} \lean{CV3Rel}. \lean{CVMod} is the universal module
\lean{UModG} of that specification, so every module with a vector satisfying \lean{CV3Rel} receives the cyclic span of that vector
as a quotient of \lean{CVMod} (Definition~\ref{l26:def:relspec}). The variables \lean{R}, \lean{w}, \lean{ε}, \lean{ν}, \lean{ϱ},
\lean{j} of \lean{sat_cvSpec_iff} are implicit section variables. Lean's \lean{ϱ} is the $\rho$ of the text.

\textbf{The link with $M_{0,N}$.} \lean{mabRel_zero_cv3}: if $w$ satisfies the relations of $M_{0,N}$ (\lean{MabRel R w 0 N},
Definition~\ref{l26:def:mab}) with $N\ge1$, then it satisfies $\mathrm{CV}(0,0,0;N)$. The highest weight is $N\theta$, and the
partitions are $(1^N)$, $(1^N)$, $(2^N)$. So the family contains the pure-adjoint case at $\varepsilon=\nu=\rho=0$, $j=N$, with
$\mathrm{cvCount}=8^N$.

\textbf{The counts.} $3,6,10,8$ are $\dim V(\omega_i)$, $\dim V(2\omega_i)$, $\dim V(3\omega_i)$, $\dim V(\theta)$. A bound
$\dim C_j(\varepsilon,\nu,\rho)\le\mathrm{cvCount}$ is what one expects if $C_j(\varepsilon,\nu,\rho)$ is the fusion product of
$|\varepsilon|$ fundamental, $|\nu|$ ``two-box'' and $|\rho|$ ``three-box'' Kirillov--Reshetikhin modules and $j$ adjoints.
That identification is \emph{not} stated in Lean; only the upper bound is. \lean{CVBound} includes finite-dimensionality, so
the \lean{finrank} in it is not the value $0$ that Lean assigns to infinite-dimensional spaces.

\textbf{Could it be vacuous?} No. \lean{CVMod} is a genuine universal module (a quotient of the free current module by the
submodule generated by the relations), not an assumed object, and the bounds are upper bounds on it. A presentation with too
many relations would make the bound trivially true but would break \lean{mabRel_zero_cv3}, which is a theorem: the pure-adjoint
member really is implied by the relations of $M_{0,N}$. The two base statements are named \lean{Prop}s here; the next theorem
proves them.
\end{dictionary}

\begin{physics}
\emph{Heuristic.} Each member of the family is a candidate Hilbert space for a chain of $\mathfrak{sl}_3$ ``particles'' of four
kinds (the three-dimensional fundamental, the six- and ten-dimensional symmetric tensors, and the adjoint) fused at a single
point, with the Chari--Venkatesh relations encoding how much energy each kind can carry along each root direction. Enlarging the
set of particle species is what makes an inductive count possible: removing one adjoint particle produces configurations with
lighter species, which the smaller pure-adjoint family cannot describe.
\end{physics}

\begin{status}
Definitions (named \lean{Prop}s, not axioms). \lean{mabRel_zero_cv3} and \lean{sat_cvSpec_iff}: proved, standard axioms
(campaign record, targeted builds; not re-run here).
\end{status}

\begin{theorem}[The $A_2$ total bound]\label{l26:thm:cvtotal}
\leavevmode
\begin{enumerate}
\item The two base families hold: \lean{CVUniv.sharpRectangularBases} and \lean{CVUniv.residualTheta3Bound}.
\item If both base families hold, then \lean{FusionTotalDim}: for every $m$, every current module and every vector $w$
  satisfying the local Weyl and fusion relations for $m\theta$, $\sum_{b<B}\dim\ltwentysixcyc(w)_b\le8^m$ for all $B$.
\item Hence \lean{FusionTotalDim} holds, with no hypothesis.
\end{enumerate}
\end{theorem}
\noindent Files \leanpath{VOA/\allowbreak DeligneAudit/\allowbreak DeligneFusionCVRectFilteredRectangular.lean} (1, first
family), \leanpath{VOA/\allowbreak DeligneAudit/\allowbreak DeligneFusionCVResidualThetaBoundFinal.lean} (1, second family),
\leanpath{VOA/\allowbreak DeligneAudit/\allowbreak DeligneFusionCVTotal.lean} (2) and
\leanpath{VOA/\allowbreak DeligneAudit/\allowbreak A2Unconditional.lean} (3).
\begin{leancode}
theorem sharpRectangularBases : SharpRectangularBases := by

theorem residualTheta3Bound : ResidualTheta3Bound := by

theorem fusionTotalDim_of_cvBases (hR : SharpRectangularBases)
    (hΘ : ResidualTheta3Bound) : FusionTotalDim :=

theorem CVUniv.fusionTotalDim : SectorWeyl.FusionTotalDim :=
  CVUniv.fusionTotalDim_of_cvBases CVUniv.sharpRectangularBases CVUniv.residualTheta3Bound
\end{leancode}

\begin{dictionary}
\textbf{Names.} The first three declarations are in the namespace \lean{VOA.DeligneAudit.A2.CVUniv}; the fourth is declared in
\lean{VOA.DeligneAudit.A2} with the dotted name \lean{CVUniv.fusionTotalDim}, so its full name is
\lean{VOA.DeligneAudit.A2.CVUniv.fusionTotalDim}. Its one-line proof is shown because it is informative: it is literally (2)
applied to the two terms of (1).

\textbf{The conclusion.} \lean{SectorWeyl.FusionTotalDim} is Definition~\ref{l26:def:mab}(b): a universally quantified
statement over \emph{all} current modules \lean{V} and vectors \lean{w} with \lean{LocalWeylRel ρ w m}, bounding the sum of the
first $B$ degree pieces of the cyclic span by $8^m$. Because the degree pieces are finite-dimensional
(Theorem~\ref{l26:thm:finite}), this is the bound $\dim F_m\le8^m$ of the text. Together with the free lower bound (the fusion
product $V(\theta)^{*m}$ of dimension $8^m$ is a quotient of $F_m$), it gives $\dim F_m=8^m$ on paper; the equality is not itself
stated in Lean.

\textbf{Hypotheses.} None for (1) and (3). The internal definitions they use (\lean{CVMod}, \lean{cvSpec},
\lean{bracketTable}) cannot weaken the conclusion of (3), because that conclusion mentions none of them: it is the previously
defined \lean{FusionTotalDim}.
\end{dictionary}

\textbf{Proof on paper.} (2) For fixed $m$, the vector $w$ satisfies $\mathrm{CV}(0,0,0;m)$ (\lean{mabRel_zero_cv3}), so its cyclic span
is a quotient of $C_m(0,0,0)$ degree by degree, and it suffices to show $\dim C_m(0,0,0)\le8^m$. One proves the stronger
$\dim C_j(\varepsilon,\nu,\rho)\le\mathrm{cvCount}(\varepsilon,\nu,\rho;j)$ for $\varepsilon\in\{(0,0),(1,0),(0,1)\}$ and all
$\nu,\rho,j$, by strong induction. The induction step is a three-term dimension inequality of Ravinder type. Killing the
top generators along the simple roots maps $C_j(\varepsilon,\nu,\rho)$ onto a member of the family in which adjoint factors are
traded for lighter ones (the \emph{quotient relations}), and the kernel of such a map is a cyclic module whose generator satisfies the
relations of a further, smaller member (the \emph{kernel relations}). Rank--nullity bounds $\dim C_j(\varepsilon,\nu,\rho)$ by the sum
of the dimensions of the smaller members, and the counts are such that the sum is exactly $\mathrm{cvCount}$; for two adjoints, for
example, the recurrence closes as $64=36+18+10$. The induction ends at $j=0$, on the rectangular base family, and, in one residual
branch, on $C_1(0,0,\rho)$; these are the two base families. The sums are taken degree by degree in the universal module before
mapping to an arbitrary model, so no independence of the pieces of the model is used.
(1) The rectangular bound is proved through the associated graded of the PBW filtration of the universal module: the CV relations
give polynomial relations among commuting symbols, and an explicit spanning set of the quotient is counted; its size is
$3^{|\varepsilon|}6^{|\nu|}10^{|\rho|}$. The residual family is bounded by the same kind of argument, with explicit quotient maps
that split the adjoint site ($8=5+3$, then $5=3+2$). (3) Compose. \qed

\begin{physics}
\emph{Heuristic.} The statement is that $m$ adjoint ``particles'' of $\widehat{\mathfrak{sl}}_3$ fused at one point have exactly
$8^m$ states, the number they have when kept apart: fusion redistributes energy but creates or destroys no states. For the
Deligne chain this is the input that bounds the number of operators of $L_{-3/2}(\mathfrak{sl}_3)$ in each $m$-particle sector
(Chapter~\ref{ch:L24}).
\end{physics}

\begin{status}
\lean{CVUniv.fusionTotalDim}: \texttt{\#print axioms} prints \lean{[propext, Classical.choice, Quot.sound]} per the targeted
build record (27 October 2026: \lean{VOA.DeligneAudit.A2Unconditional}, 9755 jobs, green); the base families
\lean{sharpRectangularBases} and \lean{residualTheta3Bound} likewise (targeted builds of 8863 and 8887 jobs). We did not re-run
these checks: this checkout has no compiled objects for the new modules. The aggregate build, which contains the umbrella \lean{VOA.Deligne},
was green on 2026-10-28 in a pre-release build (18{,}901 jobs; 18{,}907 after a merge). \textbf{No hypotheses.}
\end{status}

%======================================================================
\section[The $A_2$ total bound on paper]{The $A_2$ total bound on paper: Feigin's Cor.~2 and the slice lemma}\label{l26:sec:open}
%======================================================================

This section records, at the level of mathematics only, what is known about the last $A_2$ input; we write $D(N)$ for
$D(1,N\theta)=M_{N,0}$. The Lean proof of the bound is the different route of Section~\ref{l26:sec:cv}; nothing in this
section is in Lean beyond the reductions of Section~\ref{l26:sec:a2red}. We keep the analysis because it explains why the
published proof needed repair, and because the slice lemma below would give a second, more direct proof.

\paragraph{The statement.} $\dim V(\theta)^{*N}=8^N$ for $\mathfrak{sl}_3$, i.e.\ $\dim M_{0,N}\le8^N$ (the lower bound is
free: the fusion product is a quotient of $M_{0,N}$ of dimension $8^N$). Equivalently, with $D(N)=M_{N,0}$:
\begin{equation*}
M_{0,N}\cong D(N)/(\ltwentysixg\otimes t^N\C[t])D(N),\qquad\text{so}\qquad \dim(\ltwentysixg\otimes t^N\C[t])D(N)\ge9^N-8^N .
\end{equation*}
This is Feigin's Cor.~2 (arXiv:0806.4851) and Kus--Littelmann's Cor.~3.6(3) at level one, and exactly \lean{IdealLower}.

\paragraph{Feigin's argument and its gap.} Feigin realizes $D(N)\cong D(1)^{*N}$ (Fourier--Littelmann for simply laced types) as a fusion product
$T=\bigotimes_jD(1)(c_j)$, $D(1)=\ltwentysixg\oplus\C$, at distinct points $c_j$, and asserts (Prop.~1) that the associated
graded of $T$ for the $t^N$-filtration is a bigraded fusion product in which a second degree $u$ counts trivial factors.
In $T$ the image of the ideal $\ltwentysixg\otimes f(t)\C[t]$, $f=\prod(t-c_j)$, is exactly $K=$ the tensors with at least one
trivial factor, and $T/K\cong\ltwentysixg^{\otimes N}$. The bound needs $\operatorname{gr}K\subseteq(t^N)\cdot\operatorname{gr}T$. The
natural mechanism, that the fusion filtration is homogeneous for the $u$-grading, is \textbf{false}: for $N=2$ the degree-one
filtration piece has dimension $55$ while the sum of the dimensions of its $u$-projections is $63$; for $N=3$ the pairs are
$(161,188)$, $(313,403)$, \dots. Feigin's printed proof of Prop.~1 (labelled ``the idea of the proof'') is a change of current
operators followed by ``Proposition is proved''; it is a sketch, and a rigorous Prop.~1 is equivalent in strength to the statement itself.
Ravinder (arXiv:1409.0274, Prop.~5 in the v4 numbering) takes his terminal case from Feigin. The repair and the alternative route recorded in
the Warning of Section~\ref{l26:sec:math} (Kus--Littelmann 2015; Setia--Rani--Khandai, arXiv:2405.11944) are published. The library formalizes the Setia--Rani--Khandai route (Section~\ref{l26:sec:cv}).

\paragraph{The slice lemma.} Let $T_{a,b}=D(1)(c_1)\otimes\dots\otimes D(1)(c_a)\otimes\ltwentysixg(c_{a+1})\otimes\dots\otimes\ltwentysixg(c_{a+b})$,
cyclic on $w=e_\theta^{\otimes(a+b)}$ (on $\ltwentysixg$-sites $x(1)$ acts by $0$), with fusion filtration $\Phi^{(a,b)}$. Let
$S=\ker(T_{a,b}\to T_{a-1,b+1})$, the tensors with the trivial vector at site $a$, so $S\cong T_{a-1,b}\otimes1$; put
$\mathrm{sh}=2a+b-1$. The \textbf{slice lemma (SL)} is the equality of subspaces
\begin{equation}\label{l26:eq:SL}
\Phi^{(a,b)}_d\cap S=\Phi^{(a-1,b)}_{d-\mathrm{sh}}\otimes1\qquad\text{for all }d .
\end{equation}
\begin{itemize}
\item The inclusion $\supseteq$ \textbf{is proved on paper}: with $g_a(t)=(t-c_a)\prod_{j<a}(t-c_j)^2\prod_{j>a}(t-c_j)$, of
  degree $\mathrm{sh}$ and $\equiv t^{\mathrm{sh}}$ modulo lower terms, $f_\theta(g_a)w$ is a nonzero multiple of
  $e_\theta^{\otimes(a-1)}\otimes1\otimes e_\theta^{\otimes b}$ (all other sites see a double or simple zero), and every current
  kills the trivial factor, so words act through the smaller model.
\item The inclusion $\subseteq$ is \textbf{open}. It is equivalent to the Hilbert-series recursion
  $H_{a,b}=H_{a-1,b+1}+q^{2a+b-1}H_{a-1,b}$ for fusion products, i.e.\ Ravinder's exact sequence realized in a model, and it would
  give \lean{RavinderKernelLower} in the associated-graded model, hence a second proof of \lean{FusionTotalDim}.
\item \textbf{Numerics} (exact, modulo $2^{61}-1$): SL holds for all $(a,b)$ with $a+b\le4$ in every degree; the stronger
  claim $K\cap\Phi_d\subseteq\ltwentysixg(f)\Phi_{d-N}+\Phi_{d-1}$ holds for $N\le4$ (for $N=4$ the dimensions of
  $\operatorname{gr}K$ total $2465=6561-4096$).
\end{itemize}
\paragraph{The lattice form of the slice problem.} In $V_{A_2}$, with $v_i=S_i(-\theta)e^{(N-1)\theta}=\pm f_\theta(2N-1-i)e^{N\theta}$
(the $S_i$ are the Schur polynomials in the Heisenberg modes) and $C^{\le i}$ the cyclic span of $v_0,\dots,v_i$, the successive quotients have dimensions $9^{N-1-i}8^i$ for $N\le4$
($729,648,576,512$ at $N=4$). A proof would need a functional separating these pieces; charge and conformal-weight windows
separate them for $N=2$ but not for $N=3$.

\begin{warning}
\textbf{Approaches to the $A_2$ total bound that fail}, for the record (each with the counterexample that rules it out).
\begin{enumerate}
\item \emph{The naive truncated box.} Imposing $f_1(k)uw=0$ for $k\ge m$ and box bounds overcounts: $72$ vs $64$ at $m=2$,
  $656$ vs $512$ at $m=3$; already for $\mathfrak{sl}_2$, $\lambda=3$ with $F(k)y=0$ for $k\ge2$ gives $7$ instead of $6$: the
  box monomials with indices $i<2$ and $i+\ell\le3$ are $1$; $z_0,z_1$; $z_0^2,z_0z_1,z_1^2$; $z_0^3$, while the true dimension is
  $\prod(\xi_j+1)=6$ for $\xi=\mathrm{bal}(3,2)=(2,1)$ (Theorem~\ref{l26:thm:cv}).
\item \emph{Short identities.} The truncation identities in $M_{0,m}$ are a normal-form problem for the whole annihilator
  ideal; a single-letter recursion reproduces the rule for $m\le3$ but fails for $28$ of $256$ cases at $m=4$.
\item \emph{The presentation route} (``$M_{0,m}$ is the associated graded of the fusion model'') is circular: the
  surjection is free, its injectivity is the statement.
\item \emph{Plain tensor induction} $T_N=T_{N-1}\otimes D(1)$ fails (fusion associativity): the image of the filtration in
  $T_{N-1}\otimes\ltwentysixg$ is the mixed filtration, not a shifted tensor filtration; one must use the mixed models
  $T_{a,b}$.
\item \emph{Separating functionals} in $V_{A_2}$ by charge, $L_0$-weight or $\theta$-degree fail from $N=3$.
\item \emph{Leading-term independence} of the Chari--Loktev words against Fock monomials fails ($77/81$ at $N=2$,
  $547/729$ at $N=3$ under every tested order).
\item \emph{Ungraded counting} in the evaluation model is useless: at nonzero points $\ltwentysixg\otimes t^N\C[t]$ acts as
  all of $\ltwentysixgt$.
\end{enumerate}
\end{warning}

%======================================================================
\section{The type-generic fusion layer}\label{l26:sec:generic}
%======================================================================

The $A_2$ argument of Sections~\ref{l26:sec:squeeze}--\ref{l26:sec:univ} is redone for an arbitrary sector
$S$ (Definition~\ref{l24:def:sectordata}): $S$ supplies $n=\dim\ltwentysixg$, structure constants $S.f$ on $\{0,\dots,n-1\}$,
the index $c_0$ of $f_\theta$, predicates \lean{raising}, \lean{cartan}, \lean{cartanN}, \lean{lowering}, \lean{fusion},
the values $\theta(h)$ (\lean{thetaH}, and \lean{thetaN} for Cartan basis vectors with $\theta(h)<0$), and lowering exponents
\lean{lowExp}$=\theta(h_i)$. \lean{LocalWeylRelG S ρ w m} and \lean{FusionModuleBoundG S F} are quoted in
Definition~\ref{l24:def:U5}.

\begin{definition}[Relation specifications and type-generic universal modules]\label{l26:def:relspec}
A \emph{relation specification} $R$ consists of a set of \emph{tail words} $t$ (relations $t\cdot w=0$) and a set of
\emph{eigen-relations} $(h,\lambda)$ (relations $h(0)w=\lambda w$). A vector $w$ \emph{satisfies} $R$ if all these hold. The
universal module $U(R)$ is the free space on words modulo all prefixed commutation instances and prefixed relations of $R$;
$M^{\mathrm{univ}}(R)(q)$ is the character of the cyclic module of its generator. $R$ is \emph{finite} if every
$\ltwentysixcyc_{U(R)}(w_U)_N$ is finite-dimensional.
\end{definition}
\noindent File \leanpath{VOA/\allowbreak DeligneAudit/\allowbreak DeligneFusionUnivG.lean}.
\begin{leancode}
structure RelSpec (n : ℕ) where
  tails : Set (List (Ltr n))
  eig : Set (Fin n × ℂ)

structure Sat {V : Type*} [AddCommGroup V] [Module ℂ V] (ρ : CurrentModule f V) (w : V)
    (R : RelSpec n) : Prop where
  tail : ∀ t ∈ R.tails, wordOp ρ t w = 0
  eig : ∀ p ∈ R.eig, ρ.op p.1 0 w = p.2 • w

def relGen (R : RelSpec n) : Set (FreeW n) :=
  {g | ∃ p : List (Ltr n), (∃ s x y, g = commInst f p s x y) ∨ (∃ t ∈ R.tails, g = sw (p ++ t)) ∨
    (∃ q ∈ R.eig, g = sw (p ++ [(q.1, 0)]) - q.2 • sw p)}

def Jrel (R : RelSpec n) : Submodule ℂ (FreeW n) := Submodule.span ℂ (relGen f R)

abbrev UModG (R : RelSpec n) := FreeW n ⧸ Jrel f R

/-- **The universal module is a current module.** -/
def univRepG (R : RelSpec n) : CurrentModule f (UModG f R) where
  op c r := opU f R (c, r)

def wUG (R : RelSpec n) : UModG f R := Submodule.Quotient.mk (sw [])

def UnivFD (R : RelSpec n) : Prop :=
  ∀ N, FiniteDimensional ℂ (cycDeg (univRepG f R) (wUG f R) N)

def McharG (R : RelSpec n) : PowerSeries ℚ :=
  PowerSeries.mk fun N => (Module.finrank ℂ (cycDeg (univRepG f R) (wUG f R) N) : ℚ)

theorem finrank_le_McharG {R : RelSpec n} (hfd : UnivFD f R) {V : Type} [AddCommGroup V]
    [Module ℂ V] {ρ : CurrentModule f V} {w : V} (h : Sat f ρ w R) (N : ℕ) :
    (Module.finrank ℂ (cycDeg ρ w N) : ℚ) ≤ PowerSeries.coeff N (McharG f R) := by
\end{leancode}

\begin{dictionary}
\lean{f : Fin n → Fin n → Fin n → ℂ} is a section variable (arbitrary structure constants; no Jacobi identity is needed for
the construction). \lean{commInst f p s x y} is as in Definition~\ref{l26:def:umod} with \lean{f} in place of
\lean{bracketTable}. \lean{finrank_le_McharG} is the universal bound: a module whose generator satisfies $R$ has degree pieces
no larger than those of $U(R)$, provided $U(R)$ is finite (\lean{hfd}); without \lean{hfd} the inequality could fail only through
the convention \lean{finrank}$=0$ on infinite-dimensional spaces.
\end{dictionary}

\begin{physics}
The same ``largest Hilbert space compatible with given null vectors'', now for any Lie algebra in the Deligne series, with the
null vectors given as data.
\end{physics}

\begin{status}
Definitions and theorem; \texttt{\#print axioms} standard (record).
\end{status}

\begin{definition}[The local Weyl specification; the total bound; U7, type-generic]\label{l26:def:lwspec}
\lean{lwSpec S m} is the specification of the local Weyl and fusion relations for $m\theta$. \textbf{\lean{FusionTotalBoundG S}}:
for every $m$ and every module with these relations, the degree pieces are finite-dimensional and the total dimension is
$\le n^m$. \textbf{\lean{DemazureRealizationG S spec Dem n}}: some module whose generator satisfies \lean{spec n 0} has
$\dim\ltwentysixcyc(w)_N\ge[q^N]\mathrm{Dem}_n$ for every $N$.
\end{definition}
\noindent Files \leanpath{VOA/\allowbreak DeligneAudit/\allowbreak DeligneFusionUnivGSectorRelations.lean} (\lean{lwSpec},
\lean{localWeylRelG_iff}) and \leanpath{VOA/\allowbreak DeligneAudit/\allowbreak DeligneFusionUnivGSector.lean}
(\lean{FusionTotalBoundG}, \lean{DemazureRealizationG}).
\begin{leancode}
def lwSpec (m : ℕ) : RelSpec S.n where
  tails := {t | (∃ e r, S.raising e ∧ t = [(e, r)]) ∨ (∃ h r, S.cartan h ∧ 1 ≤ r ∧ t = [(h, r)]) ∨
    (∃ i, S.lowering i ∧ t = List.replicate (m * S.lowExp i + 1) (i, 0)) ∨
    (∃ x, S.fusion x ∧ t = [(x, m)]) ∨ (∃ h r, S.cartanN h ∧ 1 ≤ r ∧ t = [(h, r)])}
  eig := {q | (∃ h, S.cartan h ∧ q = (h, ((m * S.thetaH h : ℕ) : ℂ))) ∨
    (∃ h, S.cartanN h ∧ q = (h, -((m * S.thetaN h : ℕ) : ℂ)))}

theorem localWeylRelG_iff {V : Type*} [AddCommGroup V] [Module ℂ V] (ρ : CurrentModule S.f V)
    (w : V) (m : ℕ) : LocalWeylRelG S ρ w m ↔ Sat S.f ρ w (lwSpec S m) := by

def FusionTotalBoundG : Prop :=
  ∀ (m : ℕ) (V : Type) [AddCommGroup V] [Module ℂ V] (ρ : CurrentModule S.f V) (w : V),
    LocalWeylRelG S ρ w m → (∀ b, FiniteDimensional ℂ (cycDeg ρ w b)) ∧
      ∀ B : ℕ, ∑ b ∈ Finset.range B, Module.finrank ℂ (cycDeg ρ w b) ≤ S.n ^ m

def DemazureRealizationG (spec : ℕ → ℕ → RelSpec S.n) (Dem : ℕ → PowerSeries ℚ) (n : ℕ) : Prop :=
  ∃ (V : Type) (_ : AddCommGroup V) (_ : Module ℂ V) (ρ : CurrentModule S.f V) (w : V),
    Sat S.f ρ w (spec n 0) ∧ ∀ N, coeff N (Dem n) ≤ (Module.finrank ℂ (cycDeg ρ w N) : ℚ)
\end{leancode}

\begin{dictionary}
\lean{S : DeligneSectorData} is a section variable. In \lean{lwSpec}: raising letters at every grade; Cartan letters at
positive grade; $f_i(0)^{m\theta(h_i)+1}$ for \lean{lowering} letters; $x(m)$ for every \lean{fusion} letter; and the eigen-relations
$h(0)w=m\theta(h)w$ (and $=-m\theta_N(h)w$ for \lean{cartanN}). \lean{localWeylRelG_iff} proves this is exactly
\lean{LocalWeylRelG}. For $E_6$, $E_7$, $E_8$ the \lean{fusion} set is Ravinder's
$\{f_\alpha:\langle\alpha,\theta^\vee\rangle=1\}\cup\{f_\theta\}$; for $A_2$, $D_4$, $G_2$, $F_4$ it is all negative root vectors. The
dictionary of Definition~\ref{l24:def:U5} explains why both readings are sound for U5: more relations make the hypothesis weaker.
\lean{DemazureRealizationG} takes \lean{spec} and \lean{Dem} as parameters; in all uses \lean{spec} is Ravinder's family
\lean{ravFam S} (Section~\ref{l26:sec:kernelG}) and \lean{Dem} the Demazure series of Chapter~\ref{ch:L25}.

\textbf{Non-simply-laced caution.} For $G_2$ and $F_4$, $\ltwentysixWloc(m\lambda)\neq D(1,m\lambda)$ in general (Naoi). Along
multiples of $\theta$ (long, fundamental for $G_2$ and $F_4$) the project's numerics find no discrepancy: $\sum_N\dim D(1,n\theta)_N=15^n$
for $G_2$ and $n\le18$, and inverting the mixed recursion gives nonnegative $\ltwentysixF_m$ with $\ltwentysixF_m(1)=14^m$ for $m\le18$.
Ravinder's presentation of $D(1,n\theta)$ (his Prop.~1) and his exact sequences are stated for every simple $\ltwentysixg$ (only
his Thm.~2 and Cor.~3 assume simply laced), and Kus--Littelmann's Thm.~4.2 gives the total bound for $G_2$, $F_4$ as for every other
type. The figures above are numerical consistency checks, not proofs.
\end{dictionary}

\begin{warning}
\textbf{The fusion set convention.} For $A_2$, $D_4$, $G_2$, $F_4$ the sectors' \lean{fusion} predicate contains all negative root
vectors; for $E_6$, $E_7$, $E_8$ it is Ravinder's set. This makes \lean{LocalWeylRelG}, hence \lean{FusionTotalDimG}, differ in
form between types, but not in content. Both readings are sound for the chain (Definition~\ref{l24:def:U5}), and they impose the
same relations: Ravinder's condition is $\langle\alpha,\theta^\vee\rangle=1$, which includes the short roots of $G_2$ and $F_4$ (for which
$(\theta,\alpha^\vee)\ge2$), and the remaining positive roots, with $\langle\alpha,\theta^\vee\rangle=0$, satisfy $f_\alpha(r)w=0$ for all $r$
by the $h(1)$-bracket argument (Remark~\ref{l26:rem:fusionset}; Chapter~\ref{ch:L27}, Lemma~\ref{l27:lem:truncweyl}). An earlier
version of this Warning said that these extra relations were not known to hold in Ravinder's $F_m$; that was wrong.
\end{warning}

\begin{physics}
(Heuristic.) $m$ adjoint particles of any Deligne-series flavour symmetry, fused; the total bound says the fused space is no larger than
$\ltwentysixg^{\otimes m}$.
\end{physics}

\begin{status}
\lean{FusionTotalBoundG} and \lean{DemazureRealizationG}: \textbf{named inputs}. The first is reduced to
\lean{FusionTotalDimG} for all seven types (Section~\ref{l26:sec:finiteG}); the second is proved for $A_2$ (through the
$A_2$-specific \lean{DemazureRealization}, which is the same statement for \lean{bracketTable}) and open for the other six,
with the relations half proved on lattices for $D_4$, $E_6$, $E_7$, $E_8$ (Section~\ref{l26:sec:lattice}).
\end{status}

\begin{theorem}[The generic squeeze and the count]\label{l26:thm:squeezeG}
Let $\mathrm{Dem}_n$ be power series and $d\in\N$. (1) If each $\mathrm{Dem}_n$ is a polynomial with $\mathrm{Dem}_n(1)=(d+1)^n$,
then each $\mathrm{fusionOf}(\mathrm{Dem})_m$ is a polynomial with value $d^m$ at $q=1$. (2) If $M$ satisfies the generic squeeze
inputs (nonneg, (R), $\mathrm{Dem}_n\le M_{n,0}$, total $\le d^m$) and the count of (1) holds for \lean{fusionOf Dem}, then
$M_{0,m}=\mathrm{fusionOf}(\mathrm{Dem})_m$. (3) If in addition every module with \lean{LocalWeylRelG S} for $m\theta$ is bounded by
$M_{0,m}$ and $n=d$, then U5 holds for \lean{fusionOf Dem}.
\end{theorem}
\noindent File \leanpath{VOA/\allowbreak DeligneAudit/\allowbreak DeligneFusionSqueezeG.lean}.
\begin{leancode}
theorem fusionOf_count (Dem : ℕ → PowerSeries ℚ) (d : ℕ)
    (hcount : ∀ n, ∃ B, (∀ N, B ≤ N → coeff N (Dem n) = 0) ∧
      ∑ N ∈ Finset.range B, coeff N (Dem n) = ((d : ℚ) + 1) ^ n) (m : ℕ) :
    ∃ B, (∀ b, B ≤ b → coeff b (fusionOf Dem m) = 0) ∧
      ∑ b ∈ Finset.range B, coeff b (fusionOf Dem m) = (d : ℚ) ^ m := by

structure FusionSqueezeInputsG (Dem : ℕ → PowerSeries ℚ) (d : ℕ) (M : ℕ → ℕ → PowerSeries ℚ) :
    Prop where
  nonneg : ∀ a b N, 0 ≤ coeff N (M a b)
  rec_le : ∀ a b, CoeffLe (M (a + 1) b) (M a (b + 1) + X ^ (2 * a + b + 1) * M a b)
  dem_le : ∀ n, CoeffLe (Dem n) (M n 0)
  total_le : ∀ m B, ∑ N ∈ Finset.range B, coeff N (M 0 m) ≤ (d : ℚ) ^ m

theorem squeezeG (h : FusionSqueezeInputsG Dem d M)
    (hcount : ∀ m, ∃ B, (∀ b, B ≤ b → coeff b (fusionOf Dem m) = 0) ∧
      ∑ b ∈ Finset.range B, coeff b (fusionOf Dem m) = (d : ℚ) ^ m) (m : ℕ) :
    M 0 m = fusionOf Dem m := by

theorem fusionModuleBoundG_of_squeeze (S : DeligneSectorData) (hd : S.n = d)
    (h : FusionSqueezeInputsG Dem d M)
    (hcount : ∀ m, ∃ B, (∀ b, B ≤ b → coeff b (fusionOf Dem m) = 0) ∧
      ∑ b ∈ Finset.range B, coeff b (fusionOf Dem m) = (d : ℚ) ^ m)
    (huniv : ∀ (m : ℕ) (V : Type) [AddCommGroup V] [Module ℂ V] (ρ : CurrentModule S.f V) (w : V),
      LocalWeylRelG S ρ w m → ∀ b, (Module.finrank ℂ (cycDeg ρ w b) : ℚ) ≤ coeff b (M 0 m)) :
    FusionModuleBoundG S (fusionOf Dem) := by
\end{leancode}

\begin{dictionary}
\lean{Dem}, \lean{d}, \lean{M} are section variables in \lean{squeezeG}. The count is phrased with an explicit support bound
\lean{B}: coefficients vanish from \lean{B} on and the sum below \lean{B} is the stated power. In \lean{fusionModuleBoundG_of_squeeze},
\lean{hd : S.n = d} ties the bound to $\dim\ltwentysixg$. The proofs are the $A_2$ ones verbatim, with $8$ replaced by $d$; the
count (1) is binomial inversion: evaluating the recursion at $q=1$ gives $(d+1)^n=\sum_k\binom nk\mathrm{fusionOf}_{n-k}(1)$.
\end{dictionary}

\textbf{Proof on paper.} As Theorem~\ref{l26:thm:squeeze}; for (1), induction with $\ltwentysixqb nk\big|_{q=1}=\binom nk$ and
the binomial theorem. \qed

\begin{physics}
(Heuristic.) Same counting argument as for $A_2$, for every Deligne flavour symmetry.
\end{physics}

\begin{status}
Proved; \texttt{\#print axioms} standard (\lean{squeezeG} re-run). The count hypothesis is discharged per type through the
Demazure totals of Chapter~\ref{ch:L25} (Theorem~\ref{l25:thm:total}), or avoided by the unclamped series used in
Chapter~\ref{ch:L27}.
\end{status}

\begin{theorem}[Generic exact-sequence step; U5 from the kernel lemma]\label{l26:thm:kernelstep}
Let $R_1,R',R_0$ be specifications and $k$ a word such that every tail of $R'$ is a tail of $R_1$ or ends in $k$, and the
eigen-relations of $R'$ are among those of $R_1$ (\lean{SpecStep R₁ R' k}). If $k\cdot w_U$ satisfies $R_0$ in $U(R_1)$, then
\begin{equation*}
\dim\ltwentysixcyc_{U(R_1)}(w)_N\le\dim\ltwentysixcyc_{U(R')}(w)_N+[\deg k\le N]\dim\ltwentysixcyc_{U(R_0)}(w)_{N-\deg k},
\end{equation*}
provided the two right-hand pieces $\ltwentysixcyc_{U(R')}(w)_N$ and $\ltwentysixcyc_{U(R_0)}(w)_{N-\deg k}$ are finite-dimensional.
For a family \lean{spec} with $\mathrm{spec}(0,m)=$ \lean{lwSpec S m}, kernel words $\mathrm{kw}(a,b)$ of degree $2a+b+1$ with
\lean{SpecStep (spec (a+1) b) (spec a (b+1)) (kw a b)}, the per-type kernel lemma \lean{RavinderKernelRelG}, the total bound, U7 and
the count, U5 holds for \lean{fusionOf Dem}.
\end{theorem}
\noindent File \leanpath{VOA/\allowbreak DeligneAudit/\allowbreak DeligneFusionUnivGGrade.lean}.
\begin{leancode}
structure SpecStep (R R' : RelSpec n) (k : List (Ltr n)) : Prop where
  tail : ∀ t ∈ R'.tails, t ∈ R.tails ∨ ∃ p, t = p ++ k
  eig : R'.eig ⊆ R.eig

theorem finrank_stepG (hs : SpecStep R₁ R' k) (hK : Sat f (univRepG f R₁) (vKG f k R₁) R₀)
    (N : ℕ) [FiniteDimensional ℂ (cycDeg (univRepG f R') (wUG f R') N)]
    [FiniteDimensional ℂ (cycDeg (univRepG f R₀) (wUG f R₀) (N - deg k))] :
    Module.finrank ℂ (cycDeg (univRepG f R₁) (wUG f R₁) N) ≤
      Module.finrank ℂ (cycDeg (univRepG f R') (wUG f R') N) +
        if deg k ≤ N then
          Module.finrank ℂ (cycDeg (univRepG f R₀) (wUG f R₀) (N - deg k)) else 0 := by

def RavinderKernelRelG (spec : ℕ → ℕ → RelSpec n) (kw : ℕ → ℕ → List (Ltr n)) : Prop :=
  ∀ (a b : ℕ) (V : Type) [AddCommGroup V] [Module ℂ V] (ρ : CurrentModule f V) (w : V),
    Sat f ρ w (spec (a + 1) b) → Sat f ρ (wordOp ρ (kw a b) w) (spec a b)

theorem fusionModuleBoundG_of_kernel {Dem : ℕ → PowerSeries ℚ}
    (hspec : ∀ m, spec 0 m = lwSpec S m)
    (hs : ∀ a b, SpecStep (spec (a + 1) b) (spec a (b + 1)) (kw a b))
    (hdeg : ∀ a b, deg (kw a b) = 2 * a + b + 1)
    (hF : FusionTotalBoundG S) (hK : RavinderKernelRelG S.f spec kw)
    (hD : ∀ n, DemazureRealizationG S spec Dem n)
    (hcount : ∀ m, ∃ B, (∀ b, B ≤ b → coeff b (fusionOf Dem m) = 0) ∧
      ∑ b ∈ Finset.range B, coeff b (fusionOf Dem m) = (S.n : ℚ) ^ m) :
    FusionModuleBoundG S (fusionOf Dem) :=
\end{leancode}

\begin{dictionary}
\lean{deg} is local notation for \lean{SectorGraded.wordDeg}, the total $t$-degree of a word. Section variables:
\lean{f} and \lean{k} (and implicit \lean{R₁ R₀ R'}) for \lean{finrank_stepG}; \lean{S}, \lean{spec}, \lean{kw} (implicit) for
\lean{fusionModuleBoundG_of_kernel}. The two bracketed \lean{[FiniteDimensional ...]} binders of \lean{finrank_stepG} are
instance hypotheses: finite-dimensionality of the two right-hand pieces, exactly as in the $A_2$ step (Theorem~\ref{l26:thm:step}).
They are needed because \lean{Module.finrank} is $0$ on an infinite-dimensional space, so without them the right-hand side
could be $0$ for a spurious reason and rank--nullity would not apply. In \lean{fusionModuleBoundG_of_kernel} they are discharged
by \lean{univFD_all}. \lean{RavinderKernelRelG f spec kw} is \emph{the} per-type input: if $w$ satisfies
$\mathrm{spec}(a+1,b)$ then $\mathrm{kw}(a,b)\cdot w$ satisfies $\mathrm{spec}(a,b)$. Every \lean{RelSpec} is homogeneous (tails are
single words, eigen-relations use grade-zero letters), so the degree projection needed for the step exists with no extra
hypothesis (\lean{projFG_J}). The finiteness of all $U(\mathrm{spec}(a,b))$ is derived (\lean{univFD_all}), not assumed.
\end{dictionary}

\textbf{Proof on paper.} As Theorem~\ref{l26:thm:step}, with \lean{sat_quotG} (killing $k\cdot w$ turns $R_1$ into $R'$) in place of
\lean{mabRel_quot}. For U5: finiteness of all universal pieces by induction on $a$ (\lean{univFD_all}); (R) from the step; (D)
from U7 and the universal bound; (T) from the total bound at $a=0$; then Theorem~\ref{l26:thm:squeezeG}. \qed

\begin{physics}
(Heuristic.) Ravinder's particle decomposition, with the binding energy $2a+b+1$, for any Deligne flavour symmetry.
\end{physics}

\begin{status}
Proved; \texttt{\#print axioms} standard (\lean{fusionModuleBoundG_of_kernel} re-run). Conditional on the printed hypotheses;
\lean{RavinderKernelRelG} is discharged for every type below.
\end{status}

%======================================================================
\section{Ravinder's kernel lemma for every type}\label{l26:sec:kernelG}
%======================================================================

\begin{definition}[Ravinder's family; the kernel table; integrability]\label{l26:def:ravfam}
\leavevmode
\begin{enumerate}
\item[(a)] For $k,d\in\N$, \lean{ravSpec S k d} is the specification: local Weyl relations for $k\theta$, truncations
$x(k)w=0$ for fusion letters $x\neq f_\theta$, and $f_\theta(d)w=0$.
\item[(b)] \emph{Ravinder's family} is $\mathrm{ravFam}(a,b)=\mathrm{ravSpec}(a+b,2a+b)$, with kernel words
$\mathrm{ravKw}(a,b)=[f_\theta(2a+b+1)]$.
\item[(c)] A \emph{kernel table} (\lean{KernelTable}) is a finite list of structure-constant facts about $S$ (below).
\item[(d)] \textbf{\lean{ThetaNilG S}} is integrability of $f_\theta$: if $w$ has weight $k\theta$, is killed by the raising
zero modes, and $f_i(0)^{k\theta(h_i)+1}w=0$ for the lowering letters, then $f_\theta(0)^{2k+1}w=0$.
\end{enumerate}
\end{definition}
\noindent File \leanpath{VOA/\allowbreak DeligneAudit/\allowbreak D4FusionKernelGeneric.lean}.
\begin{leancode}
def ravSpec (k d : ℕ) : RelSpec S.n where
  tails := {t | (∃ e r, S.raising e ∧ t = [(e, r)]) ∨ (∃ h r, S.cartan h ∧ 1 ≤ r ∧ t = [(h, r)]) ∨
    (∃ i, S.lowering i ∧ t = List.replicate (k * S.lowExp i + 1) (i, 0)) ∨
    (∃ x, S.fusion x ∧ t = [(x, if x = S.c₀ then d else k)]) ∨
    (∃ h r, S.cartanN h ∧ 1 ≤ r ∧ t = [(h, r)])}
  eig := {q | (∃ h, S.cartan h ∧ q = (h, ((k * S.thetaH h : ℕ) : ℂ))) ∨
    (∃ h, S.cartanN h ∧ q = (h, -((k * S.thetaN h : ℕ) : ℂ)))}

def ravFam (a b : ℕ) : RelSpec S.n := ravSpec S (a + b) (2 * a + b)

def ravKw (a b : ℕ) : List (Ltr S.n) := [(S.c₀, 2 * a + b + 1)]

structure KernelTable where
  τ : Fin S.n → ℂ
  tauH : ∀ h, S.cartan h → τ h = (S.thetaH h : ℂ)
  tauN : ∀ h, S.cartanN h → τ h = -(S.thetaN h : ℂ)
  c0_fusion : S.fusion S.c₀
  truncUp : ∀ x, S.fusion x ∨ S.lowering x →
    ∃ h, (S.cartan h ∨ S.cartanN h) ∧ S.f h x x ≠ 0 ∧ ∀ c, c ≠ x → S.f h x c = 0
  raisTheta : ∀ e, S.raising e → ∀ c, S.f e S.c₀ c ≠ 0 →
    (S.cartan c ∨ S.cartanN c) ∨ (S.fusion c ∧ c ≠ S.c₀)
  cartTau : ∀ h, S.cartan h ∨ S.cartanN h → ∀ c, S.f h S.c₀ c = if c = S.c₀ then -τ h else 0
  lowComm : ∀ x, S.fusion x ∨ S.lowering x → ∀ c, S.f x S.c₀ c = 0
  lowExp_le : ∀ i, S.lowering i → S.lowExp i ≤ 1
  simple : ∀ i, S.lowering i → S.lowExp i = 1 → ∃ (e : Fin S.n) (s : ℂ), S.raising e ∧
    (∀ c, S.f e i c ≠ 0 → S.cartan c ∨ S.cartanN c) ∧ ∑ c, s * S.f e i c * τ c = 1 ∧
    ∀ d, ∑ c, s * S.f e i c * S.f c i d = if d = i then -2 else 0
  eθ : Fin S.n
  sθ : ℂ
  eθ_raising : S.raising eθ
  eθ_cartan : ∀ c, S.f eθ S.c₀ c ≠ 0 → S.cartan c ∨ S.cartanN c
  eθ_theta : sθ * ∑ c, S.f eθ S.c₀ c * τ c = 2
  /-- Root vectors killing the generator at grade `0` (orthogonal to `θ`), generated by brackets
  from the simple ones with `θ(h_i) = 0`; `rk` makes the generation well founded. -/
  kill0 : Fin S.n → Prop
  rk : Fin S.n → ℕ
  kill0_low : ∀ x, kill0 x → S.fusion x ∨ S.lowering x
  kill0_gen : ∀ x, kill0 x → (S.lowering x ∧ S.lowExp x = 0) ∨
    ∃ (a b : Fin S.n) (κ : ℂ), κ ≠ 0 ∧ kill0 a ∧ kill0 b ∧ rk a < rk x ∧ rk b < rk x ∧
      ∀ c, S.f a b c = if c = x then κ else 0
  fusionRoot : ∀ x, S.fusion x → x ≠ S.c₀ →
    (∃ β κ, S.raising β ∧ κ ≠ 0 ∧ ∀ c, S.f β S.c₀ c = if c = x then κ else 0) ∨ kill0 x

def ThetaNilG : Prop :=
  ∀ (k : ℕ) (V : Type) [AddCommGroup V] [Module ℂ V] (ρ : CurrentModule S.f V) (w : V),
    (∀ e, S.raising e → ρ.op e 0 w = 0) →
    (∀ h, S.cartan h → ρ.op h 0 w = ((k * S.thetaH h : ℕ) : ℂ) • w) →
    (∀ h, S.cartanN h → ρ.op h 0 w = (-((k * S.thetaN h : ℕ) : ℂ)) • w) →
    (∀ i, S.lowering i → (ρ.op i 0)^[k * S.lowExp i + 1] w = 0) →
    ((ρ.op S.c₀ 0) ^ (2 * k + 1)) w = 0
\end{leancode}

\begin{dictionary}
\textbf{The table, field by field.} \lean{τ} is the signed value $\theta(h)$ on Cartan letters. \lean{c0_fusion}: $f_\theta$ is a
fusion letter. \lean{truncUp}: each fusion or lowering letter $x$ has a diagonal Cartan $h$ with $[h,x]\in\C^\times x$ (used to
propagate a truncation $x(N_0)w=0$ to all higher grades). \lean{raisTheta}: $[e,f_\theta]$ is Cartan or a fusion letter $\neq f_\theta$.
\lean{cartTau}: $[h,f_\theta]=-\theta(h)f_\theta$. \lean{lowComm}: fusion and lowering letters commute with $f_\theta$.
\lean{lowExp_le}: $\theta(h_i)\le1$ for simple lowering letters. \lean{simple}: each simple lowering letter with $\theta(h_i)=1$ has an
$\mathfrak{sl}_2$-partner $e$ with the right normalization. \lean{eθ}, \lean{sθ}: the $\theta$-triple. \lean{kill0}, \lean{rk}: root
vectors orthogonal to $\theta$ that kill $w$ at grade $0$, generated by brackets from simple ones, with a rank making the generation
well founded. \lean{fusionRoot}: each fusion letter $x\neq f_\theta$ is $[e_\beta,f_\theta]$ up to a nonzero scalar, or is in
\lean{kill0}. All fields are decidable statements about integer tables and are checked per type by \lean{decide +kernel}. A
simply-laced version \lean{KernelTableSL} with a converter \lean{KernelTable.ofSL} is used for $D_4$, $E_6$, $E_7$, $E_8$.

\textbf{About \lean{ThetaNilG}.} Its docstring calls it a ``named input''; it is \emph{proved} for every type
(Theorem~\ref{l26:thm:thetanil}). It is the classical integrability of $V(k\theta)$ (Humphreys \S21; Kac Lemma~3.4), which for an
arbitrary current module is not automatic: $f_\theta$ is not simple.
\end{dictionary}

\begin{physics}
The table records how the energy-raising and lowering currents of each flavour algebra commute with the top lowering current
$f_\theta$; integrability says the $\theta$-direction excitations of a highest-weight state of weight $k\theta$ come in a finite
spin-$k$ multiplet.
\end{physics}

\begin{status}
Definitions. \lean{ravFam_step} (\lean{SpecStep}) and \lean{ravKw_deg} ($\deg=2a+b+1$) are proved generically by unfolding.
\end{status}

\begin{theorem}[The generic kernel lemma]\label{l26:thm:kernelG}
For every sector $S$ with a kernel table and \lean{ThetaNilG S}, Ravinder's kernel relation holds: if $w$ satisfies
$\mathrm{ravFam}(a+1,b)$ then $f_\theta(2a+b+1)w$ satisfies $\mathrm{ravFam}(a,b)$.
\end{theorem}
\noindent File \leanpath{VOA/\allowbreak DeligneAudit/\allowbreak D4FusionKernelGeneric.lean}.
\begin{leancode}
theorem ravinderKernelRelG_of_table (T : KernelTable S) (hN : ThetaNilG S) :
    RavinderKernelRelG S.f (ravFam S) (ravKw S) := by
\end{leancode}

\begin{dictionary}
\lean{S} is a section variable. The conclusion is \lean{RavinderKernelRelG} (Theorem~\ref{l26:thm:kernelstep}) for the family
\lean{ravFam S} and kernel words \lean{ravKw S}.
\end{dictionary}

\textbf{Proof on paper.} This is the $A_2$ argument of Theorem~\ref{l26:thm:kernelA2} made generic (\lean{kernel_sat}): raising and
Cartan relations for $v=f_\theta(D+1)w$ by commuting through $f_\theta$ (\lean{raisTheta}, \lean{cartTau}); lowering relations by
$\mathfrak{sl}_2$ algebra (\lean{simple}, \lean{lowComm}); the $f_\theta$-truncation from the $\theta$-Garland identity for the
triple $(s_\theta e_\theta(1),\,s_\theta[e_\theta,f_\theta](m),\,f_\theta(k))$, using \lean{ThetaNilG} for the nilpotency input of
Garland; the other fusion truncations by applying $e_\beta$ to the $\theta$-Garland identity (\lean{fusionRoot}) or directly
(\lean{kill0}). \qed

\begin{physics}
(Heuristic.) As for $A_2$: the state with one top excitation emitted behaves like the ground state of the smaller system.
\end{physics}

\begin{status}
Proved; \texttt{\#print axioms} standard (re-run).
\end{status}

\begin{theorem}[Integrability for the seven types; the kernel lemma unconditionally]\label{l26:thm:thetanil}
\lean{ThetaNilG} holds for $D_4$, $E_6$, $E_7$, $E_8$, $G_2$, $F_4$; hence \lean{RavinderKernelRelG} for \lean{ravFam} holds for
these six types with no hypothesis. Together with the $A_2$ theorem \lean{ravinderKernelRel} (Theorem~\ref{l26:thm:kernelA2}),
the Lean analogue of Ravinder's kernel lemma holds for all seven types.
\end{theorem}
\noindent\textbf{Two caveats on the word ``all seven''.} (a) $A_2$ is formalized differently: its kernel lemma is
\lean{RavinderKernelRel}, phrased with \lean{MabRel} and \lean{bracketTable}, and its integrability input is \lean{ftheta_pow}. No
\lean{ThetaNilG} or \lean{RavinderKernelRelG} statement exists for \lean{A2sector}, so the generic pipeline
\lean{fusionModuleBoundG_of_kernel} is not instantiated at $A_2$; the $A_2$ chain is the specific one of
Sections~\ref{l26:sec:squeeze}--\ref{l26:sec:garland}. (b) For $G_2$ and $F_4$
the theorems below concern the library's own family \lean{ravFam} (fusion set: all negative root vectors). This family imposes the
same relations as Ravinder's (Remark~\ref{l26:rem:fusionset}), and Ravinder's Theorem~1, Corollaries~1--2 and Proposition~1 are
stated for every simple $\ltwentysixg$ (only his Theorem~2 and Corollary~3 assume simply laced); what is not formalized is the
identification of these modules with the Demazure and fusion modules (Section~\ref{l26:sec:generic}).
\noindent File \leanpath{VOA/\allowbreak DeligneAudit/\allowbreak D4FusionKernelTheta.lean}.
\begin{leancode}
theorem sl2_sharpen (E F H : V →ₗ[ℂ] V) (hEF : ∀ x, E (F x) - F (E x) = H x)
    (hHF : ∀ x, H (F x) - F (H x) = (-2 : ℂ) • F x) (u : V) (k : ℕ) (hE : E u = 0)
    (hH : H u = (k : ℂ) • u) {N : ℕ} (hN : F^[N] u = 0) : F^[k + 1] u = 0 := by

theorem d4_thetaNil : ThetaNilG D4sector := by

theorem d4_ravinderKernelRel' : RavinderKernelRelG D4sector.f (ravFam D4sector) (ravKw D4sector) :=
\end{leancode}
\noindent File \leanpath{VOA/\allowbreak DeligneAudit/\allowbreak G2FusionKernel.lean}.
\begin{leancode}
def adL : Module.End ℂ V →ₗ[ℂ] Module.End ℂ V := LinearMap.mulLeft ℂ A - LinearMap.mulRight ℂ A

theorem nil_of_ad {w : V} (hA : A w = 0) {X : Module.End ℂ V} {r p : ℕ}
    (hX : (adL A)^[r + 1] X = 0) (hXp : (X ^ p) w = 0) : (((adL A)^[r] X) ^ p) w = 0 := by
\end{leancode}
\noindent File \leanpath{VOA/\allowbreak DeligneAudit/\allowbreak G2FusionKernelTheta.lean}.
\begin{leancode}
theorem g2_thetaNil : ThetaNilG G2sector := by

theorem g2_ravinderKernelRel : RavinderKernelRelG G2sector.f (ravFam G2sector) (ravKw G2sector) :=
\end{leancode}
\noindent File \leanpath{VOA/\allowbreak DeligneAudit/\allowbreak E6FusionKernelTheta.lean}.
\begin{leancode}
theorem e6_thetaNil : ThetaNilG E6sector := by

theorem e6_ravinderKernelRel' : RavinderKernelRelG E6sector.f (ravFam E6sector) (ravKw E6sector) :=
\end{leancode}
\noindent File \leanpath{VOA/\allowbreak DeligneAudit/\allowbreak E7FusionKernelTheta.lean}.
\begin{leancode}
theorem e7_thetaNil : ThetaNilG E7sector := by

theorem e7_ravinderKernelRel' : RavinderKernelRelG E7sector.f (ravFam E7sector) (ravKw E7sector) :=
\end{leancode}
\noindent File \leanpath{VOA/\allowbreak DeligneAudit/\allowbreak E8FusionKernelTheta.lean}.
\begin{leancode}
theorem e8_thetaNil : ThetaNilG E8sector := by

theorem e8_ravinderKernelRel' : RavinderKernelRelG E8sector.f (ravFam E8sector) (ravKw E8sector) :=
\end{leancode}
\noindent File \leanpath{VOA/\allowbreak DeligneAudit/\allowbreak F4LieKernel.lean}.
\begin{leancode}
theorem f4_thetaNil : ThetaNilG F4sector := by

theorem f4_ravinderKernelRel : RavinderKernelRelG F4sector.f (ravFam F4sector) (ravKw F4sector) :=
\end{leancode}

\begin{dictionary}
\lean{D4sector}, \lean{G2sector}, \dots\ are the sector data of Chapter~\ref{ch:L24} (Theorem~\ref{l24:thm:seven}); the
$E_6$, $E_8$ and $F_4$ declarations live in the namespaces \lean{UnivG.E6K}, \lean{UnivG.E8K}, \lean{UnivG.F4K}, the $D_4$, $E_7$ and $G_2$ ones in \lean{UnivG}.
\lean{adL A T} $=AT-TA$ in $\End(V)$ (section variable \lean{A}). \lean{sl2_sharpen}: section variable $V$; it says a highest-weight
vector of weight $k$ for an $\mathfrak{sl}_2$-triple, on which $F$ is nilpotent, satisfies $F^{k+1}u=0$. The unprimed
\lean{d4_ravinderKernelRel (hN : ThetaNilG D4sector)} etc.\ are the conditional versions; the primed ones (and
\lean{g2_ravinderKernelRel}, \lean{f4_ravinderKernelRel}) are unconditional.
\end{dictionary}

\textbf{Proof on paper.} \emph{Heisenberg chains (simply laced).} For $A_2$, $\theta=\alpha_1+\alpha_2$ and $f_1,f_2,f_\theta$ form a Heisenberg
algebra, so one application of Theorem~\ref{l26:thm:garland}(1) suffices. For $D_4$ there is no such pair of simple roots: only $\alpha_2$
has $(\theta,\alpha_2^\vee)=1$ ($\alpha_2$ the central node, $\theta=\alpha_1+2\alpha_2+\alpha_3+\alpha_4$, and $\alpha_1,\alpha_3,\alpha_4$
orthogonal to $\theta$, so that $f_{\alpha_1},f_{\alpha_3},f_{\alpha_4}$ kill $w$ at grade $0$). The chain builds $\theta$ by adding one
simple root at a time:
\begin{equation*}
\alpha_2\ \to\ \alpha_1+\alpha_2\ \to\ \alpha_1+\alpha_2+\alpha_3\ \to\ \alpha_1+\alpha_2+\alpha_3+\alpha_4\ \to\ \alpha_1+2\alpha_2+\alpha_3+\alpha_4=\theta
\end{equation*}
(in the Lean table basis: $18\to20\to23\to26\to27$, with $f_{\alpha_2}=x_{18}$ and the orthogonal simple roots $x_{16},x_{17},x_{19}$).
At each step: if $f_\gamma^{k+1}w=0$ and $f_\delta^{k+1}w=0$ with $[f_\gamma,f_\delta]=\pm f_{\gamma+\delta}$ central in the pair, the
Heisenberg lemma gives $f_{\gamma+\delta}^{2k+1}w=0$; since $(\theta,(\gamma+\delta)^\vee)=1$ and $w$ is highest weight for the
$\mathfrak{sl}_2$ of $\gamma+\delta$, \lean{sl2_sharpen} improves this to $f_{\gamma+\delta}^{k+1}w=0$; at the last step one keeps
$2k+1$. The $D_4$ chain has four steps ending at $f_\theta$; for $E_6$, $E_7$, $E_8$ the chains have $9$, $15$ and $27$ steps reaching $f_{\theta-\alpha}$ for the simple root $\alpha$ with $(\theta,\alpha^\vee)=1$, followed by one last Heisenberg step $[f_{\theta-\alpha},f_\alpha]=\pm f_\theta$; they were found by table scans
(\texttt{generators/*\_kernel\_scan.py}) and each step is a kernel-checked table fact. \emph{Non-simply laced.} For $G_2$ the long
roots form an $A_2$ but the short simple root breaks the Heisenberg pattern; \lean{nil_of_ad} replaces it: if $Aw=0$,
$(\operatorname{ad}A)^{r+1}X=0$ and $X^pw=0$, then $((\operatorname{ad}A)^rX)^pw=0$ (expand $\operatorname{ad}_A^{rp}(X^p)$ by Leibniz;
only one term survives). For $G_2$: $(\operatorname{ad}f_{\alpha_1})^3f_{\alpha_2}=6x_3$ and $(\operatorname{ad}f_{\alpha_1})^4f_{\alpha_2}=0$,
so $x_3^{k+1}w=0$, and the Heisenberg pair $(f_{\alpha_2},x_3)$ gives $f_\theta^{2k+1}w=0$. $F_4$ is similar with two
\lean{nil_of_ad} steps and one Heisenberg pair. \qed

\begin{physics}
The top $\theta$-excitation of a weight-$k\theta$ state cannot be repeated more than $2k$ times: a finite multiplet. The
chains are explicit paths through the root system from a simple root to $\theta$.
\end{physics}

\begin{status}
Proved, unconditional, for the six types $D_4$, $E_6$, $E_7$, $E_8$, $G_2$, $F_4$ (with $A_2$ from
Theorem~\ref{l26:thm:kernelA2}; caveats (a), (b) above); \texttt{\#print axioms} standard (re-run for $D_4$, $E_6$, $G_2$, $F_4$; recorded for
$E_7$, $E_8$, whose modules we did not import because of the per-process memory limit).
\end{status}

\begin{warning}
\textbf{Dead end: Heisenberg chains along short roots.} The Heisenberg route needs $\alpha=\gamma+\delta$ with $\alpha+\gamma$ and
$\alpha+\delta$ not roots. It fails for $G_2$ at every non-simple root and for 18 roots of $F_4$. Use \lean{nil_of_ad} or
\lean{top_nil} (next section) instead.
\end{warning}

%======================================================================
\section{The finiteness clause for every type}\label{l26:sec:finiteG}
%======================================================================

\begin{definition}[Splitting the total bound; the finiteness table]\label{l26:def:finG}
\lean{FusionTotalFinG S}: degree pieces of modules with \lean{LocalWeylRelG S} are finite-dimensional.
\lean{FusionTotalDimG S}: their total dimension is $\le n^m$. A letter $c$ is \emph{low} if it is not raising, Cartan or
\lean{cartanN}; it is \emph{known} for $(\rho,w)$ if $c(0)^Nw=0$ for some $N$. A \emph{finiteness table} says: every low letter has
nonzero weight, and in every module with the local Weyl relations every low letter is known.
\end{definition}
\noindent File \leanpath{VOA/\allowbreak DeligneAudit/\allowbreak DeligneFusionFiniteG.lean}.
\begin{leancode}
def FusionTotalFinG : Prop :=
  ∀ (m : ℕ) (V : Type) [AddCommGroup V] [Module ℂ V] (ρ : CurrentModule S.f V) (w : V),
    LocalWeylRelG S ρ w m → ∀ b, FiniteDimensional ℂ (cycDeg ρ w b)

def FusionTotalDimG : Prop :=
  ∀ (m : ℕ) (V : Type) [AddCommGroup V] [Module ℂ V] (ρ : CurrentModule S.f V) (w : V),
    LocalWeylRelG S ρ w m → ∀ B : ℕ, ∑ b ∈ Finset.range B, Module.finrank ℂ (cycDeg ρ w b) ≤ S.n ^ m

theorem fusionTotalBoundG_iff : FusionTotalBoundG S ↔ FusionTotalFinG S ∧ FusionTotalDimG S :=
\end{leancode}
\noindent File \leanpath{VOA/\allowbreak DeligneAudit/\allowbreak DeligneFusionFiniteGLen.lean}.
\begin{leancode}
def LowG (c : Fin S.n) : Prop := ¬ (S.raising c ∨ S.cartan c ∨ S.cartanN c)

structure FiniteTableG : Prop where
  wt_ne : ∀ c, LowG S c → S.wt c ≠ 0
  nil : ∀ (m : ℕ) (V : Type) [AddCommGroup V] [Module ℂ V] (ρ : CurrentModule S.f V) (w : V),
    LocalWeylRelG S ρ w m → ∀ c, LowG S c → Known ρ w c
\end{leancode}
\noindent File \leanpath{VOA/\allowbreak DeligneAudit/\allowbreak DeligneFusionFiniteGNil.lean}.
\begin{leancode}
def Known (w : V) (c : Fin S.n) : Prop := ∃ N, (ρ.op c 0 ^ N) w = 0
\end{leancode}

\begin{dictionary}
\lean{S} is a section variable; in \lean{Known}, \lean{ρ} is a section variable. \lean{S.wt c} is the weight of the basis letter $c$
(a vector of integers, Definition~\ref{l24:def:sectordata}). \lean{FusionTotalDimG} is the type-generic analogue of
\lean{FusionTotalDim}; with \lean{S.n} $=8$ and the $A_2$ sector it is the same statement. As for $A_2$, the dimension clause
alone would be weak on infinite-dimensional pieces; the finiteness clause is therefore proved separately.
\end{dictionary}

\begin{physics}
(Heuristic.) Finiteness of each energy level of the fused multi-particle space, for each Deligne flavour.
\end{physics}

\begin{status}
\lean{FusionTotalDimG S}: \textbf{named hypothesis} for every sector, the per-type residual of the total bound (for $A_2$
the library proves the bound in its $A_2$-specific form \lean{FusionTotalDim}, Theorem~\ref{l26:thm:cvtotal}, and states no
instance of \lean{FusionTotalDimG A2sector}).
\lean{FusionTotalFinG}: proved per type (next theorem).
\end{status}

\begin{theorem}[Kac's lemma; the top of an ad-string; finiteness for all seven types]\label{l26:thm:finiteG}
(1) If $(\operatorname{ad}A)^KB=0$ and $A^Nu=0$ then $A^{N+K}(Bu)=0$. (2) If $(\operatorname{ad}X)^{q+1}Y=0$, $X^pv=0$ and
$Y^N$ kills $X^bv$ for all $b<p$, then $((\operatorname{ad}X)^qY)^Nv=0$. (3) A finiteness table gives \lean{FusionTotalFinG};
hence \lean{FusionTotalDimG S} implies \lean{FusionTotalBoundG S}. (4) For each of the seven sectors a finiteness table is built from
a kernel-checked certificate; so for each type the total bound reduces to \lean{FusionTotalDimG}.
\end{theorem}
\noindent File \leanpath{VOA/\allowbreak DeligneAudit/\allowbreak DeligneFusionFiniteGAlg.lean}.
\begin{leancode}
def adR (X : R) : AddMonoid.End R := AddMonoidHom.mulLeft X - AddMonoidHom.mulRight X

theorem pow_kill_of_ad (A B : Module.End ℂ V) (K N : ℕ) (u : V) (hK : (adR A)^[K] B = 0)
    (hu : (A ^ N) u = 0) : (A ^ (N + K)) (B u) = 0 := by

theorem top_nil (X Y : Module.End ℂ V) (q p N : ℕ) (v : V) (hq : (adR X)^[q + 1] Y = 0)
    (hp : (X ^ p) v = 0) (hY : ∀ b < p, (Y ^ N) ((X ^ b) v) = 0) :
    (((adR X)^[q] Y) ^ N) v = 0 := by
\end{leancode}
\noindent File \leanpath{VOA/\allowbreak DeligneAudit/\allowbreak DeligneFusionFiniteGLen.lean}.
\begin{leancode}
theorem cycDeg_finiteDimensionalG (T : FiniteTableG S) {w : V} {m : ℕ}
    (hW : LocalWeylRelG S ρ w m) (b : ℕ) : FiniteDimensional ℂ (cycDeg ρ w b) := by
\end{leancode}
\noindent File \leanpath{VOA/\allowbreak DeligneAudit/\allowbreak DeligneFusionFiniteG.lean}.
\begin{leancode}
theorem fusionTotalBoundG_of_dim (T : FiniteTableG S) (h : FusionTotalDimG S) :
    FusionTotalBoundG S :=

theorem finiteTableG_of_cert {fnz : Fin S.n → Fin S.n → Fin S.n → Bool}
    (hfz : ∀ a b c, fnz a b c = true → S.f a b c ≠ 0) (kn : Fin S.n → Bool) (cert : List (Step S.n))
    (hkn : ∀ c, kn c = true → S.raising c ∨ S.lowering c) (hv : certValid S.wt fnz kn cert = true)
    (hcov : ∀ c, (certKnown kn cert c && !wEqB (S.wt c) 0) = true ∨
      S.raising c ∨ S.cartan c ∨ S.cartanN c) :
    FiniteTableG S where
\end{leancode}
\noindent File \leanpath{VOA/\allowbreak DeligneAudit/\allowbreak DeligneFusionFiniteGE8.lean}.
\begin{leancode}
theorem e8_finiteTable : FiniteTableG E8sector :=

theorem e8_fusionTotalBoundG_of_dim (h : FusionTotalDimG E8sector) :
    FusionTotalBoundG E8sector :=
\end{leancode}

\begin{dictionary}
\lean{adR X} is $Z\mapsto XZ-ZX$ on a ring (here $\End V$); \lean{(adR X)^[q]} is its $q$-fold iterate. In
\lean{cycDeg_finiteDimensionalG}, \lean{S}, \lean{V}, \lean{ρ} are section variables. \lean{finiteTableG_of_cert} reads a
certificate: \lean{fnz} marks nonzero structure constants (\lean{hfz} relates it to \lean{S.f}), \lean{kn} marks the letters known
from the start (raising letters, nilpotent with $N=1$, and simple lowering letters), \lean{cert} is a list of steps
$(X,Y,\text{path})$, \lean{certValid} is a Boolean check, and \lean{hcov} says every low letter is reached and has nonzero weight.
The per-type theorems \lean{a2_}, \lean{d4_}, \lean{e6_}, \lean{e7_}, \lean{e8_}, \lean{g2_}, \lean{f4_fusionTotalBoundG_of_dim}
(namespaces \lean{FiniteG.A2}, \dots) all have the form printed for $E_8$.
\end{dictionary}

\textbf{Proof on paper.} (1) Expand $A^{N+K}B=\sum_j\binom{N+K}j(\operatorname{ad}A)^j(B)A^{N+K-j}$; terms with $j\ge K$ vanish,
terms with $j<K$ have $A^{\ge N+1}$ on the right. (2) $(\operatorname{ad}X)^{qN}(Y^N)=c\,((\operatorname{ad}X)^qY)^N$ with $c\neq0$ (all
other terms of the Leibniz expansion contain $(\operatorname{ad}X)^{>q}Y=0$); expanding $(\operatorname{ad}X)^{qN}(Y^N)v$ as
$\sum\pm\binom{qN}jX^{qN-j}Y^NX^jv$ kills every term by the hypotheses. This is the $t^{qN}$ coefficient of
$e^{tX}Y^Ne^{-tX}$, done without exponentials. (3) As for $A_2$ (Theorem~\ref{l26:thm:finite}): letters of nonzero weight are
ad-nilpotent of a uniform order, so Kac's lemma on words makes each known letter locally nilpotent on the cyclic module; sorted words
of bounded length span each degree. (4) A certificate step $(X,Y,\text{path})$ uses (2): $X$ and $Y$ known, the $\operatorname{ad}X$-string
from $Y$ runs through basis letters with nonzero constants and stops; its top is then known. Step counts: $A_2$ 1, $G_2$ 4, $D_4$ 8,
$F_4$ 20, $E_6$ 30, $E_7$ 56, $E_8$ 112. For $G_2$ the root $2\alpha_1+\alpha_2$ is reached as the top of an ad-string starting at a
\emph{raising} letter. \qed

\begin{physics}
The zero-mode algebra acts integrably on the fused space: each lowering zero mode can be applied only finitely often to any state.
\end{physics}

\begin{status}
Proved; \texttt{\#print axioms} standard (\lean{top_nil}, \lean{fusionTotalBoundG_of_dim}, \lean{d4_fusionTotalBoundG_of_dim} re-run;
the others recorded). \textbf{After this theorem the per-type residual of U5's total bound is exactly \lean{FusionTotalDimG}.}
\end{status}

%======================================================================
\section{Lattice realizations of Demazure modules}\label{l26:sec:lattice}
%======================================================================

For $A_2$ the Demazure lower bound (D) came from the evaluation model (Theorem~\ref{l26:thm:u7}), whose last step used the
$\mathfrak{sl}_3$-specific Chari--Loktev upper bound. For the other types the natural model is the Frenkel--Kac realization of the
level-one module: $D(1,N\theta)$ is the cyclic $\ltwentysixgt$-submodule of the lattice vertex algebra $V_L$ generated by $e^{N\theta}$.
The library proves the \emph{relations half} of this for the simply-laced types and names the \emph{dimension half}.

\begin{definition}[Lattice realization data; the lattice dimension statement]\label{l26:def:latdata}
For an even positive-definite lattice $L$ and a sector $S$, lattice realization data are weight-one states $X_i\in V_L$ whose
$0$-products are the sector table, each a multiple of a root vector $e^{\mathrm{rt}(i)}$ or a coroot state $\mathrm{cv}(i)(-1)\mathbf1$,
and a root $\theta\in L$ satisfying sign conditions on $\langle\cdot,\theta\rangle$. \textbf{\lean{LatDim Dem N}}: the degree pieces of the
cyclic module of $e^{N\theta}$ dominate $\mathrm{Dem}_N$.
\end{definition}
\noindent File \leanpath{VOA/\allowbreak DeligneAudit/\allowbreak LatticeRealizationGeneric.lean}.
\begin{leancode}
abbrev WL : VertexOperatorAlgebra (LatModuleGen L) := latticeVOAGen L hE hP

def latRepG (S : DeligneSectorData) (X : Fin S.n → LatModuleGen L)
    (hX : ∀ i, X i ∈ (WL L hE hP).grading 1)
    (hbr : ∀ c d, (WL L hE hP).smul (X c) 0 (X d) = ∑ e, S.f c d e • X e) :
    CurrentModule S.f (LatModuleGen L) where
  op c n := mode ((WL L hE hP).Y (X c)) (n : ℤ)

structure LatData (S : DeligneSectorData) (L : IntegralLattice) (hE : L.Even) (hP : L.PosDef) where
  X : Fin S.n → LatModuleGen L
  mem : ∀ i, X i ∈ (WL L hE hP).grading 1
  br : ∀ c d, (WL L hE hP).smul (X c) 0 (X d) = ∑ e, S.f c d e • X e
  cart : Fin S.n → Bool
  rt : Fin S.n → L.carrier
  sg : Fin S.n → ℂ
  cv : Fin S.n → L.carrier
  root : ∀ i, cart i = false → X i = sg i • latChargeGen L (rt i)
  coroot : ∀ i, cart i = true → X i = hSt L (cv i)
  θ : L.carrier
  raising : ∀ e, S.raising e → cart e = false ∧ 0 ≤ L.form (rt e) θ
  cartan : ∀ h, S.cartan h → cart h = true ∧ L.form (cv h) θ = S.thetaH h
  cartanN : ∀ h, S.cartanN h → cart h = true ∧ L.form (cv h) θ = -(S.thetaN h : ℤ)
  lowering : ∀ i, S.lowering i → cart i = false ∧ L.form (rt i) (rt i) = 2 ∧
    -(S.lowExp i : ℤ) ≤ L.form (rt i) θ
  fusion : ∀ x, S.fusion x → cart x = false ∧ (if x = S.c₀ then (-2 : ℤ) else -1) ≤ L.form (rt x) θ

abbrev ρ : CurrentModule S.f (LatModuleGen L) := latRepG S D.X D.mem D.br

abbrev vN (N : ℕ) : LatModuleGen L := latChargeGen L ((N : ℤ) • D.θ)

def LatDim (Dem : ℕ → PowerSeries ℚ) (N : ℕ) : Prop :=
  ∀ M, PowerSeries.coeff M (Dem N) ≤ (Module.finrank ℂ (cycDeg D.ρ (D.vN N) M) : ℚ)
\end{leancode}

\begin{dictionary}
\lean{L}, \lean{hE : L.Even}, \lean{hP : L.PosDef} are section variables; inside \lean{namespace LatData}, the section
variables are \lean{{S : DeligneSectorData} (D : LatData S L hE hP)}, so \lean{ρ}, \lean{vN}, \lean{LatDim} (and \lean{sat},
\lean{demazureRealizationG} below) take \lean{D} as an \emph{explicit} first argument and \lean{S}, \lean{L}, \lean{hE}, \lean{hP}
implicitly; dot notation \lean{D.ρ}, \lean{D.vN N} supplies \lean{D}. \lean{latRepG} is printed without its \lean{comm} field
(the current-algebra commutator, proved from the commutator formula for weight-one fields as in
Theorem~\ref{l26:thm:latrep}); the elision is ours. \lean{latChargeGen L β} is $e^\beta$, \lean{hSt L v} is $v(-1)\mathbf 1$,
\lean{L.form} the bilinear form. \lean{cart i} decides whether $X_i$ is a coroot (\lean{true}) or a root vector (\lean{false});
\lean{sg i} is the complex scalar with $X_i=\mathrm{sg}(i)\,e^{\mathrm{rt}(i)}$ for root letters (field \lean{root}), i.e.\ the
cocycle sign and normalization that make the $0$-products equal the sector table (field \lean{br}); it is unconstrained otherwise,
and is ignored for coroot letters. The sign conditions: raising
$\langle\mathrm{rt}(e),\theta\rangle\ge0$; Cartan $\langle\mathrm{cv}(h),\theta\rangle=\theta(h)$; lowering
$\langle\mathrm{rt}(i),\mathrm{rt}(i)\rangle=2$ and $\langle\mathrm{rt}(i),\theta\rangle\ge-\theta(h_i)$; fusion
$\langle\mathrm{rt}(x),\theta\rangle\ge-1$ ($\ge-2$ for $x=f_\theta$). \lean{vN N} is $e^{N\theta}$.
\end{dictionary}

\begin{physics}
The Frenkel--Kac construction: the level-one currents of a simply-laced $\ltwentysixg$ are realized by $\operatorname{rank}\ltwentysixg$ free
bosons and vertex operators $e^{\alpha}(z)$ for the roots. The state $e^{N\theta}$ is the extremal state of charge $N\theta$; the sign
conditions are statements about which vertex-operator modes annihilate it.
\end{physics}

\begin{status}
Definitions. \lean{LatDim} is the \textbf{named} dimension half of U7 on the lattice side.
\end{status}

\begin{theorem}[$e^{N\theta}$ satisfies Ravinder's $M_{N,0}$ relations; lattice U7]\label{l26:thm:latsat}
(1) For lattice realization data with $\langle\theta,\theta\rangle=2$ and every $N$, $e^{N\theta}$ satisfies \lean{ravFam S N 0}.
(2) Hence \lean{LatDim Dem N} implies \lean{DemazureRealizationG S (ravFam S) Dem N}.
(3) Realization data exist for $D_4$, $E_6$, $E_7$, $E_8$ (\lean{d4LatData}, \lean{e6LatData}, \lean{e7LatData}, \lean{e8LatData},
with $\langle\theta,\theta\rangle=2$ proved as \lean{d4_hθ}, \dots, \lean{e8_hθ}); for these types the relations half of U7 is
unconditional.
(4) (Corollary for $D_4$.) If \lean{FusionTotalBoundG D4sector} holds and \lean{D4LatticeDim N} holds for every $N$, then $D_4$ is
classically free for its explicit presentation, with the expected kernel (\lean{d4_classicallyFree_of_totalBound_latDim}). The
$E_6$, $E_7$ analogues take additional ``star'' inputs (Chapter~\ref{ch:L25}).
There are no lattice data for $G_2$, $F_4$: the Frenkel--Kac construction realizes level-one currents of \emph{simply laced}
algebras only, since it needs all roots to have norm $2$ in an even lattice.
\end{theorem}
\noindent File \leanpath{VOA/\allowbreak DeligneAudit/\allowbreak LatticeRealizationGeneric.lean}.
\begin{leancode}
theorem sat (hθ : L.form D.θ D.θ = 2) (N : ℕ) : Sat S.f D.ρ (D.vN N) (ravFam S N 0) := by

theorem demazureRealizationG (hθ : L.form D.θ D.θ = 2) {Dem : ℕ → PowerSeries ℚ} {N : ℕ}
    (hdim : D.LatDim Dem N) : DemazureRealizationG S (ravFam S) Dem N :=
\end{leancode}
\noindent File \leanpath{VOA/\allowbreak DeligneAudit/\allowbreak LatticeRealizationD4.lean}.
\begin{leancode}
theorem d4_lattice_sat (N : ℕ) : Sat D4sector.f d4LatData.ρ (d4LatData.vN N) (ravFam D4sector N 0) :=

def D4LatticeDim (N : ℕ) : Prop :=
  d4LatData.LatDim DeligneCharacter.D4Wp.hilbSeriesZ N

theorem d4_classicallyFree_of_totalBound_latDim (hF : FusionTotalBoundG D4sector)
    (hdim : ∀ N, D4LatticeDim N) :
    VOA.DeligneHigherRankConditionalCapstone.CaseClassicallyFree .D4 D4Lie.W4 D4Lie.explicitPresentation4 ∧
      RingHom.ker (MvPolynomial.aeval fun i =>
        liDegreeZeroClass D4Lie.W4 1 ⟨D4Lie.explicitPresentation4.gen i,
          D4Lie.explicitPresentation4.gen_weight_one i⟩).toRingHom =
        D4Lie.explicitPresentation4.ideal :=
\end{leancode}
\noindent File \leanpath{VOA/\allowbreak DeligneAudit/\allowbreak LatticeRealizationE8.lean}.
\begin{leancode}
theorem e8_lattice_sat (N : ℕ) : Sat E8sector.f e8LatData.ρ (e8LatData.vN N) (ravFam E8sector N 0) :=

def E8LatticeDim (N : ℕ) : Prop :=
  e8LatData.LatDim (DeligneCarrier.hilbSeries DeligneCharacter.e8DemazureHilb) N

theorem e8_demazureRealization {N : ℕ} (h : E8LatticeDim N) :
    DemazureRealizationG E8sector (ravFam E8sector) (DeligneCarrier.hilbSeries DeligneCharacter.e8DemazureHilb) N :=
\end{leancode}

\begin{dictionary}
\lean{d4LatData} (\leanpath{LatticeRealizationD4.lean}) uses the $D_4$ root lattice \lean{d4} and the basis \lean{D4Lie.xs} of
Chapter~\ref{ch:L23}; \lean{table_eq_sL} (in the \lean{Bridge} files, from the row checks) identifies the sector table with the lattice
model's bracket table. $\theta$ is \lean{wt 15} for $D_4$, \lean{wt 41} for $E_6$, \lean{wt 69} for $E_7$, \lean{wt 127} for $E_8$.
\lean{D4Wp.hilbSeriesZ} and \lean{hilbSeries e8DemazureHilb} are the Demazure series of Chapter~\ref{ch:L25}. There is no $E_8$
classical-freeness corollary of the lattice route in this file; the $E_8$ final theorem (Chapter~\ref{ch:L27}) takes
\lean{DemazureRealizationG} directly. The \lean{_latDim} theorems for $E_6$ and $E_7$ (\lean{e6_}, \lean{e7_classicallyFree_of_totalBound_latDim})
have the shape of the $D_4$ one with extra star inputs.
\end{dictionary}

\textbf{Proof on paper.} Vertex-operator vanishing: $e^\alpha_{(r)}e^\beta=0$ for $r\ge-\langle\alpha,\beta\rangle$. With $\beta=N\theta$ this
kills raising modes ($\langle\mathrm{rt}(e),\theta\rangle\ge0$) and gives the truncations $x(N)$ for fusion letters
($\langle\mathrm{rt}(x),N\theta\rangle\ge-N$) and $f_\theta(2N)$. Coroot modes act by $\langle\mathrm{cv}(h),N\theta\rangle=N\theta(h)$ at grade $0$ and
by $0$ at positive grade. The lowering relations come from the weight shell: after $N\theta(h_i)+1$ applications of $f_i(0)$ the
charge $N\theta-K\alpha_i$ has $\tfrac12|N\theta-K\alpha_i|^2=N^2+K(K-N\langle\alpha_i,\theta\rangle)>N^2$, while the conformal weight is
still $N^2$; the space of such states is zero. \qed

\begin{physics}
In the free-boson picture the null vectors of the Demazure module are just energy conservation: a vertex operator cannot lower
the energy below the minimum allowed in its charge sector.
\end{physics}

\begin{status}
\lean{LatData.sat}, \lean{d4_lattice_sat}, \lean{d4_classicallyFree_of_totalBound_latDim}: \texttt{\#print axioms} standard (re-run);
\lean{e6_}, \lean{e7_}, \lean{e8_lattice_sat}: standard (record, WS-AFF). \lean{D4LatticeDim}, \lean{E6LatticeDim},
\lean{E7LatticeDim}, \lean{E8LatticeDim}: \textbf{named}, not proved. \textbf{Numerics} (exact ranks modulo two primes,
\texttt{aff\_lattice/lattice\_demazure.py}): the lattice cyclic module is graded-sharp, $\dim W_d=[q^d]\ltwentysixHilb D(1,N\theta)$ in every
degree, for $A_2$, $N\le3$ (totals $9,81,729$) and $D_4$, $N\le3$ (totals $29,841,24389$; for $N=3$:
$1925,4396,6475,5655,3590,1617,568,134,28,1$). No lattice data exist for $G_2$, $F_4$ (not simply laced).
\end{status}

\begin{warning}
\textbf{What a proof of the dimension half needs.} For $A_2$ the graded statement followed from a \emph{total} count plus a graded
\emph{upper} bound (Chari--Loktev), which is specific to $\mathfrak{sl}_n$, where $D(1,N\theta)=\ltwentysixWloc(N\theta)$ and
Gelfand--Tsetlin theory applies. For $D_4$, $E_6$, $E_7$, $E_8$, $G_2$, $F_4$ no analogous spanning set is known, and a totals-only
version of (D) does not suffice in the squeeze, which uses $M_{n,0}=\ltwentysixD_n$ coefficientwise. A proof needs the Demazure character
formula (Kumar, Mathieu) or a $D/E$ analogue of Chari--Loktev. It is one of the open inputs of Chapter~\ref{ch:L27} for the types other than $A_2$.
\end{warning}

%======================================================================
\section{How the fusion layer feeds the Deligne chain}\label{l26:sec:feed}
%======================================================================

Table~\ref{l26:tab:feed} lists, per type, what this chapter supplies to the sector bound of Chapter~\ref{ch:L24}
(Theorem~\ref{l24:thm:sfb}, which consumes \lean{FusionModuleBoundG S F}) and to the resummation D3 (which consumes
\lean{MixedFusionRecursion F Dem}), and what it leaves.
\begin{table}[ht]\centering\small
\caption{Per-type output of the fusion layer. ``K'' is the kernel lemma, ``Fin'' the finiteness clause, both proved. ``U7'' is
\lean{DemazureRealizationG} for all $n$. The last column is what Chapter~\ref{ch:L27} must still assume; the alternatives in
parentheses imply the entry before them.}\label{l26:tab:feed}
\begin{tabular}{p{0.07\textwidth}p{0.30\textwidth}p{0.26\textwidth}p{0.26\textwidth}}
\toprule
type & K and Fin (proved) & U7 & still assumed \\
\midrule
$A_2$ & \lean{ravinderKernelRel}, \lean{cycDeg_finiteDimensional} & proved (\lean{demazureRealization_all}) & nothing: \lean{FusionTotalDim} is proved (\lean{CVUniv.fusionTotalDim}) \\
$D_4$ & \lean{d4_ravinderKernelRel'}, \lean{d4_fusionTotalBoundG_of_dim} & proved (\lean{D4Char.d4_demazureRealization}) & \lean{FusionTotalDimG} only \\
$E_6$ & \lean{e6_ravinderKernelRel'}; Fin as for $E_8$ & \lean{DemazureRealizationG}; lattice relations proved & \lean{FusionTotalDimG}, U7 (or \lean{E6LatticeDim}) \\
$E_7$ & \lean{e7_ravinderKernelRel'}; Fin as for $E_8$ & as $E_6$ & as $E_6$ \\
$E_8$ & \lean{e8_ravinderKernelRel'}, \lean{e8_fusionTotalBoundG_of_dim} & as $E_6$ & as $E_6$ \\
$G_2$ & \lean{g2_ravinderKernelRel}; Fin as for $E_8$ & \lean{DemazureRealizationG} (no lattice) & \lean{FusionTotalDimG}, U7 \\
$F_4$ & \lean{f4_ravinderKernelRel}; Fin as for $E_8$ & as $G_2$ & as $G_2$ \\
\bottomrule
\end{tabular}
\end{table}
U6 is \lean{mixedFusionRecursion_fusionOf} for every type. The count ``$\mathrm{fusionOf}(\mathrm{Dem})_m(1)=(\dim\ltwentysixg)^m$'' needed by
the squeeze comes from the Demazure totals $\ltwentysixD_n(1)=(\dim\ltwentysixg+1)^n$ of Chapter~\ref{ch:L25} (Theorem~\ref{l25:thm:total});
the per-type headline wrappers that assemble everything, including the star identities of Chapter~\ref{ch:L25} (proved there; they
are not open inputs of this chapter) for $G_2$, $F_4$, $E_6$,
$E_7$, $E_8$, are the theorems \lean{X_classicallyFree_final} of Chapter~\ref{ch:L27}. An example of such a wrapper at the level of this
chapter is \lean{UnivG.e8_classicallyFree_of_totalBound (hF : FusionTotalBoundG E8sector) (hD : ...) (hcount : E8FusionCount)
(hstar : E8StarIdentity)} in \leanpath{E8FusionKernelFinal.lean}; its count input \lean{E8FusionCount} is removed in Chapter~\ref{ch:L27}
by the unclamped fusion series of Chapter~\ref{ch:L25}.

%======================================================================
\section{Supporting files and certificates, described collectively}\label{l26:sec:certs}
%======================================================================

\paragraph{How to read a certificate.} A reader who has never used Lean needs four facts. (1) The Python generators
that produced the data are \emph{untrusted}: nothing depends on their correctness.
(2) Each family has a \emph{Boolean predicate} or a decidable \lean{Prop}, defined in Lean, together with a \emph{soundness theorem}
proved once (for instance \lean{finiteTableG_of_cert}: a valid certificate yields a finiteness table). (3) The Lean kernel
\emph{evaluates} the predicate on the literal data by \lean{decide +kernel}; that evaluation is a proof, checked by the kernel itself
and not by a compiled program (no \lean{native_decide} appears in these files). (4) Hence a referee must audit only the soundness
theorems, which this chapter presents, and the identification of the data with the mathematical objects (sector tables, which
Chapter~\ref{ch:L24} certifies).

\paragraph{Families.}
\begin{itemize}
\item \textbf{Lattice bridge rows} (62 files \lean{LatticeRealization{D4,E6,E7,E8}Rows*}: 4, 10, 17, 31). Each theorem
  \lean{bridgeRow i} states, for one row $i$ of the bracket table, that the sector table entry equals the lattice model's
  certified entry for every column, by one \lean{decide +kernel}. The \lean{Bridge} files combine the rows into
  \lean{table_eq_sL}; the \lean{Checks} files ($E_7$, $E_8$) hold the index checks (which basis index is $\theta$, which are simple)
  in light files. Reason for the split: a whole-table \lean{decide} over $133^2$ entries took 41\,GB and timed out; one row takes
  about 3\,s and under 3\,GB.
\item \textbf{Finiteness certificates} (\lean{DeligneFusionFiniteG{A2,D4,E6,E7,E8,F4,G2}}, \lean{...E7Chk}, \lean{...E8Chk},
  \lean{DeligneFusionFiniteGCert}). The Boolean checks (\lean{wEqB}, \lean{singleChk}, \lean{topChk}, \lean{pathChk},
  \lean{certKnown}, \lean{certValid}) are stated for a bare weight function in \lean{FiniteGCert}, which imports almost nothing;
  the per-type files read them against the sector by definitional unfolding. Certificate sizes are the step counts of
  Theorem~\ref{l26:thm:finiteG}. Measured: $E_8$ check 12.2\,s, 4.6\,GB.
\item \textbf{Kernel tables} (\lean{D4FusionKernel}, \lean{E6FusionKernel}, \lean{E7FusionKernel}, \lean{E8FusionKernel},
  \lean{E8FusionKernelDefs}, \lean{E8FusionKernelFacts}, \lean{F4LieKernelData}). Each field of \lean{KernelTable}
  (Definition~\ref{l26:def:ravfam}) is proved from an integer table by \lean{decide +kernel}, using literal lists with \lean{getD}
  (selectors such as \lean{hselL}, \lean{bselL}, \lean{kselL} for $E_6$) rather than \lean{![…]} vectors, and, for $E_8$,
  support checks instead of per-entry facts.
\item \textbf{Heisenberg chains} (\lean{*FusionKernelTheta}). Each chain step is a three-entry table fact (\lean{chainZ}).
\end{itemize}

\paragraph{Supporting mathematics files} (all proved, standard axioms, presented in summary).
\begin{itemize}
\item \leanpath{DeligneFusionPBW.lean}: \lean{wordOp} algebra and PBW sorting of words in a current module
  (\lean{cycDeg_le_sortedLow}: degree pieces are spanned by sorted words with ``low'' letters first).
\item \leanpath{DeligneFusionTruncation.lean}: loop ideals; an ideal of current modes that kills $w$ kills the cyclic module
  (\lean{wordOp_eq_zero_of_mem}); truncations propagate upward (\lean{cycDeg_le_truncated}).
\item \leanpath{DeligneFusionGrade.lean} and \leanpath{DeligneFusionUnivGGrade.lean}: degree projections on universal modules,
  quotient current modules, rank--nullity on submodules.
\item \leanpath{DeligneFusionCount.lean}: polynomial versions of $\ltwentysixD_n$ and $\mathrm{fusionF}$, their values at $q=1$
  ($9^n$, $8^n$) and supports (\lean{demazureHilb_sum}, \lean{fusionF_sum}).
\item \leanpath{DeligneFusionECalc.lean}, \leanpath{DeligneFusionCLFilt.lean}, \leanpath{DeligneFusionCLGT.lean},
  \leanpath{DeligneFusionCLNormal.lean}, \leanpath{DeligneFusionCLTheta.lean}, \leanpath{DeligneFusionCLLift.lean}: the
  Chari--Loktev machinery at vector level (the $E$-calculus of CL 3.4, the Gelfand--Tsetlin filtration and \lean{part_i},
  the normal-order lemma and \lean{prop333}, normal words and the statistic $\sigma$, the $\theta$-box \lean{theta_box}, and the
  degree--weight lift). One \lean{simp} over $512$ cases in \lean{table_gw} costs about 15\,s and 2.6\,GB.
\item \leanpath{DeligneFusionCVsl2.lean}, \leanpath{DeligneFusionCVPart.lean}: Chari--Venkatesh Lemma~6.6 made precise
  (\lean{shift_aux}, \lean{shift_iter}), the partitions $\xi^\pm$ and the product identity (\lean{arith_plus}, \lean{prod_split}).
\item \leanpath{G2Fusion.lean}: \lean{g2FusionF}, its U6, and a $G_2$ count statement \lean{G2DemazureCount}.
\item \leanpath{E6FusionKernelFinal.lean}, \leanpath{E7FusionKernelFinal.lean}, \leanpath{E8FusionKernelFinal.lean}: per-type
  wrappers from \lean{FusionTotalBoundG}, the realizations, a count and a star input.
\end{itemize}

%======================================================================
\section{Verification record}\label{l26:sec:verification}
%======================================================================

\paragraph{What the project records.} At an earlier commit the umbrella target \lean{VOA.Deligne} was built green
(\texttt{campaign\_build.sh VOA.Deligne}, 11099 jobs), and every \texttt{\#print axioms} in these files reported a subset of the
three standard axioms (the builds of these files are logged in \texttt{WS-MLD.md},
\texttt{WS-JTP.md}, \texttt{WS-REL.md}, \texttt{WS-CAR.md}, \texttt{WS-AFF.md}). On 2026-10-28 the aggregate build \texttt{bash scripts/campaign\_build.sh VOA} (the full \texttt{lake build VOA} under the project's memory-capped wrapper) was green on 2026-10-28 (a pre-release build): 18{,}901 jobs, and again 18{,}907 jobs after merging the then-current \texttt{origin/main}; the aggregate target contains the umbrella. Since that umbrella build the library gained the Chari--Venkatesh files of
Section~\ref{l26:sec:cv} and the unconditional $A_2$ file \leanpath{VOA/DeligneAudit/A2Unconditional.lean}, and some certificate
files were split for memory; every statement quoted in this chapter occurs verbatim in the state this chapter was written against. For the new files the
record is a set of targeted builds (27 October 2026): \lean{VOA.DeligneAudit.A2Unconditional} green, 9755 jobs,
maximum resident memory 6.13\,GB, with \lean{CVUniv.fusionTotalDim} and \lean{A2.classicallyFree} printing exactly
\lean{[propext, Classical.choice, Quot.sound]}; \lean{CVUniv.sharpRectangularBases} (8863 jobs) and
\lean{CVUniv.residualTheta3Bound} (8887 jobs) green with standard axioms. So these results are
\texttt{\#print axioms}-clean per the targeted build record; the aggregate build was green on 2026-10-28, and
\texttt{\#print axioms} on \lean{A2.classicallyFree} reports \lean{[propext, Classical.choice, Quot.sound]} (Chapter~\ref{ch:L27}). A grep of these
files finds no \lean{sorry}, \lean{admit}, \lean{native_decide}, \lean{implemented_by} and no \lean{axiom} declaration.

\paragraph{What we re-ran for this chapter.} From the repository root, with \texttt{lake env lean -DmaxSynthPendingDepth=3} on
scratch files importing one module each (peak resident memory 5.6--5.9\,GB; each run about 4\,s, using the existing build products):
\begin{plaincode}
import VOA.DeligneAudit.DeligneFusionU7:
  A2.classicallyFree_of_kernelLower', A2.classicallyFree_of_totalDim',
  EvalModel.demazureRealization_all, CLAssembly.clSpanning,
  CLAssembly.finrank_cycDeg_le, A2.fusionTotalDim_of_kernelLower,
  SectorWeyl.ravinderKernelRel, SectorWeyl.fusionTotalBound_of_dim,
  A2.squeeze, FusionLattice.latRep_mabRel, GrFilt.demazureRealization_of_filt
import VOA.DeligneAudit.DeligneFusionIdeal:
  A2.fusionTotalDim_of_idealLower, A2.trunc_all
import VOA.DeligneAudit.DeligneFusionCVTrunc:
  CVsl2.cv_bound, CVsl2.tw_bound, Garland.garland, Sl2Weyl.box_span
import VOA.DeligneAudit.D4FusionKernelTheta:
  UnivG.d4_thetaNil, UnivG.d4_ravinderKernelRel', UnivG.ravinderKernelRelG_of_table,
  UnivG.fusionModuleBoundG_of_kernel, UnivG.dem_le_of_realizationG,
  DeligneCarrier.squeezeG, DeligneCarrier.mixedFusionRecursion_fusionOf,
  UnivG.d4_classicallyFree_of_totalBound
import VOA.DeligneAudit.LatticeRealizationD4:
  LatticeRealization.LatData.sat, LatticeRealization.D4.d4_lattice_sat,
  LatticeRealization.D4.d4_classicallyFree_of_totalBound_latDim
import VOA.DeligneAudit.DeligneFusionFiniteGD4:
  FiniteG.D4.d4_fusionTotalBoundG_of_dim, FiniteG.fusionTotalBoundG_of_dim,
  FiniteG.top_nil
import VOA.DeligneAudit.G2FusionKernelTheta:
  UnivG.g2_thetaNil, UnivG.nil_of_ad, UnivG.g2_ravinderKernelRel
import VOA.DeligneAudit.E6FusionKernelTheta:
  UnivG.E6K.e6_thetaNil, UnivG.E6K.e6_ravinderKernelRel'
import VOA.DeligneAudit.F4LieKernel:
  UnivG.F4K.f4_thetaNil, UnivG.F4K.f4_ravinderKernelRel
\end{plaincode}
Every one of these (all names prefixed by \lean{VOA.DeligneAudit.}) printed
\begin{plaincode}
depends on axioms: [propext, Classical.choice, Quot.sound]
\end{plaincode}
These runs were made at two earlier commits (the second for
\lean{A2.classicallyFree_of_totalDim'}, \lean{EvalModel.demazureRealization_all} and \lean{SectorWeyl.ravinderKernelRel}, with the
same output, peak resident memory 5.8\,GB); the statements they concern are unchanged in the state this chapter was written against. They were \emph{not}
repeated there, nor for this edition, and the axiom checks of the Chari--Venkatesh declarations of Section~\ref{l26:sec:cv} were not run
by us: this checkout has no compiled objects for those modules. We did not import the $E_7$ or $E_8$ kernel and lattice modules, because of the project's per-process memory limit;
their axiom lists are from the project record.
A clean \texttt{\#print axioms} does not report explicit hypotheses: the classical-freeness theorems above remain
\emph{conditional} on the \lean{Prop}s in their signatures, as stated; \lean{CVUniv.fusionTotalDim} has none. We did not run \texttt{lake build}.

%======================================================================
\section{Findings, oddities and Lean defects}\label{l26:sec:defects}
%======================================================================

\begin{warning}
\textbf{Stale docstrings} (no effect on any statement; severity S4).
\begin{enumerate}
\item \lean{A2.DemazureRealization} (\leanpath{DeligneFusionDemazure.lean}): the docstring says that neither the presentation nor the
  character of $D(1,n\theta)$ ``is formalized'', implying the \lean{Prop} is an open input. It is proved for all $n$ by
  \lean{EvalModel.demazureRealization_all}.
\item \leanpath{DeligneFusionAssembly.lean}, module docstring: lists \lean{RavinderKernelTrunc} and \lean{FusionTotalBound} as named
  inputs of \lean{rec_le} and \lean{total_le}. \lean{RavinderKernelTrunc} is a theorem (\lean{ravinderKernelTrunc}), and the finiteness
  half of \lean{FusionTotalBound} is a theorem (\lean{cycDeg_finiteDimensional}). The docstring of \lean{RavinderKernelTrunc} itself
  still calls it a ``named literature input''.
\item \lean{UnivG.ThetaNilG} (\leanpath{D4FusionKernelGeneric.lean}): docstring and module header call it a ``named input''; it is
  proved for every type (Theorem~\ref{l26:thm:thetanil}).
\item \lean{DeligneCarrier.D4DemazureCount} (\leanpath{DeligneFusionSqueezeG.lean}): docstring says ``checked numerically for
  $n\le7$''; the corresponding total is proved (\lean{d4_demazure_total}, \lean{d4DemazureCountZ}, Chapter~\ref{ch:L25}). The
  \lean{Prop} \lean{D4DemazureCount} itself (over $\N$-valued \lean{d4DemazureHilb}) is derived only from a positivity hypothesis
  (\lean{d4DemazureCount_of_coeff}); it is not used by the final theorems.
\item \lean{UnivG.E6K.e6_ravinderKernelRel} (\leanpath{E6FusionKernel.lean}, line~187): the docstring reads ``Ravinder's kernel lemma
  for \texttt{E₇}''. The statement is correct (it is about \lean{E6sector}); only the docstring names the wrong type.
\item \leanpath{DeligneFusionCVTotal.lean}: \lean{SharpRectangularBases} and \lean{ResidualTheta3Bound} are documented as
  ``UNPROVED'' explicit inputs, and the module docstring says the file does not assert unconditional $A_2$ classical freeness. Both
  are proved (Theorem~\ref{l26:thm:cvtotal}); the file is correct as an implication.
\item \leanpath{DeligneFusionU7.lean}, module docstring: ``$A_2$ classical freeness rests on the single named residual
  \lean{RavinderKernelLower}''. It no longer rests on any named statement (Chapter~\ref{ch:L27}).
\end{enumerate}
\end{warning}



\begin{warning}
\textbf{Unused open \lean{Prop}s in the library.} \lean{CL.CLIndependence} (unproved, unused) and \lean{CL.CLCount} (proved, Theorem~\ref{l26:thm:cl}) and the
older interfaces \lean{demazureRealization_of_CL}, \lean{Mchar_eq_of_CL} remain in \leanpath{DeligneFusionCL.lean}. They are not on the path
of any headline theorem.
\end{warning}

%======================================================================
\section{Referee's checklist}\label{l26:sec:checklist}
%======================================================================

\begingroup\small\raggedright
\begin{longtable}{p{0.27\textwidth}p{0.27\textwidth}p{0.24\textwidth}p{0.12\textwidth}}
\toprule
declaration & intended meaning & conditional on / status & check \\
\midrule
\endhead
\multicolumn{4}{l}{\textbf{Headline results}}\\
\lean{CVUniv.}\allowbreak\lean{fusionTotalDim} & $\dim F_m\le8^m$ for $\mathfrak{sl}_3$ & \textbf{none}; proved (record) & Thm.~\ref{l26:thm:cvtotal}\\
\lean{CVUniv.}\allowbreak\lean{sharpRectangularBases}, \lean{residualTheta3Bound} & the two CV base families & none; proved (record) & Thm.~\ref{l26:thm:cvtotal}\\
\lean{CVUniv.}\allowbreak\lean{fusionTotalDim_}\allowbreak\lean{of_}\allowbreak\lean{cvBases} & base families $\Rightarrow$ total bound & the two base statements & Thm.~\ref{l26:thm:cvtotal}\\
\lean{classicallyFree_}\allowbreak\lean{of_}\allowbreak\lean{totalDim'}, \lean{classicallyFree_}\allowbreak\lean{of_}\allowbreak\lean{kernelLower'} & $A_2$ classical freeness & \textbf{\lean{FusionTotalDim}} (proved, so unconditional in Ch.~\ref{ch:L27}), resp.\ \textbf{\lean{RavinderKernelLower}} & Thm.~\ref{l26:thm:u7final}\\
\lean{d4_}\allowbreak\lean{classicallyFree_}\allowbreak\lean{of_}\allowbreak\lean{totalBound_}\allowbreak\lean{latDim} & $D_4$ classical freeness & \textbf{\lean{FusionTotalBoundG}}, \textbf{\lean{D4LatticeDim}} & Thm.~\ref{l26:thm:latsat}\\
\midrule
\multicolumn{4}{l}{\textbf{The $A_2$ fusion layer}}\\
\lean{SectorWeyl.}\allowbreak\lean{MabRel} & relations of Ravinder's $M_{a,b}$, $\mathfrak{sl}_3$ & definition & Def.~\ref{l26:def:mab}\\
\lean{FusionTotalDim} & $\dim F_m\le8^m$ for $\mathfrak{sl}_3$ & definition; proved by \lean{CVUniv.fusionTotalDim} & Def.~\ref{l26:def:mab}\\
\lean{cycDeg_}\allowbreak\lean{finiteDimensional} & degree pieces of fusion modules finite & none; proved & Thm.~\ref{l26:thm:finite}\\
\lean{fusionOf}, \lean{mixedFusionRecursion_}\allowbreak\lean{fusionOf} & U6 by construction & none; proved & Def.~\ref{l26:def:fusionOf}\\
\lean{classicallyFree_}\allowbreak\lean{of_}\allowbreak\lean{fusionBound} & $A_2$ case from U5 & U5 for \lean{fusionF} & Thm.~\ref{l26:thm:a2u5}\\
\lean{squeeze}, \lean{fusionModuleBound_}\allowbreak\lean{of_}\allowbreak\lean{squeeze} & Ravinder's Cor.~1 & squeeze inputs; proved & Thm.~\ref{l26:thm:squeeze}\\
\lean{DemazureRealization} & U7 at $n$ ($A_2$) & proved for all $n$ & Def.~\ref{l26:def:squeezeinputs}\\
\lean{univRep}, \lean{finrank_}\allowbreak\lean{cycDeg_}\allowbreak\lean{le_}\allowbreak\lean{univ} & universal $M_{a,b}$ & none; proved & Def.~\ref{l26:def:umod}\\
\lean{finrank_step} & spanning half of Ravinder's sequence & (KL), which is proved & Thm.~\ref{l26:thm:step}\\
\lean{classicallyFree_}\allowbreak\lean{of_}\allowbreak\lean{totalDim} & $A_2$ from total bound $+$ U7 & \lean{FusionTotalDim}, U7 & Thm.~\ref{l26:thm:a2assembly}\\
\lean{Garland.garland} & Garland identity, $P(u)=1$ & five \lean{include}d hyps & Thm.~\ref{l26:thm:garland}\\
\lean{ravinderKernelRel} & kernel lemma, $\mathfrak{sl}_3$ & none; proved & Thm.~\ref{l26:thm:kernelA2}\\
\lean{Sl2Weyl.box_span} & Chari--Pressley Prop.~6.1 & none; proved & Thm.~\ref{l26:thm:box}\\
\lean{CLAssembly.}\allowbreak\lean{clSpanning}, \lean{finrank_}\allowbreak\lean{cycDeg_}\allowbreak\lean{le} & Chari--Loktev spanning, $\mathfrak{sl}_3$ & none; proved & Thm.~\ref{l26:thm:cl}\\
\lean{CL.}\allowbreak\lean{CLIndependence} & Chari--Loktev independence & \textbf{named}, unused & Def.~\ref{l26:def:cl}\\
\lean{CVsl2.cv_bound}, \lean{tw_bound} & Chari--Venkatesh $\mathfrak{sl}_2$ bound & none; proved; not on headline path & Thm.~\ref{l26:thm:cv}\\
\lean{EvalModel.}\allowbreak\lean{demazureRealization_}\allowbreak\lean{all} & U7 for $A_2$ & none; proved & Thm.~\ref{l26:thm:u7}\\
\lean{FusionLattice.}\allowbreak\lean{latRep_}\allowbreak\lean{mabRel} & $e^{N\theta}\in V_{A_2}$ has $M_{N,0}$ relations & none; proved & Thm.~\ref{l26:thm:latrep}\\
\lean{RavinderKernelLower}, \lean{IdealLower} & lower bounds implying the total bound & \textbf{named} & Def.~\ref{l26:def:s2}\\
\lean{fusionTotalDim_}\allowbreak\lean{of_}\allowbreak\lean{kernelLower}, \lean{_of_idealLower} & reductions & the named \lean{Prop}; proved & Thm.~\ref{l26:thm:a2red}\\
\midrule
\multicolumn{4}{l}{\textbf{The type-generic layer}}\\
\lean{UnivG.univRepG}, \lean{finrank_}\allowbreak\lean{le_}\allowbreak\lean{McharG} & universal module for any specification & \lean{UnivFD}; proved & Def.~\ref{l26:def:relspec}\\
\lean{FusionTotalBoundG}, \lean{DemazureRealizationG} & total bound, U7, any type & \textbf{named} (U7 proved for $D_4$) & Def.~\ref{l26:def:lwspec}\\
\lean{squeezeG}, \lean{fusionOf_count} & generic squeeze and count & squeeze inputs; proved & Thm.~\ref{l26:thm:squeezeG}\\
\lean{finrank_stepG}, \lean{fusionModuleBoundG_}\allowbreak\lean{of_}\allowbreak\lean{kernel} & generic step; U5 from kernel lemma & printed hyps; proved & Thm.~\ref{l26:thm:kernelstep}\\
\lean{KernelTable}, \lean{ThetaNilG} & table facts; integrability & definitions & Def.~\ref{l26:def:ravfam}\\
\lean{ravinderKernelRelG_}\allowbreak\lean{of_}\allowbreak\lean{table} & generic kernel lemma & table, \lean{ThetaNilG}; proved & Thm.~\ref{l26:thm:kernelG}\\
\lean{d4_thetaNil}, \dots, \lean{f4_}\allowbreak\lean{ravinderKernelRel} & kernel lemma, every type & none; proved & Thm.~\ref{l26:thm:thetanil}\\
\lean{FusionTotalDimG} & $\dim\le(\dim\ltwentysixg)^m$ & \textbf{named}, every type other than $A_2$ & Def.~\ref{l26:def:finG}\\
\lean{top_nil}, \lean{fusionTotalBoundG_}\allowbreak\lean{of_}\allowbreak\lean{dim}, \lean{e8_}\allowbreak\lean{fusionTotalBoundG_}\allowbreak\lean{of_}\allowbreak\lean{dim} & finiteness, every type & certificate; proved & Thm.~\ref{l26:thm:finiteG}\\
\midrule
\multicolumn{4}{l}{\textbf{Lattice realizations}}\\
\lean{LatData}, \lean{LatDim} & lattice data; dimension half & definitions; \lean{LatDim} \textbf{named} & Def.~\ref{l26:def:latdata}\\
\lean{LatData.sat}, \lean{d4_lattice_sat}, \lean{e8_lattice_sat} & relations half of lattice U7 & none; proved & Thm.~\ref{l26:thm:latsat}\\
\bottomrule
\end{longtable}
\endgroup

%======================================================================
\section{Notes and references}\label{l26:sec:notes}
%======================================================================

\textbf{Current algebras, Weyl and Demazure modules.} V.~Chari and A.~Pressley, \emph{Weyl modules for classical and quantum affine
algebras}, Represent.\ Theory 5 (2001) 191--223, arXiv:math/0004174 (Proposition~6.1, Lemmas~6.5--6.7). H.~Garland, \emph{The arithmetic
theory of loop algebras}, J.\ Algebra 53 (1978) 480--551. V.~Chari and S.~Loktev, \emph{Weyl, Demazure and fusion modules for the current
algebra of $\mathfrak{sl}_{r+1}$}, Adv.\ Math.\ 207 (2006) 928--960, arXiv:math/0502165. G.~Fourier and P.~Littelmann, \emph{Weyl modules,
Demazure modules, KR-modules, crystals, fusion products and limit constructions}, Adv.\ Math.\ 211 (2007) 566--593. K.~Naoi, \emph{Weyl modules,
Demazure modules and finite crystals for non-simply laced type}, Adv.\ Math.\ 229 (2012) 875--934. S.~Kumar, \emph{Demazure character formula in
arbitrary Kac--Moody setting}, Invent.\ Math.\ 89 (1987); O.~Mathieu, \emph{Formules de caract\`eres pour les alg\`ebres de Kac--Moody
g\'en\'erales}, Ast\'erisque 159--160 (1988). Y.~Sanderson, \emph{On the connection between Macdonald polynomials and Demazure characters}, J.\
Algebraic Combin.\ 11 (2000); B.~Ion, \emph{Nonsymmetric Macdonald polynomials and Demazure characters}, Duke Math.\ J.\ 116 (2003).

\textbf{Fusion products.} B.~Feigin and S.~Loktev, \emph{On generalized Kostka polynomials and the quantum Verlinde rule}, AMS Transl.\ 194
(1999), arXiv:math/9812093. E.~Feigin, \emph{The PBW filtration}, SIGMA 4 (2008) 070, arXiv:0806.4851 (Prop.~1, Cor.~2 in the SIGMA and
arXiv v2 numbering, Prop.~2.1, Cor.~2.2 in arXiv v1: the step with the gap discussed in Section~\ref{l26:sec:open}). B.~Ravinder, \emph{Demazure modules, Chari--Venkatesh modules and fusion products}, SIGMA 10
(2014) 110, arXiv:1409.0274 (Theorem~1, Proposition~1, Corollaries~1--3, Proposition~5, in the v4 numbering; stated for every simple
$\ltwentysixg$ except Theorem~2 and Corollary~3, which assume simply laced). V.~Chari and R.~Venkatesh, \emph{Demazure modules,
fusion products and $Q$-systems}, Comm.\ Math.\ Phys.\ 333 (2015) 799--830, arXiv:1305.2523 (\S6, Lemma~6.6). D.~Kus and P.~Littelmann, \emph{Fusion products and toroidal algebras}, Pacific J.\ Math.\ 278 (2015) 427--445, arXiv:1411.5272 (Theorem~3.5,
Corollary~3.6(3), Definition~4.1 and Theorem~4.2, the remark on p.~438 on Feigin's gap, and the step on p.~437 flagged by the project). D.~Setia, S.~Rani and T.~Khandai, \emph{Filtration of tensor product of local
Weyl modules for $\mathfrak{sl}(n+1)[t]$}, arXiv:2405.11944 (\S8, Theorems~8.1.2 and 8.4.1; we
cite the 2024 arXiv version, which the project inspected; the journal version of 2025 was not checked). S.~Rani et al., \emph{Graded character formula for fusion products in type $A_2$}, arXiv:2303.03337 (two factors, gap-free). M.~Brito,
V.~Chari and A.~Moura, \emph{Demazure modules of level two and prime representations}, arXiv:1504.00178 (does not cover $V(\theta)^{*N}$).

\textbf{Physics.} C.~Beem, M.~Lemos, P.~Liendo, W.~Peelaers, L.~Rastelli and B.~C.~van Rees, \emph{Infinite chiral symmetry in four
dimensions}, Commun.\ Math.\ Phys.\ 336 (2015) 1359--1433, arXiv:1312.5344 (the chiral algebra of a four-dimensional $\mathcal N=2$
theory; its vacuum character is the Schur index). C.~Beem and L.~Rastelli, \emph{Vertex operator algebras, Higgs branches and modular
differential equations}, JHEP 08 (2018) 114, arXiv:1707.07679 (the comparison of the Schur index with arc-space Hilbert series).

\textbf{Lattice realizations.} I.~Frenkel and V.~Kac, \emph{Basic representations of affine Lie algebras and dual resonance models}, Invent.\
Math.\ 62 (1980) 23--66. V.~Kac, \emph{Infinite-dimensional Lie algebras}, 3rd ed., Cambridge (1990), Lemma~3.4 and Ch.~12; J.~Humphreys,
\emph{Introduction to Lie algebras and representation theory}, \S21.

\textbf{Project records.} The campaign logs \texttt{WS-MLD.md}
(the $A_2$ fusion layer, U7, the reductions and the slice lemma), \texttt{WS-JTP.md} (the generic squeeze, universal modules,
kernel lemmas), \texttt{WS-REL.md} (finiteness certificates), \texttt{WS-CAR.md} ($E_6$--$E_8$ kernel chains),
\texttt{WS-AFF.md} (lattice realizations; Feigin's route and the lattice form of the slice problem),
\texttt{A2\_SOURCE\_AUDIT\_2026\_10\_27.md} (the source audit of the $A_2$ literature).

\endgroup
